% !TEX program = xelatex
% !TEX encoding = UTF-8 Unicode
\documentclass[11pt,a4paper]{article}

% ---- preamble ----
\usepackage[margin=1in]{geometry}
\usepackage{setspace}
\doublespacing
\usepackage{fontspec}
\setmainfont{TeX Gyre Termes}
\usepackage{hyperref}
\hypersetup{colorlinks=true,linkcolor=blue,citecolor=blue,urlcolor=blue}
\usepackage{booktabs}
\usepackage{longtable}
\usepackage{enumitem}
\usepackage{titlesec}
\usepackage{tocloft}
% IPA via direct Unicode input (XeLaTeX native)
\usepackage{morefloats} % extra float capacity

% ---- partial build support ----
\ifdefined\includeonlysections
  \includeonly{\includeonlysections}
\fi

% ---- metadata ----
\title{Orthography to Phoneme Conversion in Southern New England Algonquian:\\A Multi-Source Research Document}
\author{Michael Conrad\\Brothertown-Language Project}
\date{\today}

\begin{document}

\maketitle
\tableofcontents
\newpage

% ============================================================
% SECTION 1 — Conversion Overview
% ============================================================
\section{Conversion Overview}
\label{sec:overview}

% Section 1 — Conversion Overview
% TLDR for linguist reviewers. Proto-examples for pending specs,
% real examples for completed specs, technique comparison, status table,
% and related work on the colonial recorder perception problem.

\subsection{Conversion Examples and Methodologies}
\label{sec:examples}

This section presents concrete orthography-to-phoneme conversions across the
eight source languages. For sources where implementation is complete, real
conversion results are shown. For sources where only research has been
completed, \textbf{[PROTO]} examples are given --- these represent the current
best hypothesis grounded in the research evidence, with documented uncertainty.

\subsubsection*{Narragansett (Williams 1643) --- \texttt{[PROTO]}}

\begin{quote}
\textbf{Input:} \texttt{w\'etu} (house) \\
\textbf{Output:} /wi:tə w/ \\
\textbf{Method:} Comparative calibration via PA *wi:kiwa:mi, Massachusett
cognate \textit{wetu} (Eliot). The initial \textit{w\'e-} maps to /wi:/ via
Williams' circumflex-as-length convention; final \textit{-tu} maps to /tə w/
via PA reflex. \\
\textbf{Uncertainty:} Vowel length in the first syllable is inferred from the
circumflex accent, which Williams uses for primary stress --- whether this
marks phonemic length or stress prominence is debated.
\end{quote}

\begin{quote}
\textbf{Input:} \texttt{Aaunchem\'okaw} (Record 3834) \\
\textbf{Output:} /ã:čimo:hka:w/ \texttt{[PROTO]} \\
\textbf{Method:} FSM with preaspiration inference. Williams writes \textit{k}
but the only phonetic record (\textbackslash ph field) shows /hk/. The
preaspiration gap --- English-speaking recorders routinely fail to hear /h/
before stops --- is the single most important finding for this source. \\
\textbf{Uncertainty:} The nasalization of the initial vowel is inferred from
the \textit{Aaun-} spelling, which may represent /ã:/ or /aun/.
\end{quote}

\subsubsection*{Wampanoag (Trumbull/Eliot 1903) --- \texttt{[PROTO]}}

\begin{quote}
\textbf{Input:} \texttt{matta} (no, not) \\
\textbf{Output:} /mata/ \\
\textbf{Method:} Direct Eliot orthography mapping. Eliot's system is the most
regularized colonial orthography in the corpus. Confirmed via PA *mati-.
\end{quote}

\begin{quote}
\textbf{Input:} \texttt{achmōwonk} (devil) \\
\textbf{Output:} /čomu:wa\ng k/ \texttt{[PROTO]} \\
\textbf{Method:} 4-stage FSM pipeline: (1) normalize ō (Eliot's oo-ligature) to /u:/, (2) map
\textit{ach-} to /č/ via Eliot's c-before-a rule, (3) resolve the
\textit{-wonk} suffix via PA *-a:kani, (4) comparative fallback against
Pokanoket \textit{hobbamock} (Winslow) to check for dialect divergence. \\
\textbf{Uncertainty:} The ō symbol (Eliot's oo-ligature,
U+A74F) was historically misidentified as an infinity symbol. Its phonetic
value is long /u:/, but the precise quality (tense vs.\ lax) is debated.
\end{quote}

\subsubsection*{Mohegan-Pequot (Fielding 2013) --- \texttt{[PROTO]}}

\begin{quote}
\textbf{Input:} \texttt{wetu} (house) \\
\textbf{Output:} /wetu/ \\
\textbf{Method:} Direct 1:1 phoneme mapping. Fielding's modern practical
orthography is essentially phonemic. Confirmed via PA *wi:kiwa:mi and
cognate with Narragansett \textit{w\'etu} and Massachusett \textit{wetu}.
\end{quote}

\subsubsection*{Wampanoag (Wood 1634) --- \texttt{[PROTO]}}

\begin{quote}
\textbf{Input:} \texttt{abbona} (five) \\
\textbf{Output:} /apana/ \texttt{[PROTO]} \\
\textbf{Method:} Normalize \textit{b} to /p/ (Wood writes \textit{b} where
Eliot/Trumbull use \textit{p} --- direct evidence of voicing neutralization).
Segment the initial \textit{a-} as a demonstrative prefix. \\
\textbf{Uncertainty:} Wood's triple spelling for 'devil'
(\textit{abbomocho}/\textit{abomocho}/\textit{abamacho}) is the strongest
evidence in the corpus for colonial vowel instability. The medial vowel
could be /a/, /ə/, or /o/.
\end{quote}

\subsubsection*{Mahican (Edwards 1787) --- \texttt{[PROTO]}}

\begin{quote}
\textbf{Input:} \texttt{ghusooh} (eight) \\
\textbf{Output:} /xə sə w/ \texttt{[PROTO]} \\
\textbf{Method:} The \textit{gh} grapheme is the most unusual across all eight
sources. Candidates include /x/ (voiceless velar fricative), /h/, or /k/.
PA *nešwa:šika and Wampanoag cognate \textit{shwoshuk} suggest /x/
is most likely --- a Mahican phoneme not present in SNEA languages. \\
\textbf{Uncertainty:} The \textit{gh} value cannot be confirmed without
cross-linguistic cognate evidence from Munsee Delaware or other Mahican-adjacent
languages.
\end{quote}

\subsubsection*{Mohegan-Pequot (Prince-Speck 1904) --- \texttt{[PROTO]}}

\begin{quote}
\textbf{Input:} \texttt{bahkeder} \\
\textbf{Output:} /pahketə r/ \texttt{[PROTO]} \\
\textbf{Method:} German orthography normalizer: \textit{b-} initial = /p/
(German final devoicing rule applied to perception), \textit{-hk-} = /hk/
(preaspiration preserved --- Prince-Speck is one of the few sources that
writes preaspiration explicitly), \textit{-der} = /tə r/ (German
\textit{d} = /t/ in final position). \\
\textbf{Uncertainty:} The \textit{-er} suffix may represent /ə r/ or
a reduced syllable. Loanword detection needed --- \textasciitilde20--30\% of
the 446 entries are English loanwords.
\end{quote}

\subsection{Technique Comparison and Selection Rationale}
\label{sec:techniques}

The technique selection for each source is itself a research output, not a
predetermined input. Each source's unique characteristics --- corpus size,
orthographic regularity, ambiguity profile, available comparative evidence,
recorder linguistic background --- determine which approach is appropriate.
A one-size-fits-all approach would fail because the sources span 263 years
of colonial documentation with radically different orthographic systems.

Table~\ref{tab:techniques} compares the approaches considered across all
sources, with expected accuracy and rationale for selection or rejection.
The per-source rationale is documented in detail in Part~I.

\begin{table}[ht]
\centering
\caption{Technique comparison across SNEA sources}
\label{tab:techniques}
\small
\begin{tabular}{lp{3cm}p{2.5cm}p{4.5cm}}
\toprule
\textbf{Technique} & \textbf{Sources} & \textbf{Expected Accuracy} & \textbf{Rationale} \\
\midrule
Regex-based conversion & None & Low &
Rejected --- ambiguities are fundamentally context-dependent; regex cannot
distinguish /a/ from /ə/ from /a:/ in Williams' orthography without
comparative evidence. \\
\midrule
FSM with weighted transitions & Narragansett, Wampanoag, Mahican & Moderate &
Selected --- handles context-dependent ambiguity via state transitions;
weights encode probability from PA cognate frequency. \\
\midrule
1:1 phoneme mapping & Mohegan-Pequot (Fielding) & High &
Selected --- Fielding's modern practical orthography is already phonemic;
conversion is a direct lookup. \\
\midrule
ML classifier & Narragansett & Unknown &
Selected but blocked --- only 3 records have phonetic (\textbackslash ph)
fields; training data insufficient without hand-annotating \textasciitilde50
entries. \\
\midrule
Comparative calibration & All colonial sources & High for confirmed cognates &
Selected --- PA reconstructions and sister-language evidence (Massachusett,
Mohegan) disambiguate vowel quality and preaspiration. \\
\midrule
4-stage FSM pipeline & Wampanoag (Trumbull) & Moderate--High &
Selected --- preprocessing $\rightarrow$ phonemic conversion $\rightarrow$
comparative fallback $\rightarrow$ statistical verification. \\
\midrule
Loanword detection FSM & Mohegan-Pequot (Prince-Speck) & High &
Selected --- \textasciitilde20--30\% English loanwords must be identified
before phonemic conversion to avoid false cognate matches. \\
\midrule
German orthography normalizer & Mohegan-Pequot (Prince-Speck) & High &
Selected --- Prince's Neogrammarian training imposed German conventions
(\textit{tsh} for /ʃ/, \textit{z} for /ts/, Dehnungs-\textit{h}). \\
\midrule
Morphological segmentation & Mahican (Edwards) & Moderate &
Selected --- sentence-level entries require prefix/suffix parsing
(\textit{k-} 2nd person, \textit{n-} 1st person) before conversion. \\
\bottomrule
\end{tabular}
\end{table}

\subsection{Source Status and Expected KPIs}
\label{sec:status}

Table~\ref{tab:status} shows the current status of each source. Research is
complete for all eight sources. Implementation is pending for all sources.
Proto-KPI targets are research-grounded predictions based on the known
ambiguity profile of each source.

\begin{table}[ht]
\centering
\caption{Source status and proto-KPI targets}
\label{tab:status}
\small
\begin{tabular}{lcccc}
\toprule
\textbf{Source} & \textbf{Research} & \textbf{Implementation} &
\textbf{Expected Recall} & \textbf{Expected Accuracy} \\
\midrule
Narragansett (Williams 1643) & Complete & Pending & 85--90\% & 70--80\% \\
Wampanoag (Trumbull 1903) & Complete & Pending & 90--95\% & 80--90\% \\
Mohegan-Pequot (Fielding 2013) & Complete & Pending & 95--98\% & 90--95\% \\
Wampanoag multi-source & Complete & Pending & 85--90\% & 75--85\% \\
Wampanoag (Winslow 1624) & Complete & Pending & 70--80\% & 60--70\% \\
Wampanoag (Wood 1634) & Complete & Pending & 80--85\% & 65--75\% \\
Mohegan-Pequot (Prince-Speck 1904) & Complete & Pending & 85--90\% & 75--85\% \\
Mahican (Edwards 1787) & Complete & Pending & 80--85\% & 70--80\% \\
\bottomrule
\end{tabular}
\end{table}

\noindent\textbf{Note:} Expected accuracy is lower for sources with extreme
vowel ambiguity (Williams, Wood) or tiny corpora (Winslow). Mohegan-Pequot
(Fielding) is expected to have the highest accuracy because its modern
practical orthography is already phonemic. These targets will be updated as
implementation validates or refutes the research predictions.

\subsection{The Colonial Recorder Perception Problem}
\label{sec:related-work}

Every source in this study shares a fundamental challenge: English-speaking
recorders with different linguistic backgrounds hearing Algonquian phonemes
through their native phonological filter. This is not unique to SNEA --- it is
a well-documented phenomenon across colonial documentation of North American
languages.

\subsubsection*{Dr.\ Keith Cunningham --- Nanticoke}

Dr.\ Keith Cunningham (Linguist, Nanticoke Indian Tribe) completed his
doctoral dissertation \emph{A Phonological Analysis of Nanticoke With
Practical Applications for Language Revitalization} at Georgetown University
(2024). The dissertation analyzes three 18th-century Nanticoke word lists
(Heckewelder 1785) using the same methodology employed here: historical
phonology from colonial recorders, with the same challenges of multilingual
informants and orthographic complexity. His 2024 Algonquian Conference
presentation \emph{A Revised Phonological Analysis of the Heckewelder
Vocabulary of Nanticoke} directly parallels the analyses in this document.
Cunningham has also co-authored a Nanticoke language textbook, bridging
directly into practical revitalization --- the same goal as the SNEA
orthography-to-phoneme conversion pipeline.

Cunningham's work demonstrates that the colonial recorder perception problem
extends beyond SNEA into the broader Eastern Algonquian family (Nanticoke,
Maryland Algonquian). His methodology --- PA comparative calibration,
recorder-bias identification, cross-source verification --- is directly
transferable to the SNEA sources.

\subsubsection*{Dr.\ Craig Kopris --- Wyandot}

Dr.\ Craig Kopris (Independent scholar) has nearly three decades of experience
in language description and endangered language revitalization, with emphasis
on Iroquoian languages. His research note \emph{Les sons wyandots per\c{c}us
par des oreilles \'etrang\`eres} (Wyandot Sounds through Foreign Ears,
\emph{Recherches am\'erindiennes au Qu\'ebec}, 2014) is the exact framing of
the core problem addressed in this document: how European recorders
systematically misperceived and mis-transcribed indigenous phonemes. Kopris
demonstrates that knowing the sound system of the transcriber's first language
is the key to recovering the phonetic details of the target language.

His paper \emph{Wyandot Phonology: Recovering the Sound System of an Extinct
Language} (Proceedings of the Second Annual High Desert Linguistics Society
Conference, 1999) applies the same recovery methodology to Wyandot, a
Northern Iroquoian language unrelated to Algonquian. This demonstrates that
the foreign-ear perception problem is universal across colonial documentation
of North American languages, not limited to any single language family. His
Dictionary of Wyandot project (NEH award FN-255576-17) shows the long-term
application of these methods to language revitalization.

\subsubsection*{Relevance to This Project}

The work of Cunningham and Kopris establishes that the methodological
framework for addressing the colonial recorder perception problem is well
established. The SNEA orthography-to-phoneme project extends this framework
to eight sources across four languages (Narragansett, Wampanoag, Mohegan-Pequot,
Mahican) spanning 263 years of colonial documentation (1624--1904). The
techniques developed here --- FSM with weighted transitions, comparative
calibration via PA, colonial English phonology baselines, preaspiration
detection --- contribute to an ongoing methodological conversation about
recovering phonology from colonial sources across Eastern North America.


% ============================================================
% PART I — Source Workflows
% ============================================================
\part{Source Workflows}

% Narragansett (Williams 1643) — Source Workflow
% Issue #1362
% Status: Research complete, implementation pending

\section{Narragansett: Williams 1643}
\label{sec:1362}

\subsection{Source Description}

\textbf{Recorder:} Roger Williams (1603--1683), founder of Rhode Island Colony.
\textbf{Source:} \emph{A Key into the Language of America} (1643), London: Gregory Dexter.
\textbf{Language:} Narragansett [xnt] (Southern New England Algonquian).
\textbf{Corpus:} 2,134 records.
\textbf{Orthography:} English-based ad-hoc spelling, seventeenth-century Early
Modern English conventions. Williams marks stress with diacritics (á primary,
à low/falling, â rising-falling) --- a major advantage over other colonial
sources such as Eliot's Massachusett Bible (1663), which uses no stress
marking.

\subsubsection{Recorder Biography}

Roger Williams (1603--1683) was a Puritan minister, theologian, and founder
of the Colony of Rhode Island and Providence Plantations. A graduate of
Pembroke College, Cambridge (B.A. 1627), Williams was a skilled linguist who
reportedly mastered Latin, Greek, Hebrew, Dutch, French, and several
Algonquian languages. He arrived in Massachusetts Bay in 1631 and began
intensive work with the Narragansett people after his banishment from the
Massachusetts Bay Colony in 1635. His \emph{A Key into the Language of
America} (1643) was the first extended study of a Native American language
published in English. The work takes the form of a phrase book organized by
topic (greetings, family, hunting, etc.), with each Narragansett phrase
accompanied by an English translation and occasional commentary.

\subsubsection{Orthography System}

Williams used an English-based ad-hoc spelling system with no standardized
orthography. He applied English sound-spelling conventions (seventeenth-century
Early Modern English) to Narragansett, an Algonquian language with a radically
different phoneme inventory. The system was designed to be usable by English
colonists, not by Algonquianists --- it systematically collapses phonemic
distinctions that English lacks.

\subsubsection{Stress Marking}

Williams' orthography is \textbf{stress-marked}, unlike Eliot's Massachusett
--- this is a major advantage. Eliot's Massachusett Bible (1663) uses no
accent marks, making stress reconstruction for Massachusett far more difficult.
Williams' accents provide critical data about Narragansett prosody.

\begin{table}[h]
\centering
\begin{tabular}{ll}
\toprule
Mark & Meaning \\
\midrule
á & Primary stress (or high pitch --- debated) \\
à & Low/falling tone (secondary stress rarely written) \\
â & Rising-falling tone \\
\bottomrule
\end{tabular}
\caption{Stress marking in Williams' orthography}
\label{tab:1362-stress}
\end{table}

Aubin's proposed stress rule (from personal correspondence) posits that all
long vowels are stressed and all alternate even-numbered short vowels are
stressed. However, this rule requires correct vowel length classification,
which has not been achieved for extinct SNEA languages. The problem is
circular: vowel length is needed to assign stress, but stress is one of the
cues for vowel length. Costa (2007) discusses the prosodic structure of SNEA
languages and notes that Narragansett stress patterns appear similar to Loup A
and Massachusett, with a preference for stress on long syllables and
alternating rhythm on short syllables --- but the documentary record is too
thin to confirm this definitively.

\subsection{Vowel Inventory}

\subsubsection{Simple Vowels}

\begin{table}[h]
\centering
\begin{tabular}{llll}
\toprule
Williams & IPA & Narragansett Phoneme & Notes \\
\midrule
a & /ə/, /ɔː/, /aː/ & PA *a *eː *aː & Three-way ambiguity \\
ah & /aː/, /ɔː/ & PA *aː & Long variants \\
e & /iː/, /ɛ/, /ə/, ∅ & PA *iː *e *a & Four-way ambiguity; silent word-final \\
ê & /iː/ & PA *iː & Circumflex = long i \\
i & /ə/, /iː/, /ɪ/ & PA *a *iː *i & Often = schwa in unstressed \\
o & /ə/, /aː/, /uː/; after w: /aː/, /ɔː/ & PA *a *aː *oː & Broad range \\
oo, ô & /uː/ & PA *oː & Consistent long u \\
u & /ə/, /a/; word-initially possibly /w/ & PA *a & Ambiguous schwa-ish \\
ie & /iː/ & PA *iː & \\
ee & /iː/ & PA *iː & Consistent \\
ei & /iː/, /ə/ & PA *iː *a & \\
ih & /ə/ & PA *a & \\
oh & /ə/, /oʊ/ & PA *a & \\
eu & /juː/ & --- & Rare, from English loan phonology \\
ea & /iː/, /ɛə/, /aː/ & PA *iː *aː & \\
\bottomrule
\end{tabular}
\caption{Simple vowel mappings in Williams' orthography}
\label{tab:1362-vowels}
\end{table}

\subsubsection{Nasalized Vowels}

A critical feature of SNEA phonology: Proto-Algonquian *aː and *oː developed
nasalized reflexes in Narragansett before certain consonants. Williams does
not mark nasality directly; he uses nasal coda consonant digraphs to imply it.

\begin{table}[h]
\centering
\begin{tabular}{ll}
\toprule
Williams & IPA \\
\midrule
an, aum, aun & /ã/ (PA *aː → ą in some environments) \\
om, on & /õ/ (PA *oː → ǫ) \\
ook & /oːk/, /oːkʷ/ (may be nasalized in context) \\
\bottomrule
\end{tabular}
\caption{Nasalized vowel patterns}
\label{tab:1362-nasal}
\end{table}

The nasalization is phonemic: \textit{n now} (yes) vs. \textit{núnow} (no)
--- only nasality distinguishes them, and Williams collapses this distinction
orthographically in related forms.

\subsubsection{Diphthong-Like Sequences}

\begin{table}[h]
\centering
\begin{tabular}{lll}
\toprule
Williams & IPA & Notes \\
\midrule
au & /ɔː/, /aʊ/ & \\
aw & /ɔː/, /ɒ/ & \\
ai & /aː/, /eɪ/ & \\
aü & /aː.uː/ & Rare; hiatal /aː/ + /uː/, not a true diphthong \\
\bottomrule
\end{tabular}
\caption{Diphthong-like sequences}
\label{tab:1362-diphthongs}
\end{table}

\subsection{Consonant Inventory}

SNEA languages neutralized the PA voicing distinction in stops. Williams'
English ear interpreted this as a gradient between voiced and voiceless,
writing both \textit{p} and \textit{b} (though \textit{b} is rare),
\textit{t} and \textit{d}, \textit{k} and \textit{g} (nearly absent). In
almost all cases, \textit{p}, \textit{t}, \textit{k} represent an unaspirated
stop --- between English /b/ and /p/, /d/ and /t/, /ɡ/ and /k/.

\subsubsection{Stops}

\begin{table}[h]
\centering
\begin{tabular}{lll}
\toprule
Williams & IPA & Notes \\
\midrule
p, pp & /b/ ~ /p/ [p] & Unaspirated lenis stop; alternates with \textit{b} \\
t & /d/ ~ /t/ ~ /tʲ/ & Before /i/ may palatalize to /tʃ/ \\
tt & /t/ ~ /d/ ~ /tʲ/ & Geminate, also palatalized in palatal contexts \\
k, kk & /k/ ~ /kʷ/ ~ /x/ & Labialized before /o/ /u/; guttural before /a/ \\
c, cc & /k/ ~ /kʷ/ & Before a/o/u = /k/; before e/i = /s/ (English soft-c) \\
\bottomrule
\end{tabular}
\caption{Stop consonants}
\label{tab:1362-stops}
\end{table}

Cowan (1973) demonstrated that PA *-k- (simple velar) and *-hk- (preaspirated
velar) merged phonetically in Narragansett as /hk/ medially, but can sometimes
be distinguished through the behavior of preceding vowels. PA *-k- shortens a
preceding long vowel; PA *-hk- does not. This provides an indirect method for
reconstructing preaspiration from vowel length patterns --- even though
Williams never writes the /h/.

\subsubsection{Affricates}

\begin{table}[h]
\centering
\begin{tabular}{lll}
\toprule
Williams & IPA & Notes \\
\midrule
ch & /tʃ/ & As in English ``chair'' --- may also represent /t/ palatalized \\
tch & /tʃ/ & As in English ``itch'' \\
ddt, dt & /d/ ~ /t/ ~ /tʲ/ & Complex ch-t sound; among the most ambiguous \\
\bottomrule
\end{tabular}
\caption{Affricate consonants}
\label{tab:1362-affricates}
\end{table}

\subsubsection{Nasals}

\begin{table}[h]
\centering
\begin{tabular}{lll}
\toprule
Williams & IPA & Notes \\
\midrule
m, mm & /m/ & \\
n, nn & /n/ & Before consonant can be syllabic /n̩/ (→ /nə/ for English speakers) \\
n (syllabic) & /n̩/ → /nə/ & Unwritten schwa before n in consonant clusters \\
\bottomrule
\end{tabular}
\caption{Nasal consonants}
\label{tab:1362-nasals}
\end{table}

\subsubsection{Fricatives}

\begin{table}[h]
\centering
\begin{tabular}{lll}
\toprule
Williams & IPA & Notes \\
\midrule
s, ss & /s/ & Single \textit{s} after vowel; geminate \textit{ss} also /s/ \\
sh & /ʃ/ word-initially; /s/ before consonant & Context-dependent ambiguity \\
z & /s/ (not /z/) & English-style voiced s for fortis /s/ \\
sc, shk, sk, shq & /sk/ & Cluster spellings \\
sq & /skʷ/ ~ /x/ & Labialized or guttural depending on following vowel \\
\bottomrule
\end{tabular}
\caption{Fricative consonants}
\label{tab:1362-fricatives}
\end{table}

Goddard \& Bragdon (1988) note that in Massachusett, the cognate system shows
*/s/ leniting to /h/ in intervocallic position --- a sound change that may
also operate in Narragansett but is invisible in Williams' orthography (since
he never writes /h/).

\subsubsection{Glides}

\begin{table}[h]
\centering
\begin{tabular}{lll}
\toprule
Williams & IPA & Notes \\
\midrule
w, ww & /w/ ~ /j/ before e/i & Before e/i can be /j/ (yes) \\
q (before vowel) & /kʷ/ & As in English ``queen'' --- consistent \\
q (before w) & /k/ ~ /kʷ/ ~ /x/ & Context-dependent; may represent guttural \\
\bottomrule
\end{tabular}
\caption{Glide consonants}
\label{tab:1362-glides}
\end{table}

\subsubsection{The Silent /h/ --- Preaspiration}

\textbf{English speakers routinely fail to hear preaspirated /h/ before stops.}
This is the single most important finding for this source. Williams almost
never wrote /h/ where modern reconstructions put preaspiration:

\begin{itemize}
  \item In Narragansett (and other SNEA languages), /h/ commonly precedes /p/,
    /t/, /k/ in word-medial position (e.g., /hkaː/, /hp/, /ht/)
  \item English speakers perceive these as single stops /k/, /p/, /t/ ---
    preaspiration is not a phonemic distinction in English
  \item Williams wrote \textit{k} where the actual pronunciation was /hk/
    (preaspirated k)
  \item \textbf{Confirmed by \textbackslash ph records:} The 3 records with
    \textbackslash ph fields show missing /h/ in Williams
\end{itemize}

\begin{table}[h]
\centering
\begin{tabular}{llll}
\toprule
ID & Williams & \textbackslash ph (phonetic) & Pattern \\
\midrule
3834 & Aaunchemókaw & ąːčimoːhkaːw & \textit{k} → hk \\
4286 & Anakáusichick & anoːhkaːsiːčiːk & \textit{k} → hk \\
4285 & Anakáusu & anoːhkaːsəw & \textit{k} → hk \\
\bottomrule
\end{tabular}
\caption{Evidence from the 3 \textbackslash ph entries in the database}
\label{tab:1362-ph}
\end{table}

All three show the same pattern: Williams writes \textit{k} where the
reconstruction has /hk/. The /h/ is invisible to English ears.

\subsection{Problems Encountered}

\subsubsection{Vowel Quality Ambiguity}

Williams' vowels \textit{a}, \textit{e}, \textit{i}, \textit{o}, \textit{u}
each map to 2--4 possible phonemes. The vowel \textit{a} alone can represent
/ə/, /ɔː/, or /aː/. This extreme ambiguity means that any single Williams
spelling is consistent with multiple phonological interpretations.

\subsubsection{Nasalization Implied but Not Marked}

Sequences like \textit{an}, \textit{aun}, \textit{om} may indicate nasal
vowels /ã/, /õ/ or oral vowel + consonantal \textit{n}. Disambiguating
requires:
\begin{enumerate}
  \item Checking the cognate in Massachusett (Eliot marks nasal /ô/ explicitly)
  \item Checking PA reconstruction (do expected reflexes predict nasalization
    in this environment?)
  \item Checking morphological structure (is the \textit{n} part of a suffix,
    or a nasalization marker?)
\end{enumerate}

\subsubsection{Stop Voicing Neutralized}

\textit{p}, \textit{t}, \textit{k} represent unaspirated stops that English
speakers perceived variably as voiced or voiceless. Williams writes both
\textit{p} and \textit{b}, \textit{t} and \textit{d}, \textit{k} and
\textit{g} in free variation.

\subsubsection{English Spelling Conventions}

English digraphs (\textit{ea}, \textit{oo}, \textit{ch}, \textit{sh}) follow
English values, but inconsistently due to 17th-century English pronunciation
being different from modern English. For example, \textit{ea} in 1643 was
/ɛː/, not modern /iː/.

\subsubsection{Silent e}

Word-final \textit{e} is often silent, but not always. This affects at least
15--20\% of entries. Comparison with Massachusett cognates (Eliot) is the
primary disambiguation strategy.

\subsubsection{Preaspiration /h/ Invisible}

Only 3 of 2,134 records have phonetic (\textbackslash ph) fields. All three
show the same pattern: Williams writes \textit{k} where the reconstruction has
/hk/. If this pattern is representative, >50\% of Williams' stops may be
missing a preceding /h/.

\subsubsection{\textit{sh} Before Consonants}

Williams writes \textit{sh} before consonants where Massachusett cognates
suggest plain /s/: Williams \textit{shp-} corresponds to Massachusett
\textit{sp-}. Costa (2007) notes that this palatalization before consonants is
also attested in Loup A records, suggesting it is a genuine SNEA dialect
feature, not a Williams idiosyncrasy.

\subsubsection{The \textit{an}/\textit{aun} Ambiguity for /ã/}

The most consequential ambiguity: Williams writes both \textit{an} and
\textit{aun} for /ã/, but also writes \textit{an} for plain /an/ (oral a + n).
Disambiguating requires comparative evidence from Massachusett (Eliot marks
nasal /ô/ explicitly), PA reconstruction, and morphological analysis.

\subsection{Research Approach}

\subsubsection{Recommended Methodology}

\textbf{DO NOT USE REGEX} --- the ambiguities are fundamentally
context-dependent. Recommended approach:

\begin{enumerate}
  \item \textbf{Finite State Machine (FSM):} Encode known spelling-to-phoneme
    rules as weighted state transitions, with fallback probabilities for
    ambiguous mappings.
  \item \textbf{ML classifier:} Train on the 3 \textbackslash ph seed examples
    (plus indirect evidence from comparative Algonquian) to predict
    preaspiration placement and vowel quality.
  \item \textbf{Morphological analysis:} Use known Narragansett morphology
    (prefixes, stems, suffixes from Aubin/Hagenau) to segment Williams'
    spellings before applying conversion.
  \item \textbf{Comparative evidence:} Use Proto-Algonquian reconstructions
    and sister languages (Massachusett, Mohegan-Pequot) as priors for
    ambiguous mappings.
\end{enumerate}

\subsubsection{Annotation Plan}

The 3 \textbackslash ph entries alone are insufficient for ML training.
Recommended annotation plan:

\begin{enumerate}
  \item Hand-annotate ~50 representative entries with conversions
  \item Use the annotations as training data
  \item Evaluate on the remaining entries with spot-check verification
  \item Cross-validate against Massachusett cognates (Goddard \& Bragdon 1988)
    for at least 200 entries to calibrate the Narragansett-specific isoglosses
\end{enumerate}

\subsubsection{Cross-Linguistic Cognate Evidence}

Massachusett, the closest sister language to Narragansett, was documented by
John Eliot about 20 years after Williams. Eliot used a different orthographic
convention --- less English-phonetic, more systematic (he devised it for Bible
translation).

\begin{table}[h]
\centering
\begin{tabular}{llll}
\toprule
English & Williams Narragansett & Eliot Massachusett & PA reconstruction \\
\midrule
dog & ayim & arum & *aɬemwa \\
man & weenawash & wénawaš & \\
water & nipp & nippe & *nepyi \\
fire & apeisuwau & apeisuwâ & \\
house & wétu & wétu & *wiːkiwaːmi \\
I (say) & neen & neen & *niːra \\
\bottomrule
\end{tabular}
\caption{Cognate pairs for calibration}
\label{tab:1362-cognates}
\end{table}

\subsubsection{Open Questions for Revitalization}

\begin{enumerate}
  \item \textbf{Vowel length:} Was Narragansett vowel length phonemic (as in
    Massachusett)? Williams' orthography suggests yes (áh vs. a for long vs.
    short /a/), but the marking is inconsistent.
  \item \textbf{Preaspiration:} How pervasive was /h/ before stops? The 3
    \textbackslash ph entries suggest it was common.
  \item \textbf{Nasalization phonemicity:} Was nasalization phonemic (as in
    Massachusett) or allophonic? Costa (2007) suggests phonemic in at least
    some SNEA dialects.
  \item \textbf{Final vowel status:} Did Narragansett drop word-final vowels
    (like Mahican) or retain them (like Massachusett)? Williams' variable
    practice suggests either retention with devoicing or transition state.
\end{enumerate}

\subsection{Conversion Workflow}
\texttt{[Pending implementation --- see Issue \#1362]}

\subsection{Linguist Feedback}
\texttt{[Template --- content pending review]}

\subsection{Related Work: Colonial Recorder Perception Problem}

The core challenge of this research card --- English-speaking recorders
misperceiving and mis-transcribing Algonquian phonemes --- is a known
phenomenon documented across Eastern North America. Two researchers' work is
directly relevant:

\textbf{Dr. Keith Cunningham} (Linguist, Nanticoke Indian Tribe). His doctoral
dissertation \emph{A Phonological Analysis of Nanticoke With Practical
Applications for Language Revitalization} (Georgetown University) analyzes
three 18th-century Nanticoke word lists (Heckewelder 1785) using the same
methodology: historical phonology from colonial recorders, with the same
challenges of multilingual informants and orthographic complexity. His 2024
Algonquian Conference presentation \emph{A Revised Phonological Analysis of
the Heckewelder Vocabulary of Nanticoke} directly parallels the Williams
analysis --- both involve English-speaking recorders working without
standardized orthography, both require PA comparative calibration, and both
feed into practical language revitalization. Cunningham's work demonstrates
that the colonial recorder perception problem extends beyond SNEA into the
broader Eastern Algonquian family (Nanticoke, Maryland Algonquian).

\textbf{Dr. Craig Kopris} (Independent scholar). His paper \emph{Les sons
wyandots perçus par des oreilles étrangères} (Wyandot sounds perceived by
foreign ears, \emph{Recherches amérindiennes au Québec}, 2014) is the exact
framing of the Williams problem: how European recorders systematically
misperceived and mis-transcribed indigenous phonemes. His \emph{Wyandot
Phonology: Recovering the Sound System of an Extinct Language} applies the
same recovery methodology to an unrelated language family (Iroquoian),
demonstrating that the foreign-ear perception problem is universal across
colonial documentation of North American languages. His Dictionary of Wyandot
project (NEH award FN-255576-17) shows the long-term application of these
methods to language revitalization.

The Williams Narragansett corpus (2,134 records, only 3 with phonetic
transcription) represents one of the most extreme cases of this problem ---
but the methodological framework for addressing it is established in the
parallel work of Cunningham (Algonquian) and Kopris (Iroquoian).

\subsection{Bibliography}

\begin{enumerate}

\item Bloomfield, Leonard. (1925). ``On the Sound-System of Central
Algonquian.'' \emph{Language} 1(4): 130--156.

\item Bloomfield, Leonard. (1946). ``Algonquian.'' In Harry Hoijer et al.,
\emph{Linguistic Structures of Native America}, pp. 85--129. Viking Fund
Publications in Anthropology 6. New York: Wenner-Gren Foundation.

\item Bragdon, Kathleen J. (1999). \emph{Native People of Southern New
England, 1500--1650}. Norman: University of Oklahoma Press.

\item Bragdon, Kathleen J. (2009). \emph{Native Writings in Massachusett,
Part II: Linguistic Analysis}. Memoirs of the American Philosophical Society
263. Philadelphia: American Philosophical Society.

\item Costa, David J. (2007). ``The Dialectology of Southern New England
Algonquian.'' In H.C. Wolfart, ed., \emph{Papers of the 38th Algonquian
Conference}, pp. 81--127. Winnipeg: University of Manitoba.

\item Cowan, William. (1973). ``The Development of PA *-k- and *-hk- in
Narragansett.'' \emph{International Journal of American Linguistics} 39(2):
90--93.

\item Cowan, William. (1975). ``PA *-θ- and *-t- in Narragansett.''
\emph{International Journal of American Linguistics} 41(2): 178--181.

\item Goddard, Ives. (1978). ``Eastern Algonquian Languages.'' In Bruce G.
Trigger, ed., \emph{Handbook of North American Indians, Vol. 15: Northeast},
pp. 70--77. Washington: Smithsonian Institution.

\item Goddard, Ives. (1996). ``The Description of the Native Languages of
North America Before Boas.'' In Ives Goddard, ed., \emph{Handbook of North
American Indians, Vol. 17: Languages}, pp. 17--42. Washington: Smithsonian
Institution.

\item Goddard, Ives, and Kathleen J. Bragdon. (1988). \emph{Native Writings
in Massachusett.} 2 vols. Memoirs of the American Philosophical Society 185.
Philadelphia: American Philosophical Society.

\item Hagenau, Walter. (1962). \emph{A Morphological Study of Narragansett.}
M.A. thesis, University of North Dakota.

\item O'Brien, Frank Waobu. \emph{Appendix: Guide to Historical Spellings
and Sounds in New England Algonquian Languages.} Aquidneck Indian Council.

\item Pentland, David H. (1979). \emph{Algonquian Historical Phonology.}
Ph.D. dissertation, University of Toronto.

\item Siebert, Frank T., Jr. (1975). ``Resurrecting Virginia Algonquian from
the Dead: The Reconstituted and Historical Phonology of Powhatan.'' In James
M. Crawford, ed., \emph{Studies in Southeastern Indian Languages}, pp.
285--453. Athens: University of Georgia Press.

\item Trumbull, James Hammond. (1870). ``Notes on Forty Versions of the
Lord's Prayer in Algonkin Languages.'' \emph{Transactions of the American
Philological Association} (1869--1896): 1--20.

\item Williams, Roger. (1643). \emph{A Key into the Language of America.}
London: Gregory Dexter.

\end{enumerate}

% Wampanoag (Trumbull 1903) — Source Workflow
% Issue #1363
% Status: Research complete, implementation pending

\section{Wampanoag: Trumbull/Eliot 1903}
\label{sec:1363}

\subsection{Source Description}

\textbf{Recorder:} James Hammond Trumbull (1821--1897), drawing on John Eliot's
17th-century Massachusett writing system.
\textbf{Source:} \emph{Natick Dictionary}, Bureau of American Ethnology
Bulletin 25 (1903, posthumous).
\textbf{Language:} Wampanoag (Massachusett) [wam].
\textbf{Corpus:} 4,491 records (largest in the SNEA collection).
\textbf{Orthography:} Eliot orthography --- the first standardized writing
system for an Algonquian language, developed by John Eliot and Massachusett
collaborators (Job Nesutan, John Sassamon) in the 1650s--1660s.

\subsubsection{Recorder Biography: James Hammond Trumbull}

James Hammond Trumbull (1821--1897) was an American philologist, historian,
and bibliographer. He served as Connecticut State Librarian (1854--1897) and
was a founding member of the American Philological Association. His work on
the Algonquian languages of southern New England spanned four decades,
beginning with his early comparative study \emph{On the Best Method of
Studying the North American Languages} (1859).

Trumbull's major Algonquianist works include:
\begin{itemize}
  \item \emph{The Composition of Indian Geographical Names} (1870) ---
    Connecticut Historical Society
  \item \emph{Indian Names of Places in Connecticut} (1881) --- systematic
    geographic toponymy
  \item \emph{Natick Dictionary} (published posthumously 1903) --- his
    magnum opus
\end{itemize}

\subsubsection{The Lost Thesis Controversy}

Trumbull delivered a landmark paper before the American Philological
Association titled \emph{The Method of the Indian Languages as Exemplified by
the Massachusett Dialect} --- based on a ``lost thesis'' he had written at
Yale in the 1840s. The paper argued for the polysynthetic structure of
Algonquian languages and established Trumbull's reputation as the leading
American philologist of Native languages before the Boasian revolution.

Trumbull's methodology combined:
\begin{enumerate}
  \item \textbf{Text-based extraction:} Exhaustively reading through all of
    Eliot's Massachusett translations (Bible, primers, catechisms) and
    manually extracting every lexical item onto slips
  \item \textbf{Cross-referencing with cognates:} Comparing Massachusett
    forms with related languages (Narragansett via Williams 1643,
    Mohegan-Pequot via field notes, Abenaki, Micmac, etc.)
  \item \textbf{Morphemic analysis:} Breaking down polysynthetic forms into
    constituent morphemes --- a radical approach for its time
  \item \textbf{Etymological comparison:} Identifying cognate sets across the
    Eastern Algonquian languages
\end{enumerate}

\subsubsection{Posthumous Publication}

Trumbull died in 1897 with the dictionary manuscript substantially complete
but not yet in press. Edward Everett Hale (the Unitarian minister and author)
supervised the final editing at the Bureau of American Ethnology. The preface
acknowledges that some inconsistencies in cross-references and source
abbreviations may reflect the incomplete final revision by Trumbull himself.
The published dictionary includes 4,491 main entries.

\subsubsection{The Eliot Orthography --- Historical Origins}

John Eliot (1604--1690) arrived in Massachusetts Bay in 1631 and began
preaching to the Massachusett people in the 1640s. His language teacher was
\textbf{Job Nesutan}, a Massachusett man who served as Eliot's primary
linguistic informant and translator. Eliot also worked with \textbf{John
Sassamon} (a Massachusett convert educated at Harvard) who later became a
crucial intermediary.

Eliot's orthography was developed in the 1650s as he prepared his translation
of the Bible. It was:
\begin{itemize}
  \item \textbf{Phonemically oriented} for its time --- Eliot attempted to
    represent the sounds of Massachusett systematically, not just through
    English ear-approximations
  \item \textbf{English-based} --- using English letter values as the starting
    point, then extending them
  \item \textbf{Regularized} --- unlike Wood (1634) or Winslow (1624), Eliot
    chose consistent spellings for the same morpheme across occurrences
  \item \textbf{Influenced by Hebrew and Greek grammar} --- Eliot's
    \emph{Indian Grammar Begun} (1666) explicitly draws parallels between
    Massachusett structures and Biblical Hebrew
\end{itemize}

The \emph{Mamusse Wunneetupanatamwe Up-Biblum God} (1,180 pages) was printed
by Samuel Green and Marmaduke Johnson in Cambridge, Massachusetts. It was the
first complete Bible printed in the Americas in a Native language. Production
required a specially cast type font (some characters unique to Massachusett),
translators and typesetters literate in the orthography, and Massachusett
readers to proofread.

\subsection{Vowel Inventory}

\subsubsection{Simple Vowels}

The core challenge of the Eliot orthography is that 5 vowel letters must
represent a larger Massachusett vowel inventory (likely 8--10 vowel phonemes
including length distinctions).

\begin{table}[h]
\centering
\begin{tabular}{lllll}
\toprule
Eliot & Phoneme & IPA & Example & Notes \\
\midrule
a & Short a / Schwa & /a/ ~ /ə/ & \textit{wadchu} (mountain) & Unstressed → [ə] \\
a & Long a & /aː/ & \textit{mahche} (already) & Lengthened \\
e & Long e & /iː/ & \textit{wetu} (house) & Two different underlying vowels \\
e & Short e & /ɛ/ & \textit{penashimwog} & \\
e & Final schwa & /ə/ & \textit{nish} (two) & Often silent but written \\
i & Long i & /iː/ & \textit{pish} (future) & \\
i & Short i & /ɪ/ & \textit{nippe} (water) & \\
o & Long o & /uː/ & \textit{noh} (father) & \\
o & Short o & /ɔ/ & \textit{nuppa} (sleep) & \\
o & Schwa & /ə/ & \textit{onk} (and) & Unstressed reduction \\
u & Schwa & /ə/ & \textit{ummon} (breast) & Unstressed position \\
u & Short u & /ʌ/ & \textit{musqui} (grass) & \\
\bottomrule
\end{tabular}
\caption{Simple vowel mappings in Eliot's orthography}
\label{tab:1363-vowels}
\end{table}

\subsubsection{Vowel Digraphs}

\begin{table}[h]
\centering
\begin{tabular}{llll}
\toprule
Eliot Digraph & IPA & Example & Notes \\
\midrule
ee & /iː/ & \textit{neen} (I), \textit{keetompauog} (men) & Consistent --- most reliable \\
oo & /uː/ & \textit{woskot} (firewood), \textit{sook} (pouring) & Consistent \\
⟨ꝏ⟩ / ∞ & /uː/ or complex & See §\ref{sec:1363-infinity} & The debated symbol \\
aa & /aː/ & \textit{waam} (long), \textit{maangu} (wolf) & \\
ii & /iː/ & \textit{niin} (my) & Less common \\
au & /ɔ/ & \textit{mauche} (come) & Rare \\
ai & /aːj/ & \textit{wai} (woe) & Diphthong \\
eu & /eːw/ & \textit{seup} (river) & Possibly /əw/ \\
\bottomrule
\end{tabular}
\caption{Vowel digraphs}
\label{tab:1363-digraphs}
\end{table}

\subsection{Consonant Inventory}

\begin{table}[h]
\centering
\begin{tabular}{llll}
\toprule
Eliot & IPA & Example & Notes \\
\midrule
p & /p/ ~ /b/ & \textit{peche} (tomorrow), \textit{nuppa} (sleep) & Unaspirated, lenis medial \\
t & /t/ ~ /d/ & \textit{tummunk} (worm), \textit{wetu} (house) & Unaspirated, medially lenis \\
k & /k/ ~ /ɡ/ & \textit{kenug} (fish), \textit{askkuhnk} (serpent) & Unaspirated \\
c (a/o/u) & /k/ & \textit{commaco} (feast), \textit{cutsh} (speak) & Hard c --- English convention \\
c (e/i) & /s/ & \textit{ce} (thou), \textit{nuppacoh} (chaff) & Soft c --- English convention \\
ch & /tʃ/ & \textit{choh} (child), \textit{mache} (go) & Affricate \\
sh & /ʃ/ & \textit{asuhsh} (fine), \textit{shuckquan} (rain) & \\
s & /s/ & \textit{simma} (various) & Word-initial s often written ss \\
ss & /sː/ ~ /s/ & \textit{massa} (great) & Geminate \\
m & /m/ & \textit{matta} (no) & \\
n & /n/ & \textit{neen} (I) & \\
w & /w/ & \textit{wetu} (house), \textit{waantam} (wise) & \\
y & /j/ & \textit{yeuyeu} (gradually) & \\
q & /kʷ/ & \textit{qussuk} (stone) & Labialized k --- from PA *kw \\
h & /h/ & \textit{ohke} (earth), \textit{kehtan} (great) & Also marks preaspiration \\
ht & /ht/ ~ /h/ & \textit{kehtan} (great), \textit{kehte} (truly) & Preaspirated /t/ \\
hp & /hp/ ~ /h/ & \textit{wuhpom} (honey) & Preaspirated /p/ \\
hk & /hk/ ~ /h/ & \textit{ohke} (earth) & Preaspirated /k/ \\
' (apostrophe) & /h/ ~ /ʔ/ & \textit{m'tah} (be tired) & Elided syllable, often = /h/ \\
\bottomrule
\end{tabular}
\caption{Consonant inventory}
\label{tab:1363-consonants}
\end{table}

\subsection{The ∞ Symbol --- Historical Typography and Phonetics}
\label{sec:1363-infinity}

\subsubsection{The Misidentification}

The ∞ symbol (Unicode INFINITY, U+221E) is one of the most distinctive
features of the database's Trumbull entries. However, this is almost
certainly a \textbf{typographic misidentification} in the digitization
process. The original Trumbull dictionary uses a different symbol:

The actual symbol: Eliot's Bible and Trumbull's dictionary use a ligature of
two ⟨o⟩ letters joined --- what typographers call an ``oo ligature'' ---
Unicode ⟨ꝏ⟩ (U+A74F, LATIN SMALL LETTER OO). In the early BAE Bulletin 25
printed type, this appeared as a figure-8-like symbol because the two o's
were cast as a single piece of type:

\begin{quote}
Original Eliot/Trumbull: ⟨ꝏ⟩ (oo ligature)\\
What it looks like at small sizes: ∞ (infinity-like)\\
What DB has: ∞ (Unicode INFINITY)
\end{quote}

The digitizer apparently did not recognize the typographic convention and
encoded it as the visually-similar infinity symbol.

\subsubsection{What ⟨ꝏ⟩ Actually Represents}

In Eliot's orthography, ⟨ꝏ⟩ represents \textbf{long /uː/} --- specifically,
it is Eliot's way of writing the long u sound, contrasting with ⟨oo⟩ which
may represent a different vowel quality or length. Eliot's own description in
\emph{The Indian Grammar Begun} (1666):

\begin{quote}
``oo, and ꝏ, are two differing vowels, or rather the same vowel differently
sounded; oo is short and ꝏ is long.''
\end{quote}

This suggests ⟨oo⟩ = /u/ (short) and ⟨ꝏ⟩ = /uː/ (long). However, database
evidence showing 654 occurrences of the symbol (now encoded as ∞) suggests a
more complex distribution. In many entries, ∞ appears where a
preaspirated/breathy vowel sequence would be expected phonologically.

\subsubsection{The ∞→oozzz Pipeline Mapping}

Per the lessons-learned documentation, the mapping ∞ → oozzz was established
empirically based on:
\begin{enumerate}
  \item \textbf{Distributional patterns:} ∞ appears in contexts where /u/ +
    /h/ + /z/ would be expected from PA reconstruction
  \item \textbf{Morphophonemic alternation:} Where cognate forms in other SNEA
    languages show /u/ + preaspiration + sibilant sequences
  \item \textbf{Goddard \& Bragdon (1988) transcription conventions:} Their
    modern phonemic transcription suggests underlying forms with /uːh/
    sequences
\end{enumerate}

The \texttt{oozzz} string is therefore a \textbf{working normalization} --- it
captures the phonetic complexity (length + breathiness + sibilant quality)
that the original ⟨ꝏ⟩ was designed to encode, even if the exact phonetics
remain debated.

\begin{table}[h]
\centering
\begin{tabular}{lll}
\toprule
Pattern & Count (approx.) & Example \\
\midrule
∞ (bare) & ~500 & \textit{achm∞wonk} (news) \\
∞̌ (with breve) & ~50 & \textit{∞̌ppa} (various) --- phonetically reduced \\
∞ in compounds & ~100 & Various morpheme-internal occurrences \\
\midrule
Total & ~654 & 14.6\% of 4,491 records \\
\bottomrule
\end{tabular}
\caption{∞ in the database corpus}
\label{tab:1363-infinity-counts}
\end{table}

\subsection{Entry Structure --- Detailed Parsing Guide}

Trumbull entries have a complex multi-part structure parsed into a single
\texttt{lx} field. Understanding this structure is essential for FSM-based
parsing.

\subsubsection{Standard Entry Template}

\begin{quote}
\texttt{*|achm∞wonk|, vbl. n. `news'    C.}
\end{quote}

Field breakdown:
\begin{itemize}
  \item \texttt{*} --- Cross-reference marker
  \item \texttt{|...|} --- Citation form delimiters (morphemic boundaries)
  \item \texttt{,} --- Field separator
  \item \texttt{vbl. n.} --- Grammatical label (verbal noun)
  \item \texttt{`news'} --- Gloss in backtick quotes
  \item \texttt{C.} --- Source abbreviation
\end{itemize}

\subsubsection{Entry Variants}

Root/compound form (no gloss):
\begin{quote}
\texttt{|-adchaubuk|, in comp. words, root, or roots}
\end{quote}

Cross-reference entry:
\begin{quote}
\texttt{|asqhutt∞che| `whilst'    C. = |asq-utt∞che|.}
\end{quote}

Multiple glosses:
\begin{quote}
\texttt{|wame|, adj. `all, whole, every'    C.}
\end{quote}

\subsubsection{Source Abbreviations}

Trumbull used a system of abbreviations to tag each entry's source:

\begin{table}[h]
\centering
\begin{tabular}{lll}
\toprule
Abbreviation & Source & Approx. \% \\
\midrule
C. & Eliot's Catechism (various editions) & ~40\% \\
B. & Eliot's Bible (1663, 1685) & ~30\% \\
P. & Psalms translated by Eliot & ~10\% \\
G. & Eliot's Grammar (1666) & ~5\% \\
Pr. & Eliot's Primer(s) & ~5\% \\
Pl. & Plymouth/colonial records & ~3\% \\
(unmarked) & Unknown or Trumbull's own reconstruction & ~7\% \\
\bottomrule
\end{tabular}
\caption{Source abbreviations in Trumbull}
\label{tab:1363-sources}
\end{table}

\subsubsection{Grammatical Labels}

\begin{table}[h]
\centering
\begin{tabular}{lll}
\toprule
Abbreviation & Full Form & Meaning \\
\midrule
vbl. n. & verbal noun & Action nominalization \\
adj. & adjective & Modifier \\
pro. & pronoun & Personal/demonstrative pronoun \\
adv. & adverb & Modifier of verb/adjective \\
prep. & preposition & Relational word \\
conj. & conjunction & Connective \\
interj. & interjection & Exclamation \\
n. & noun & Substantive \\
v. & verb & Verb (unmarked for transitivity) \\
v. a. & verb active & Transitive verb \\
v. n. & verb neuter & Intransitive verb \\
part. & particle & Uninflected word \\
s. & substantive & Noun (alternative to n.) \\
pl. & plural & Plural form \\
comp. & compound & Compound word \\
in comp. & in compound & Used only in compounds \\
dim. & diminutive & Diminutive form \\
\bottomrule
\end{tabular}
\caption{Grammatical label abbreviations}
\label{tab:1363-grammar}
\end{table}

\subsection{Problems Encountered}

\subsubsection{The ∞ (ꝏ) Symbol Misidentification}

The symbol was historically misidentified as an infinity symbol (∞). It is
Eliot's ⟨oo⟩ ligature (U+A74F), representing long /uː/. The digitization
error encoding it as Unicode INFINITY (U+221E) has propagated through the
database, requiring normalization in the conversion pipeline.

\subsubsection{Preaspiration Gap}

Eliot writes \textit{hk/hp/ht} only ~15\% of expected positions. Edwards
(1787 Mahican) is the only colonial source that consistently marks
preaspiration. For Trumbull/Eliot entries, preaspiration /h/ must be inferred
from comparative evidence (Edwards' Mahican data, PA reconstructions,
morphotactic position) rather than read directly from the orthography.

\subsubsection{Apostrophe Ambiguity}

The apostrophe may represent /h/, /ʔ/, or vowel elision --- no consistent
rule. The apostrophe function is ambiguous without comparative evidence and
cannot be resolved solely from within the Trumbull entry.

\subsubsection{Dialect Stratification}

Variation between Natick (Eliot) and Pokanoket (Winslow/Wood) may be
orthographic, phonological, or dialectal. The \textit{achm∞wonk} vs.
\textit{hobbamock} divergence (the ``devil problem'') suggests real dialect
differentiation between the ``praying town'' dialect of Natick and the
Plymouth-area Pokanoket dialect.

\subsubsection{Vowel Quality}

The five-letter vowel system (a, e, i, o, u) must represent 8--10 phonemes.
Disambiguation requires distributional analysis, PA correspondence, stress
position, and morphological position.

\subsubsection{\textit{c} Ambiguity}

Before back vowels (a, o, u) → /k/; before front vowels (e, i) → /s/.
Exception: some irregular forms where \textit{c} = /k/ before e. Word-final
\textit{c} → /k/.

\subsubsection{Boundary Markers}

The \texttt{|...|} pipe convention is near-universal (4,488 of 4,491 entries)
but embedded glosses inside pipes are rare, cross-reference syntax (=) may
span pipe boundaries, and hyphens within pipes indicate morpheme boundaries,
not word boundaries.

\subsubsection{The Devil Problem: \textit{achm∞wonk} vs. \textit{hobbamock}}

One of the most revealing lexical puzzles in the corpus involves the word for
`devil':

\begin{table}[h]
\centering
\begin{tabular}{lll}
\toprule
Source & Word & Notes \\
\midrule
Winslow (1624) & \textit{hobbamock} & Initial h written; also a proper name \\
Wood (1634) & \textit{abamocho}, \textit{abomocho} & No initial h; o-final \\
Williams (1643) & \textit{ábammocho} & Narragansett; consistent with Wood \\
Trumbull/Eliot & \textit{achm∞wonk} & Different root; \textit{ch} not \textit{b/p} \\
\bottomrule
\end{tabular}
\caption{The devil problem across sources}
\label{tab:1363-devil}
\end{table}

This shows that Trumbull/Eliot may be drawing from a different lexical
stratum (the ``praying town'' dialect of Natick) than Winslow/Wood
(Plymouth-area Pokanoket). The lexical difference is not just orthographic
--- it may reflect real dialect divergence.

\subsection{Research Approach}

\subsubsection{Recommended Methodology}

\textbf{DO NOT USE REGEX} --- the ∞ symbol, pipe boundaries, embedded
glosses, and cross-reference syntax require context-sensitive parsing.

\textbf{Stage 1: Preprocessing FSM.} Tokenize Trumbull entries by:
\begin{enumerate}
  \item Recognizing \texttt{|...|} citation boundaries and stripping pipes
  \item Identifying ∞ and replacing with \texttt{oozzz} (first normalization)
  \item Extracting backtick-delimited glosses
  \item Isolating grammatical abbreviations (vbl. n., C., etc.)
  \item Handling \texttt{*} prefix and \texttt{=} cross-reference syntax
  \item Separating source abbreviations (C., B., etc.) for provenance tracking
\end{enumerate}

\textbf{Stage 2: Phonemic Conversion FSM.} Apply Eliot-to-phoneme mapping
rules:
\begin{enumerate}
  \item Map Eliot vowel digraphs (ee, oo, aa, etc.) to phonemes
  \item Handle \textit{c}→\textit{k}/\textit{s} disambiguation based on
    following vowel
  \item Infer preaspiration /h/ from written \textit{hk}, \textit{hp},
    \textit{ht} clusters (direct evidence), morphotactic position (coda of
    stressed syllable, before /p t k/), comparative evidence from Edwards
    (1787) Mahican data, and PA reconstructions showing */h/ in corresponding
    position
  \item Resolve ∞→\texttt{oozzz} (already mapped in Stage 1)
  \item Resolve apostrophe to /h/ or /ʔ/ based on context and comparative
    evidence
\end{enumerate}

\textbf{Stage 3: Comparative Fallback.} For ambiguous vowel quality (the
5-vowel→8--10-phoneme mapping problem):
\begin{enumerate}
  \item PA comparison: Use Aubin (1975) reconstructions as primary
    disambiguation
  \item Sister-language evidence: Williams (Narragansett), Edwards (Mahican),
    Fielding/Prince-Speck (Mohegan-Pequot)
  \item Known cognate patterns: Eliot orthography ↔ modern Wampanoag
    revitalization (Weeden, Baird)
  \item Goddard \& Bragdon (1988) transcriptions: Modern phonemic analyses of
    Massachusett texts
\end{enumerate}

\textbf{Stage 4: Statistical/Distributional Analysis.} The large corpus (4,491
entries) enables statistical approaches not viable for smaller sources:
\begin{enumerate}
  \item Vowel quality prediction model based on surrounding consonant
    environment and stress position
  \item Preaspiration probability model based on morphotactic position and PA
    correspondence
  \item ∞ behavior model: Distributional patterns of ∞ across phonological
    environments
\end{enumerate}

\subsubsection{Proto-Algonquian Correspondences}

\begin{table}[h]
\centering
\begin{tabular}{llll}
\toprule
PA & Massachusett Reflex & Eliot Grapheme & Example \\
\midrule
*a & /a/ & a & *aʔlapy → \textit{anupp} (bowstring) \\
*e & /ə/ & u, a & *eːθkweːwa → \textit{squa} (woman) \\
*i & /i/ & i, e & *iːli- → \textit{en-} (thus) \\
*o & /ɔ/ & o, a & *oːθaːna → \textit{-adchan} \\
*aː & /aː/ & aa, a & *aːmoːwa → \textit{amaug} (bee) \\
*eː & /iː/ & ee & *eːhšwa → \textit{ash} (three) \\
*iː & /iː/ & ee & *iːniwa → \textit{neen} (I) \\
*oː & /uː/ & oo, ꝏ & *oːtaːna → \textit{adchan} \\
*p & /p/ & p & *pexpenya → \textit{pepeng} \\
*t & /t/ & t & *taːnya → \textit{ta} (there) \\
*k & /k/ & k, c & *kaːkinoː → \textit{kakine} (all) \\
*m & /m/ & m & *maškyekwi → \textit{maskiht} (medicine) \\
*n & /n/ & n & *naːpeːwa → \textit{nap} (man) \\
*s & /s/ & s, ss & *soːkya → \textit{sohki} (hard) \\
*š & /ʃ/ & sh & *eːšpowaːki → \textit{ashpowau} (above) \\
*θ & /s/ & s & *eːθkweːwa → \textit{squa} (woman) \\
*l & /n/ & n & *laːkeːwa → \textit{noh} (father) \\
*w & /w/ & w & *waːposwa → \textit{womp} (white) \\
*y & /j/ & y & *ayaːpeːwa → \textit{ayim} (man?) \\
*h & /h/ & h, (unwritten) & *aːhšoːwa → \textit{ashsh} (?) \\
\bottomrule
\end{tabular}
\caption{PA to Massachusett regular correspondences}
\label{tab:1363-pa}
\end{table}

\subsubsection{Confirmed Cognate Mappings for Pipeline Testing}

\begin{table}[h]
\centering
\begin{tabular}{llll}
\toprule
Gloss & Trumbull/Eliot (expected) & Confidence & PA Source \\
\midrule
`no, not' & \textit{matta} & High & PA *mati- \\
`woman' & \textit{squa} & High & PA *eːθkweːwa \\
`man' & \textit{nap} & High & PA *naːpeːwa \\
`house' & \textit{wetu} & High & PA *wiːkiwaːmi \\
`water' & \textit{nippe} & High & PA *nepyi \\
`stone' & \textit{qussuk} & High & PA *aːsenya \\
`fire' & \textit{nuta} (?) & Moderate & PA *eškw- \\
`three' & \textit{ash} & High & PA *nšwaːśika \\
`white' & \textit{womp(ish)} & High & PA *waːposwa (rabbit → white) \\
`devil' & \textit{achm∞wonk} & Moderate & See §\ref{tab:1363-devil} \\
`snake' & \textit{askkuhnk} (?) & Moderate & PA *ašk-? \\
\bottomrule
\end{tabular}
\caption{Confirmed cognate mappings}
\label{tab:1363-cognates}
\end{table}

\subsection{Conversion Workflow}
\texttt{[Pending implementation --- see Issue \#1363]}

\subsection{Linguist Feedback}
\texttt{[Template --- content pending review]}

\subsection{Related Work: Colonial Recorder Perception Problem}

The core challenge of this research card --- English-speaking recorders
misperceiving and mis-transcribing Algonquian phonemes --- is a known
phenomenon documented across Eastern North America. Two researchers' work is
directly relevant:

\textbf{Dr. Keith Cunningham} (Linguist, Nanticoke Indian Tribe). His doctoral
dissertation \emph{A Phonological Analysis of Nanticoke With Practical
Applications for Language Revitalization} (Georgetown University) analyzes
three 18th-century Nanticoke word lists (Heckewelder 1785) using the same
methodology: historical phonology from colonial recorders, with the same
challenges of multilingual informants and orthographic complexity. His 2024
Algonquian Conference presentation \emph{A Revised Phonological Analysis of
the Heckewelder Vocabulary of Nanticoke} directly parallels the Trumbull/Eliot
analysis --- both involve working with standardized missionary orthographies
(Eliot for Massachusett, Heckewelder for Nanticoke) and both require PA
comparative calibration. Cunningham's work demonstrates that the colonial
recorder perception problem extends beyond SNEA into the broader Eastern
Algonquian family.

\textbf{Dr. Craig Kopris} (Independent scholar). His paper \emph{Les sons
wyandots perçus par des oreilles étrangères} (Wyandot sounds perceived by
foreign ears, \emph{Recherches amérindiennes au Québec}, 2014) is the exact
framing of the Trumbull/Eliot problem: how European recorders systematically
misperceived and mis-transcribed indigenous phonemes. His \emph{Wyandot
Phonology: Recovering the Sound System of an Extinct Language} applies the
same recovery methodology to an unrelated language family (Iroquoian),
demonstrating that the foreign-ear perception problem is universal across
colonial documentation of North American languages. His Dictionary of Wyandot
project (NEH award FN-255576-17) shows the long-term application of these
methods to language revitalization.

The Trumbull corpus (4,491 records, the largest in the SNEA collection)
benefits from Eliot's relatively systematic orthography --- but the
preaspiration gap and vowel ambiguity problems documented in this section are
the same foreign-ear perception issues that Cunningham and Kopris address in
their respective language families.

\subsection{Open Questions for Linguist Review}

\begin{enumerate}
  \item \textbf{∞(ꝏ) phonetics:} Is the \texttt{oozzz} mapping phonetically
    accurate, or should it be revised to \texttt{uu} with a diacritic for
    preaspiration? The \texttt{zzz} component seems ad-hoc.
  \item \textbf{Preaspiration inference rules:} Can a reliable set of
    morphotactic rules predict when preaspiration /h/ was present but unwritten
    in Eliot's orthography? Or does this require per-entry comparative
    verification?
  \item \textbf{Dialect stratification:} How much of the variation between
    Trumbull/Eliot (Natick) and Winslow/Wood (Pokanoket) is orthographic vs.
    phonological vs. dialectal? The \textit{achm∞wonk} vs. \textit{hobbamock}
    divergence suggests real dialect differentiation.
  \item \textbf{Apostrophe resolution:} Is there a predictable pattern for when
    \texttt{'} = /h/ vs. /ʔ/ vs. simple vowel elision? Edwards (Mahican)
    apostrophe usage may provide clues.
  \item \textbf{Eliot's \textit{c} before \textit{e}:} Are there identifiable
    exceptions to the \textit{c} before front vowels → /s/ rule?
  \item \textbf{Edwards cross-over:} How much does Edwards' (1787) Mahican
    preaspiration system overlap with Massachusett? If the systems are
    parallel, Edwards' data can directly inform Trumbull /h/ inference.
\end{enumerate}

\subsection{Bibliography}

\begin{enumerate}

\item Aubin, George F. (1975). \emph{A Proto-Algonquian Dictionary.} National
Museum of Man, Mercury Series, Canadian Ethnology Service Paper No. 39.
Ottawa: National Museums of Canada.

\item Costa, David J. (2007). ``The Dialectology of Southern New England
Algonquian.'' In H.C. Wolfart, ed., \emph{Papers of the 38th Algonquian
Conference}, pp. 81--127. Winnipeg: University of Manitoba.

\item Eliot, John. (1663). \emph{Mamusse Wunneetupanatamwe Up-Biblum God.}
Cambridge, Massachusetts: Samuel Green and Marmaduke Johnson.

\item Eliot, John. (1666). \emph{The Indian Grammar Begun.} Cambridge,
Massachusetts.

\item Goddard, Ives, and Kathleen J. Bragdon. (1988). \emph{Native Writings
in Massachusett.} 2 vols. Memoirs of the American Philosophical Society 185.
Philadelphia: American Philosophical Society.

\item O'Brien, Frank Waobu. \emph{Appendix: Guide to Historical Spellings
and Sounds in New England Algonquian Languages.} Aquidneck Indian Council.

\item Trumbull, James Hammond. (1903). \emph{Natick Dictionary.} Bureau of
American Ethnology Bulletin 25. Washington: Smithsonian Institution.

\end{enumerate}

% Mohegan-Pequot (Fielding 2013) — Source Workflow
% Issue #1364
% Status: Research complete, implementation pending

\section{Mohegan-Pequot: Fielding 2013}
\label{sec:1364}

\subsection{Source Description}

\textbf{Recorder:} Stephanie Fielding (Cits Pátunáhshô), Mohegan tribal
member and linguist, with David Costa (Myaamia Center, Miami University of
Ohio).
\textbf{Source:} \emph{Modern Mohegan Dictionary} (2006, revised 2013).
\textbf{Language:} Mohegan-Pequot [xpq] (Y-dialect of SNEA).
\textbf{Corpus:} 145 records (FF-marked entries from Fidelia Fielding's
diaries).
\textbf{Orthography:} Modern practical orthography --- 12 consonant + 6 vowel
symbols, 1:1 phoneme-grapheme correspondence, circumflex for length.

\subsubsection{Fidelia A. H. Fielding (1827--1908)}

Fidelia A. Hoscott Fielding (Mohegan name: \textbf{Jits Bodunaxa} ``Flying
Bird'') was the last fluent speaker of the Mohegan-Pequot language. Born in
1827 at Mohegan, near Norwich, Connecticut, she learned Mohegan as a child
from her grandmother, \textbf{Martha Uncas Shantup} (1761--1859), a direct
descendant of Sachem Uncas. As the number of fluent speakers dwindled through
the late 19th century, Fielding maintained her language by speaking with her
sister and, later, by talking to herself (Prince \& Speck 1904:18--19).

\subsubsection{The Fielding Diaries}

Between 1902 and 1904, Fielding kept a series of diaries written in the
Mohegan language. These are the only known surviving documents produced by a
native speaker of Mohegan-Pequot. The corpus consists of:

\begin{itemize}
  \item \textbf{Three diaries} and a copy of the Lord's Prayer in Mohegan
    (Cornell University Library, acquired 2004 from the Huntington Free
    Library)
  \item A fourth diary was lost in a fire at Columbia University around 1906
    (Granberry 2003:10)
  \item A sermon text published by Prince and Speck (1903) in \emph{American
    Anthropologist}, vol. 5, pp. 193--212
  \item A long narrative tale recorded in phonetic notation by Frank Speck
    (Speck 1904)
\end{itemize}

The diaries were written in Fielding's own orthography --- an English-based
ad-hoc spelling system reflecting her pronunciation. After her death in 1908,
the diaries passed to Frank Speck at the University of Pennsylvania, and later
to Cornell University. In November 2020, Cornell repatriated the diaries to
the Mohegan Tribe in a formal ceremony, where they now reside in the Mohegan
Archives at Uncasville, Connecticut.

\subsubsection{Corpus Structure}

The Brothertown corpus under analysis contains \textbf{145 records} from the
Fielding dictionary. This is a subset of the full dictionary (which numbers
several thousand entries, many reconstructed via comparative Algonquian
methods by David Costa). Records marked \textbf{FF} in the dictionary
indicate words attested directly in Fidelia Fielding's diaries. The 145
records constitute our working corpus because they have explicit
Mohegan-to-phoneme field mappings, whereas the larger set has modern
orthography forms requiring parsing.

\subsubsection{Comparative Significance}

Prince and Speck (1904:18) explicitly compare Fielding to \textbf{Dorothy
Pentreath}, the last speaker of Cornish, noting the ``strangely pathetic''
nature of a dying language spoken only by a single elderly woman who kept it
alive by talking to herself. Despite the ``last stages of decay'' they
observed, they found that ``enough of the original phonetics and grammatical
phenomena'' remained to judge the language's structure.

\subsubsection{Stephanie Fielding's Grammatical Analysis}

Stephanie Fielding (Cits Pátunáhshô), a full-blooded Mohegan tribal member
and descendant of Fidelia Fielding by marriage (her great-great-grandfather's
aunt), served as the Mohegan Tribe's linguist and authored the \emph{Modern
Mohegan Dictionary} (2006, revised 2013). The dictionary was reviewed by David
Costa who researched the lexicon and grammar, developed the Mohegan alphabet,
translated scripts and wrote sample sentences, and proofed the grammar
paradigms.

The dictionary includes:
\begin{enumerate}
  \item \textbf{Pronunciation Guide} --- the practical orthography with
    phoneme mappings
  \item \textbf{Mohegan Grammar Paradigms} --- a systematic grammatical
    sketch covering pronouns, nouns, verbs, and particles
  \item \textbf{Mohegan to English Dictionary} --- stems with part-of-speech
    codes and inflections
  \item \textbf{English to Mohegan Word Finder} --- English lookup with
    Mohegan stems
\end{enumerate}

Fielding's grammar covers the major Algonquian grammatical categories:
pronouns (prefixes and suffixes attached to verb stems, pro-drop/polysynthetic),
nouns (two genders: Animate vs. Inanimate, with plural, obviative, locative,
and absentative inflections), verbs (four orders: Independent, Conjunct,
Imperative, Participial, with transitivity marking VAI, VII, VTA, VTI), and
derivation (person-marking prefixes \textit{nu-} 1st person, \textit{ku-}
2nd person).

\subsubsection{Fielding's Methodology}

Stephanie Fielding worked extensively from texts, particularly the
\textbf{Eliot Bible} (John Eliot's 1663 Massachusett translation), using a
concordance to find Wôpanâak cognates and applying Mohegan-specific sound
changes to reconstruct Mohegan forms. This comparative method is standard for
Algonquian language reclamation. She also drew from Fidelia Fielding's diaries
(primary source, marked FF), sister languages (Wôpanâak, Narragansett, Munsee
Delaware), Proto-Algonquian reconstructions with sound-law application, and
David Costa's grammatical charts and comparative research.

Stephanie Fielding spent the 2017--2018 academic year as a presidential
visiting fellow at Yale University, teaching a course comparing the Mohegan
and Delaware languages. She has also worked with the Unkechaug people of Long
Island to help reclaim their nearly identical language (the primary difference
being the Mohegan /y/ vs. Unkechaug /r/ correspondence).

\subsection{Vowel Inventory}

\subsubsection{Vowel System}

The Mohegan vowel system in the Fielding practical orthography consists of
six plain vowels, four long vowels, and three diphthongs:

\begin{table}[h]
\centering
\begin{tabular}{llll}
\toprule
Grapheme & IPA & Length & Notes \\
\midrule
a & /a/ & short & Low central \\
â & /aː/ & long & \\
e & /iː/ ~ /e/ & short & Front, historically /iː/ lowered; can be /e/ \\
ê & /eː/ & long & Rare; mid front \\
i & /i/ & short & High front \\
o & (not used as plain vowel) & --- & Only occurs in diphthongs \\
ô & /oː/ & long & Very common; mid back rounded \\
u & /u/ & short & High back rounded \\
 & /ə/ & short & Schwa, not written consistently \\
\bottomrule
\end{tabular}
\caption{Mohegan vowel system}
\label{tab:1364-vowels}
\end{table}

\subsubsection{Diphthongs}

\begin{table}[h]
\centering
\begin{tabular}{lll}
\toprule
Grapheme & IPA & Notes \\
\midrule
ay & /aj/ & \\
aw & /aw/ & \\
uy & /uj/ & \\
\bottomrule
\end{tabular}
\caption{Mohegan diphthongs}
\label{tab:1364-diphthongs}
\end{table}

Granberry (2003:14) also reconstructs /oi/ /ɔj/ as marginal diphthongs (as in
English ``boy'') from certain sources.

\subsection{Consonant Inventory}

\begin{table}[h]
\centering
\begin{tabular}{llllll}
\toprule
Grapheme & IPA & Voicing & Place & Manner & Notes \\
\midrule
p & /p/ & Voiceless & Bilabial & Stop & Unaspirated lenis \\
t & /t/ & Voiceless & Dental & Stop & Unaspirated lenis \\
k & /k/ & Voiceless & Velar & Stop & Unaspirated lenis \\
c & /tʃ/ & Voiceless & Palatal & Affricate & Always /tʃ/, never /k/ \\
s & /s/ & Voiceless & Alveolar & Fricative & Fortis/lenis not marked \\
sh & /ʃ/ & Voiceless & Palatal & Fricative & \\
h & /h/ & Voiceless & Glottal & Fricative & Explicitly written \\
m & /m/ & Voiced & Bilabial & Nasal & \\
n & /n/ & Voiced & Alveolar & Nasal & \\
w & /w/ & Voiced & Labiovelar & Approximant & \\
y & /j/ & Voiced & Palatal & Approximant & \\
kw & /kʷ/ & Voiceless & Labiovelar & Stop & Labialized \\
q & /kʷ/ & Voiceless & Labiovelar & Stop & Word-final /kʷ/ \\
' & /ʔ/ & --- & Glottal & Stop & Okina \\
\bottomrule
\end{tabular}
\caption{Mohegan consonant inventory}
\label{tab:1364-consonants}
\end{table}

\subsubsection{Allophonic Variation (from Granberry 2003)}

Granberry's phonological reconstruction (using the ``reconstituted phonetic
units'' methodology of Mary Haas 1954) identifies allophonic patterns:

\begin{itemize}
  \item Stops /p, t, k/ are \textbf{lenis-fortis pairs} overlapping in
    distribution --- the practical orthography uses voiceless symbols but
    actual realization varies
  \item /kw/ word-finally often reduced to [k] before consonants
  \item Vowels show positional length variation: pre-stress short, stressed
    medium-long, post-stress reduced
  \item The phoneme /e/ (plain \textit{e}) ranges from [i] to [e] depending
    on environment, likely reflecting the historical merger of PA *i and *e
    in initial syllables (Oxford 2015)
\end{itemize}

\subsection{Historical Phonology: Proto-Algonquian Reflexes}

\subsubsection{PA Vowel System}

Proto-Algonquian is reconstructed with an eight-vowel system: four qualities
with a length contrast (Bloomfield 1946).

\begin{table}[h]
\centering
\begin{tabular}{lll}
\toprule
PA & Mohegan Reflex & Example \\
\midrule
*a & a & \\
*aː & â & \\
*e & i, a, ə (variable) & \\
*eː & ê & \\
*i & i & \\
*iː & e & \\
*o & u & \\
*oː & ô & \\
\bottomrule
\end{tabular}
\caption{PA to Mohegan vowel reflexes}
\label{tab:1364-pa-vowels}
\end{table}

\subsubsection{The Y-Dialect Classification}

Mohegan-Pequot is classified as a \textbf{Y-dialect} of Southern New England
Algonquian (SNEA). In the traditional classification based on the reflex of
PA *r (sometimes reconstructed as *l):

\begin{table}[h]
\centering
\begin{tabular}{lll}
\toprule
Dialect Type & Reflex of PA *r & Example Language \\
\midrule
N-dialect & /n/ & Massachusett (Wôpanâak), Narragansett, Nipmuc \\
Y-dialect & /y/ & \textbf{Mohegan-Pequot}, Montauk \\
R-dialect & /r/ & Western Connecticut (Loup area) \\
L-dialect & /l/ & Munsee Delaware, Unami \\
\bottomrule
\end{tabular}
\caption{SNEA dialect classification}
\label{tab:1364-dialects}
\end{table}

This means that where Wôpanâak has \textit{n}, Mohegan typically has
\textit{y}. Stephanie Fielding noted this correspondence operating even in
modern Connecticut English pronunciation (e.g., ``Sullivan'' →
``Suh-yuh-van''). Granberry (2003:13) explicitly documents the core
correspondences: Pequot \textit{p, t, k, ch, s} → Mohegan \textit{b, d, g,
j, z} form-initially, and final Pequot \textit{p, t, k, ch, s} → Mohegan
form-final \textit{b, d, g, j, z}.

\subsubsection{Consonant Reflexes}

\begin{table}[h]
\centering
\begin{tabular}{lll}
\toprule
PA & Mohegan & Notes \\
\midrule
*p & p & Unaspirated \\
*t & t & Unaspirated; palatalized to c /tʃ/ before *i, *iː, *y \\
*č & c /tʃ/ & \\
*k & k & Unaspirated \\
*m & m & \\
*n & n & \\
*θ & n (N-dialects), y (Y-dialects) & Key PA distinction \\
*r (controversial) & y & The Y-dialect defining reflex \\
*s & s & \\
*š & sh /ʃ/ & \\
*h & h & Explicit in modern orthography \\
*w & w & \\
\bottomrule
\end{tabular}
\caption{PA to Mohegan consonant reflexes}
\label{tab:1364-pa-consonants}
\end{table}

\subsubsection{Cluster Reflexes}

\begin{table}[h]
\centering
\begin{tabular}{lll}
\toprule
PA Cluster & Mohegan Reflex & Notes \\
\midrule
*ht & ht & Retained in some forms, simplified in others \\
*hk & k & \\
*hp & p & \\
*šk & shk /ʃk/ & \\
*θp & p & With devoicing \\
*mp & p & Nasal cluster simplification \\
*nt & t & \\
*nk & k & \\
*nč & c /tʃ/ & \\
\bottomrule
\end{tabular}
\caption{PA cluster reflexes in Mohegan}
\label{tab:1364-clusters}
\end{table}

\subsection{Comparison with Prince \& Speck (1904)}

\subsubsection{The Prince-Speck Glossary}

J. Dyneley Prince (Columbia University) and Frank G. Speck published
``Glossary of the Mohegan-Pequot Language'' in \emph{American Anthropologist}
(1904, vol. 6, no. 1, pp. 18--44). This glossary contains \textbf{446 words
and forms} collected by Speck from Fidelia Fielding between 1902 and 1903,
submitted to Prince in Fielding's own spelling.

Prince and Speck deliberately preserved Fielding's ``imperfect''
English-orthography system for ``sentimental reasons'' (1904:19), while
providing phonetic transcriptions in parentheses.

\subsubsection{Orthographic Differences}

\begin{table}[h]
\centering
\begin{tabular}{lll}
\toprule
Feature & Fielding (2006) & Prince-Speck (1904) \\
\midrule
Affricate /tʃ/ & \textit{c} always & \textit{ch} or contextual \\
Velar /k/ & \textit{k} & \textit{g} initially, \textit{k, c} elsewhere \\
Voicing & Unaspirated voiceless & Inconsistent (g/k, b/p variation) \\
Vowel length & Circumflex (â, ê, ô) & Colon or context \\
Glottal stop & \texttt{'} (okina) & Inverted comma \\
Schwa & Not written & \texttt{'} (apostrophe) \\
/ʃ/ & \textit{sh} & \textit{sh} \\
Stops & p, t, k (unaspirated) & b/g/p/t/k/d (inconsistent) \\
\bottomrule
\end{tabular}
\caption{Orthographic differences between Fielding and Prince-Speck}
\label{tab:1364-orth-diff}
\end{table}

\subsubsection{The Brothertown Material}

A subset of Prince-Speck's glossary consists of ``Brothertown words'' ---
items collected from an old Mohegan man who had lived at Brothertown, near
Green Bay, Wisconsin, where relocated New England tribes (Tunxis, Wampanoag,
Mohegans, Montauks) had settled. Prince and Speck note these words are
``merely corruptions of Ojibwe forms'' (1904:19) --- relevant since our
corpus bears the name ``Brothertown.''

\subsubsection{Convergence and Divergence}

Most Fielding (2006) entries with FF attestation can be mapped to cognates in
Prince-Speck (1904). The primary differences are systematic orthographic
conventions: Prince-Speck's \textit{g} initial corresponds to Fielding's
\textit{k} (voicing neutralization); Prince-Speck's \textit{ch} corresponds to
Fielding's \textit{c} (just a spelling difference); Prince-Speck's vowel
notation is inconsistent (reflecting English spelling conventions) while
Fielding's is phonemic; and Prince-Speck uses \texttt{'} for schwa while
Fielding omits schwa or encodes it as \textit{u}, \textit{a}, or \textit{i}.

\subsection{Current Orthographic Conventions}

\subsubsection{The Modern Practical Orthography}

The orthography used by Fielding (2006) and adopted by the Mohegan Language
Project is a \textbf{Latin-based phonemic alphabet} with the following
conventions:

\textbf{Alphabet (12 consonant symbols, 6 vowel symbols):}
\texttt{a â c e ê h i k m n ô p q s sh t u w y '}

Key design principles:
\begin{enumerate}
  \item \textbf{1:1 Phoneme-to-grapheme correspondence} --- each symbol maps
    to exactly one phoneme
  \item \textbf{Circumflex for length} --- â, ê, ô are long; a, e, i, u are
    short
  \item \textbf{\textit{c} for /tʃ/} --- avoids the English \textit{ch}
    digraph; never represents /k/
  \item \textbf{\textit{q} for /kʷ/} --- distinguishes from \textit{kw}
    cluster; occurs word-finally
  \item \textbf{\texttt{'} (okina) for glottal stop} --- follows Hawaiian and
    other indigenous orthography conventions
  \item \textbf{\textit{sh} for /ʃ/} --- standard English digraph
  \item \textbf{No silent letters} --- every written symbol is pronounced
  \item \textbf{Bound morphemes marked with \texttt{-}} --- prefixes:
    \textit{nu-}, \textit{ku-}; suffixes: \textit{-côq}
  \item \textbf{Part of speech tags} --- VAI, VII, VTA, VTI, NI, NA following
    Algonquianist conventions
  \item \textbf{Alternate spellings in parentheses} --- showing earlier
    English-based forms
\end{enumerate}

\subsubsection{Comparison with Granberry's Orthography}

Julian Granberry (2003) uses a somewhat different practical orthography,
differing primarily in representing lenis-fortis distinctions:

\begin{table}[h]
\centering
\begin{tabular}{lll}
\toprule
Fielding & Granberry & IPA \\
\midrule
p & b, p & /p/ \\
t & d, t & /t/ \\
k & g, k & /k/ \\
c & j, ch & /tʃ/ \\
s & z, s & /s/ \\
sh & zh, sh & /ʃ/ \\
\bottomrule
\end{tabular}
\caption{Fielding vs. Granberry orthography}
\label{tab:1364-granberry}
\end{table}

Granberry's orthography explicitly notes the vowel /ə/ (schwa) as \textit{u}
--- e.g., \textit{wutche} vs. Fielding's \textit{wicki}. He also uses
\textit{x} for /h/ in some environments. However, the Fielding orthography
has become the standard through the Mohegan Language Project's adoption.

\subsection{Problems Encountered}

\subsubsection{Small Corpus}

Only 145 records with explicit Mohegan-to-phoneme field mappings. This is the
smallest corpus among the three SNEA sources, limiting statistical approaches.

\subsubsection{Alternate Spellings}

The primary conversion challenge is English-based alternate spellings →
Fielding orthography, not Fielding → IPA. The alternate spellings provide
direct training data for converting Prince-Speck (1904) English-orthography
forms.

\subsubsection{Circularity Risk}

Some Fielding dictionary entries are reconstructed FROM Prince-Speck (1904)
--- using Fielding to verify Prince-Speck is circular for those entries.
Careful provenance tracking is required to distinguish entries with
independent FF attestation from those reconstructed via comparative methods.

\subsubsection{The Pequot Extinction and the Treaty of Hartford (1638)}

The Pequot War (1636--1638) culminated in the \textbf{Mystic Massacre} (May
26, 1637), where English forces and allies set fire to a Pequot fortified
village, blocking the exits and shooting anyone attempting to escape.
Approximately 700 Pequots were killed or captured; survivors were sold into
slavery in Bermuda and the West Indies, while others were dispersed as
captives to the victorious tribes.

The \textbf{Treaty of Hartford} (September 21, 1638) explicitly sought to
eradicate Pequot cultural identity:
\begin{itemize}
  \item \textbf{Name prohibition:} Survivors could no longer call themselves
    Pequots --- they were to be referred to as Mohegans or Narragansetts
    ``for ever''
  \item \textbf{Land prohibition:} Pequots were forbidden from returning to
    their native country
  \item \textbf{Dispersion:} Surviving Pequots were divided: ~80 to Mohegan
    Sachem Uncas, ~80 to Narragansett Sachem Miantonomo, ~20 to Niantic
    Sachem Ninigret
  \item \textbf{English control of Pequot country:} The Pequot lands went to
    the Connecticut River towns
\end{itemize}

The treaty effectively declared the Pequot nation extinct as a sovereign
polity. A 1743 copy of the treaty survives and is housed at the Connecticut
State Archives. Daragh Grant (2015, \emph{William and Mary Quarterly})
reconstructed the original 13-provision text from multiple surviving copies,
including a newly rediscovered 1665 copy in the British Library's Lansdowne
Manuscripts.

The Treaty of Hartford directly caused the language's near-extinction. By
outlawing Pequot speech and dispersing the population, the colonial government
ensured the language would cease to be transmitted to children. Mohegan
survived longer (until 1908) precisely because the Mohegan tribe was not
subject to the same dispersal --- their territory west of the Thames River
remained intact.

\subsection{Research Approach}

\subsubsection{Recommended Methodology}

Direct 1:1 phoneme mapping. Fielding's orthography is already phonemic ---
conversion is a direct lookup. The primary challenge is mapping alternate
English-based spellings to the standard Fielding form.

For the 145 records in our corpus, the practical orthography's transparency
means:
\begin{itemize}
  \item Phonemic conversion from the orthography is essentially 1:1 (98\%+
    regular)
  \item The primary challenge is not decoding Fielding orthography → phonemes,
    but mapping alternate spellings → Fielding orthography
  \item The alternate spellings provide direct training data for converting
    Prince-Speck (1904) English-orthography forms
  \item Vowel quality for /e/ (plain \textit{e}) is the single remaining
    ambiguity --- the phoneme ranges from /i/ to /e/ and its exact quality
    may need contextual disambiguation
\end{itemize}

\subsection{Conversion Workflow}
\texttt{[Pending implementation --- see Issue \#1364]}

\subsection{Linguist Feedback}
\texttt{[Template --- content pending review]}

\subsection{Related Work: Colonial Recorder Perception Problem}

The core challenge of this research card --- English-speaking recorders
misperceiving and mis-transcribing Algonquian phonemes --- is a known
phenomenon documented across Eastern North America. Two researchers' work is
directly relevant:

\textbf{Dr. Keith Cunningham} (Linguist, Nanticoke Indian Tribe). His
doctoral dissertation \emph{A Phonological Analysis of Nanticoke With
Practical Applications for Language Revitalization} (Georgetown University)
analyzes three 18th-century Nanticoke word lists (Heckewelder 1785) using the
same methodology: historical phonology from colonial recorders, with the same
challenges of multilingual informants and orthographic complexity. His 2024
Algonquian Conference presentation \emph{A Revised Phonological Analysis of
the Heckewelder Vocabulary of Nanticoke} demonstrates that the colonial
recorder perception problem extends beyond SNEA into the broader Eastern
Algonquian family. Cunningham's work on Nanticoke revitalization (co-authoring
a Nanticoke textbook) parallels the practical application of the Fielding
dictionary for Mohegan language reclamation.

\textbf{Dr. Craig Kopris} (Independent scholar). His paper \emph{Les sons
wyandots perçus par des oreilles étrangères} (Wyandot sounds perceived by
foreign ears, \emph{Recherches amérindiennes au Québec}, 2014) is the exact
framing of the Fielding/Prince-Speck problem: how European recorders
systematically misperceived and mis-transcribed indigenous phonemes. His
\emph{Wyandot Phonology: Recovering the Sound System of an Extinct Language}
applies the same recovery methodology to an unrelated language family
(Iroquoian), demonstrating that the foreign-ear perception problem is
universal across colonial documentation of North American languages. His
Dictionary of Wyandot project (NEH award FN-255576-17) shows the long-term
application of these methods to language revitalization.

The Mohegan-Pequot corpus is unique among the SNEA sources in having a
last-speaker corpus (Fidelia Fielding's diaries) alongside a modern practical
orthography (Fielding 2013). This positions it as the most phonologically
reliable source --- but the Prince-Speck (1904) glossary, with its
German-influenced transcription, reintroduces the foreign-ear perception
problem that Kopris documents for Wyandot.

\subsection{Mohegan Revitalization}

\subsubsection{The Mohegan Language Project}

Following the death of Fidelia Fielding in 1908, Mohegan had no fluent
speakers for nearly a century. In April 1998, the Council of Mohegan Tribal
Elders resolved to use Jits Bodunaxa's speech as the basis for a restored
Modern Mohegan (Granberry 2003:6). A Mohegan Language Committee was formed,
and Julian Granberry was appointed Tribal Linguist from 1998 to 2001.

In the 21st century, the \textbf{Mohegan Language Project} was established,
creating lessons, a dictionary, and online learning materials. The project
developed a practical orthography using the Latin alphabet, based on the four
phonetic diaries left by Fielding.

\subsubsection{Current Revitalization Efforts (as of 2024)}

The Mohegan Tribe operates an active \textbf{Language Restoration Project}
that ramped up significantly in 2017. Key initiatives include:
\begin{itemize}
  \item \textbf{Master-Apprentice Program:} Immersion-style learning with
    passive speaking → circumlocution → active speaking
  \item \textbf{Immersion Nests:} Yearlong cohort programs scaffolded to
    seasonal activities
  \item \textbf{Workbook Development:} First Mohegan language workbooks and
    grammar curricula
  \item \textbf{Two fluent teachers} (the most since Fidelia Fielding's time)
    leading classes
  \item \textbf{Online classes} open to all Mohegan tribal members and Mohegan
    household members
\end{itemize}

\subsubsection{Collaboration with Jessie Little Doe Baird}

MacArthur Fellow \textbf{Jessie Little Doe Baird} (Mashpee Wampanoag),
renowned for her Wôpanâak language reclamation, has served as a guide for the
Mohegan project. Using Wôpanâak (an N-dialect) as a sister language, the
Mohegan teachers cross-reference sound changes (N→Y correspondence) to build
vocabulary and grammar.

\subsubsection{The Samson Occom Papers}

In 2022, Dartmouth College repatriated the papers of \textbf{Samson Occom}
(1723--1792), a Mohegan minister and scholar who wrote extensively in his
native language. These 18th-century documents provide another primary source
layer for the Language Restoration Project, offering examples of the language
written by a fluent speaker during the period when Mohegan was still actively
spoken.

\subsubsection{Terminology: Reclamation vs. Revitalization}

Mohegan language teachers \textbf{Allyson Corbin} and \textbf{Autumn Cholewa},
who created much of the current curriculum, emphasize the term
\textbf{reclamation} over revitalization or revival:

\begin{quote}
``Reclaim is much more empowering for a community. The connotation is more
hands-on. More like, `This is up to you to do.' To reclaim it is to bring it
home. It's a prideful thing instead of something people felt pressure not to
do.'' (CT Insider, Nov 2024)
\end{quote}

\subsection{Bibliography}

\begin{enumerate}

\item Baird, Jessie Little Doe. (2010). MacArthur Fellowship for Wôpanâak
Language Reclamation.

\item Bloomfield, Leonard. (1925). ``On the Sound-System of Central
Algonquian.'' \emph{Language} 1:130--156.

\item Bloomfield, Leonard. (1946). ``Algonquian.'' In \emph{Linguistic
Structures of Native America}, ed. Harry Hoijer, 85--129. Viking Fund
Publications in Anthropology 6.

\item Costa, David. (2006). Lexical and grammatical research for the Modern
Mohegan Dictionary (unpublished ms).

\item Fielding, Stephanie. (2006). \emph{A Modern Mohegan Dictionary.}
Mohegan Tribe.

\item Fielding, Stephanie. (2013). \emph{Modern Mohegan Dictionary} (revised
edition). Mohegan Tribe.

\item Goddard, Ives. (1978). ``Eastern Algonquian Languages.'' In
\emph{Handbook of North American Indians}, vol. 15: Northeast, ed. Bruce G.
Trigger, 70--77. Smithsonian Institution.

\item Goddard, Ives. (1979). ``Comparative Algonquian.'' In \emph{The
Languages of Native America}, ed. Lyle Campbell and Marianne Mithun, 70--132.
University of Texas Press.

\item Granberry, Julian. (2003). \emph{Modern Mohegan: The Dialect of Jits
Bodunaxa.} Languages of the World/Materials 430. LINCOM Europa.

\item Granberry, Julian. (2003). \emph{A Lexicon of Modern Mohegan: The
Dialect of Jits Bodunaxa.} Languages of the World/Dictionaries 28. LINCOM
Europa.

\item Grant, Daragh. (2015). ``The Treaty of Hartford (1638): Reconsidering
Jurisdiction in Southern New England.'' \emph{William and Mary Quarterly}
72(3):461--498.

\item Haas, Mary. (1954). ``The Proto-Algonquian Word for `Sun.''' In
\emph{Language in Culture and Society}, ed. Dell Hymes. Harper \& Row.

\item Oxford, Will. (2015). ``Patterns of Contrast in Phonological Change:
Evidence from Algonquian Vowel Systems.'' \emph{Language} 91(2):308--339.

\item Oxford, Will. (2023). ``Consonant Inventories from Proto-Algonquian to
the Daughter Languages.'' University of Manitoba Algonquian Linguistics
Resources.

\item Pentland, David H. (1979). \emph{Algonquian Historical Phonology.} PhD
dissertation, University of Toronto.

\item Prince, J. Dyneley \& Speck, Frank G. (1903). ``The Mohegan-Pequot
Language.'' \emph{American Anthropologist} 5:193--212.

\item Prince, J. Dyneley \& Speck, Frank G. (1904). ``Glossary of the
Mohegan-Pequot Language.'' \emph{American Anthropologist} 6(1):18--44.

\item Siebert, Frank T. (1975). ``Resurrecting Virginia Algonquian from the
Dead.'' In \emph{Studies in Southeastern Indian Languages}, ed. James M.
Crawford, 285--453. University of Georgia Press.

\item Speck, Frank G. (1928). ``Native Tribes and Dialects of Connecticut.''
\emph{Annual Report of the Bureau of American Ethnology} 43:199--287.

\item Valentine, J. Randolph. (2001). \emph{Nishnaabemwin Reference
Grammar.} University of Toronto Press.

\end{enumerate}

% Wampanoag Multi-Source — Source Workflow
% Issue #1365
% Status: Research complete, implementation pending
% !TEX program = xelatex
% !TEX encoding = UTF-8 Unicode

\section{Wampanoag: Multi-Source (Wood, Winslow, Shepard, Eliot)}
\label{sec:1365}

\subsection{Source Description}

This section covers all historical Wampanoag/Massachusett source vocabularies
beyond the Trumbull compilation (covered in \S\ref{sec:1363}). Four primary
sources span the critical transition from raw English-hearer perception through
systematic missionary orthography:

\begin{description}
  \item[William Wood (1634)] \emph{New Englands Prospect}, \textasciitilde262
    words. Pre-Eliot English ad-hoc orthography. Wood was an English colonist
    living in the Massachusetts Bay Colony with no linguistic training. His
    transcriptions reflect the earliest English-hearer perception of Wampanoag
    sounds, before any regularization system existed.
  \item[Edward Winslow (1624)] \emph{Good Newes from New England}, 22 scattered
    words and phrases. The earliest published Wampanoag vocabulary. Winslow was
    a Mayflower passenger and Plymouth Colony governor who described the
    language as ``very copious, large, and difficult'' and was unusually honest
    about his linguistic limitations.
  \item[Thomas Shepard (1647)] \emph{The Day-Breaking, if not the sun-rising
    of the Gospell with the Indians in New-England}, 11 native words plus two
    one-sentence prayers. Predates Eliot's grammar by 19 years and his Indian
    Primer by 22 years. The prayers represent Eliot's earliest work in the
    language, before he had fully analyzed its grammar.
  \item[John Eliot (1666)] \emph{The Indian Grammar Begun}, the first published
    grammatical analysis of any Algonquian language. Eliot's orthography is a
    deliberately constructed system, not a casual English-based spelling. He
    distinguishes aspirated vs.\ unaspirated stops using English voiced/voiceless
    pairs (p\textasciitilde b, t\textasciitilde d, k\textasciitilde g).
\end{description}

\subsection{Orthography Systems}

\subsubsection{William Wood (1634)}

Wood's vocabulary is the largest pre-Eliot Massachusett word list. Its
orthography follows English sound-spelling conventions of the 1630s:

\begin{itemize}
  \item \textbf{Initial /w/ retained:} \emph{Wigwam}, \emph{Web}, \emph{Wompey}
  \item \textbf{$\langle$ea$\rangle$ for /iː/:} \emph{Kean} (cf.\ Eliot's \emph{keen})
  \item \textbf{$\langle$s$\rangle$ for /ʃ/ inconsistently:} \emph{Sawup} where
    later sources might write \emph{sh-}
  \item \textbf{$\langle$k$\rangle$ preferred over $\langle$c$\rangle$:}
    Initial /k/ written \emph{K-} even before front vowels (\emph{Kean},
    \emph{Ksitta})
  \item \textbf{No diacritics:} Wood uses plain English letters only
  \item \textbf{Vowel length not marked:} Short vs.\ long vowel distinction
    invisible
  \item \textbf{Pidginization noted:} Goddard (1977) observes a ``good deal of
    pidginization'' in Wood's data (Salvucci 2002)
\end{itemize}

\subsubsection{Edward Winslow (1624)}

Winslow's orthography is purely English-hearer perception with zero
regularization. He was recording an entirely unfamiliar language with no
reference system. The vocabulary overwhelmingly represents the Pokanoket
dialect of Wampanoag (Massasoit's people), with some possible influence from
the Patuxet band through Tisquantum.

Key characteristics:
\begin{itemize}
  \item \textbf{Pre-Eliot:} No systematic orthography; words are ad-hoc English
    approximations
  \item \textbf{$\langle$h$\rangle$ written reliably:} Winslow writes /h/ more
    consistently than any other early Wampanoag source except Edwards (1787
    Mahican). This is one of the most valuable aspects of his small corpus.
  \item \textbf{$\langle$l$\rangle$ $\rightarrow$ $\langle$n$\rangle$:}
    \emph{Winsnow} for ``Winslow'' --- no /l/ in Wampanoag
  \item \textbf{Pidginization:} Goddard (1977) identifies clear evidence of
    pidginization, including use of transitive inanimate \emph{namen} where
    transitive animate \emph{nunaw} would be grammatically correct
\end{itemize}

\subsubsection{Thomas Shepard (1647)}

The Wampanoag words in this pamphlet show unstandardized English orthography
with the following characteristics:

\begin{itemize}
  \item \textbf{$\langle$ch$\rangle$ for /č/:} \emph{Chechainuppan},
    \emph{Chepian}, \emph{Cheehesom}
  \item \textbf{$\langle$sh$\rangle$ and $\langle$s$\rangle$ for /š/:}
    \emph{Evangenesh}, \emph{Iehovah} (given English biblical name spelling)
  \item \textbf{$\langle$w$\rangle$ for /w/:} \emph{Wigwam}, \emph{Wowein},
    \emph{wicke}
  \item \textbf{Final -an for animate nouns:} \emph{Chepian},
    \emph{Chechainuppan}
  \item \textbf{$\langle$oo$\rangle$ for /uː/:} \emph{Noonatomen}
  \item \textbf{$\langle$p$\rangle$ for /b/:} \emph{Chepian} (cf.\ Eliot's
    \emph{chepiân})
  \item \textbf{Preaspiration:} \emph{cke} in \emph{wicke} may indicate
    preaspirated stop
\end{itemize}

The pamphlet contains the earliest recorded Wampanoag Christian prayers:

\begin{quote}
  \emph{Amanaomen Iehovah tahassen metagh} --- ``Take away Lord my Stony heart''
\end{quote}

\begin{quote}
  \emph{Cheehesom Iehovah kekowhogkew} --- ``Wash Lord my soule''
\end{quote}

\subsubsection{John Eliot (1666)}

Eliot's grammar presents a deliberately constructed orthography. His key
innovation is distinguishing aspirated vs.\ unaspirated stops using English
voiced/voiceless pairs (p\textasciitilde b, t\textasciitilde d,
k\textasciitilde g). This is the most important orthographic feature for
phonological reconstruction.

Eliot describes 5 primary vowels with length distinctions:
\begin{itemize}
  \item \textbf{a} (short) vs.\ \textbf{â} (long)
  \item \textbf{e} /iː/ (always long)
  \item \textbf{i} (short) vs.\ \textbf{ee}/\textbf{ie} (long --- inconsistently
    marked)
  \item \textbf{o} (short) vs.\ \textbf{ô} (long)
  \item \textbf{u} (short schwa) vs.\ \textbf{oo} (long /uː/)
\end{itemize}

\subsection{Vowel and Consonant Tables}

\begin{table}[htbp]
\centering
\caption{Comparative Orthography Table}
\label{tab:1365-compare}
\begin{tabular}{llllll}
\hline
English & Wood (1634) & Shepard (1647) & Eliot (1666+) & Cotton (1708) & Trumbull (1903) \\
\hline
man & --- & --- & wosketomp & wosketomp & wosketomp \\
woman & --- & --- & mittamwossis & mittamwo˘ssis & mittámwossis \\
house & Wigwam & Wigwam & wetu & wetu & wétu \\
devil & --- & Chepian & chepiân & chepiân & chepeân \\
king/sachem & --- & Sachem & sôchim & sôchim & sôchim \\
I & Kean & --- & keen & keen & kēn \\
very angry & --- & musquantum & --- & --- & --- \\
good & Winnet & --- & wunne & wunne & wunne \\
white & Wompey & --- & wompi & wompi & wômpei \\
sleep (classifier) & Sawup & --- & --- & --- & --- \\
\hline
\end{tabular}
\end{table}

\begin{table}[htbp]
\centering
\caption{Vowel Representations by Source}
\label{tab:1365-vowels}
\begin{tabular}{llllll}
\hline
Phoneme & Wood & Shepard & Eliot & Cotton & Trumbull \\
\hline
/iː/ & ea, ee & e, ee & e, ee & e, ee & ē \\
/a/ & a & a & a & a & a \\
/aː/ & a, aw & a & â & â & â \\
/uː/ & oo & o, oo & o, ô, ꝏ & o, ô & ô, u \\
/ə/ & --- & a, e & u & u, u˘ & u \\
/o/ & o & o & o & o & o \\
\hline
\end{tabular}
\end{table}

\begin{table}[htbp]
\centering
\caption{Consonant Representations by Source}
\label{tab:1365-consonants}
\begin{tabular}{llllll}
\hline
Phoneme & Wood & Shepard & Eliot & Cotton & Trumbull \\
\hline
/p/ & p & p & p\textasciitilde b & p\textasciitilde b & p \\
/pʰ/ & --- & --- & p & p & p (with dot) \\
/t/ & t & t & t\textasciitilde d & t\textasciitilde d & t \\
/tʰ/ & --- & --- & t & t & t (with dot) \\
/k/ & k & k & k\textasciitilde g & k\textasciitilde g & k \\
/kʰ/ & --- & --- & k & k & k (with dot) \\
/č/ & ch & ch & ch & ch & tc \\
/š/ & sh & sh & sh & sh & sh \\
/w/ & w & w & w & w & w \\
/j/ & --- & --- & y & y & y \\
\hline
\end{tabular}
\end{table}

\subsection{Source-Specific Orthographic Patterns}

Each source has characteristic ``errors'' or systematic deviations from a
standardized phonemic representation:

\begin{description}
  \item[Wood (1634)] No vowel length; English digraphs predominate; treats
    unaspirated stops as voiced; pidginized forms; word-initial \emph{a-}
    prefix frequently attached to roots.
  \item[Shepard (1647)] No consistent system; words are ad-hoc English
    approximations; preaspiration inconsistently shown; $\langle$ck$\rangle$
    for /k/.
  \item[Eliot (1666+)] Most consistent system; treats unaspirated stops as
    voiced English letters (b, d, g) but means unaspirated; vowel length
    inconsistently marked on some vowels; macron on long vowels (â, ô);
    distinctive $\langle$ꝏ$\rangle$ (oo ligature) for /uː/.
  \item[Cotton (1708)] Extends Eliot's system; adds breve for ultra-short
    vowels; shows late shift in some vowels; more consistent vowel marking.
  \item[Trumbull (1903)] Standardized historical reconstruction; marks length
    and aspiration consistently; normalizes variation.
\end{description}

\subsection{Problems Encountered}

\begin{itemize}
  \item \textbf{Source provenance correction:} The ``Anonymous 1647'' vocabulary
    in ALR Vol.\ 27 is actually Thomas Shepard's \emph{The Day-Breaking}.
    Furthermore, the sample words in the original ALR listing (\emph{kean},
    \emph{kenomptom}, \emph{wean}, \emph{nukus}, \emph{kodtomp}, \emph{ompmom},
    \emph{sawop}, \emph{acheta co}) do NOT match the words actually found in
    the 1647 text. These words originate from William Wood's 1634 vocabulary.
    The correct Shepard word list contains 11 items including the two prayers.
  \item \textbf{Pre-Eliot vs.\ post-Eliot:} Sources span the transition from
    raw English-hearer perception (Wood, Winslow) through systematic missionary
    orthography (Eliot). The same letter sequence may have different phonetic
    values in different sources.
  \item \textbf{Eliot's stop aspiration system:} Eliot is the only colonial
    source who systematically distinguishes aspirated from unaspirated stops
    using orthographic pairs (b/p, d/t, k/g). His $\langle$b, d, g$\rangle$
    represent unaspirated lenis stops (not voiced), while $\langle$p, t,
    k$\rangle$ represent aspirated fortis stops.
  \item \textbf{Pidginization:} Goddard (1977) identifies pidginization in both
    Wood and Winslow. If the English were learning a simplified contact variety,
    word order may not reflect full Wampanoag syntax, verb forms may be bare
    stems rather than inflected forms, and vocabulary may be mixed from multiple
    Algonquian languages.
  \item \textbf{Intermediary filter:} Winslow's data passed through Tisquantum
    (Squanto) and Hobbamock as interpreters. Tisquantum's English was learned
    in London, not from the Pilgrims, and his Wampanoag was Patuxet, a
    sub-dialect closely related to but not identical with Pokanoket.
  \item \textbf{Dialect uncertainty:} Pokanoket (Winslow/Wood) vs.\ Natick
    (Eliot/Trumbull) differences are unresolved. The distinction may represent
    real dialect differences, register differences (pidgin vs.\ full language),
    chronological differences (1620s vs.\ 1660s), or orthographic convergence.
  \item \textbf{Vowel instability:} Wood's triple spelling
    \emph{abbomocho}/\emph{abomocho}/\emph{abamacho} for ``devil'' is the
    strongest evidence in the entire SNEA corpus for colonial vowel instability.
  \item \textbf{Preaspiration /h/ unreported:} Like all pre-Eliot sources, Wood
    does not write preaspirated /h/ before stops. Winslow is the exception ---
    he writes /h/ more consistently than any other early source.
\end{itemize}

\subsection{Research Approach}

The multi-source parsing strategy follows a four-stage pipeline:

\begin{enumerate}
  \item \textbf{Source detection:} Identify which source's orthographic
    fingerprint is present. Each source has a characteristic orthographic
    texture that should be identifiable from spelling patterns.
  \item \textbf{Source normalizer:} Convert to Eliot-consistent form. Apply
    source-specific decoding rules --- the same letter sequence may have
    different phonetic values in different sources.
  \item \textbf{Eliot parser:} Apply Grammar rules to produce a phonemic
    representation. Eliot's orthography is the most systematic and serves as
    the intermediate representation.
  \item \textbf{Unified output:} Produce standard modern transcription.
\end{enumerate}

\subsubsection{Source-Adaptive Parsing Strategy}

\begin{description}
  \item[Wood] English sound-spelling $\rightarrow$ phoneme via lookup table.
    Key mappings: ea$\rightarrow$/iː/, oo$\rightarrow$/uː/, k$\rightarrow$/k/,
    s$\rightarrow$/s\textasciitilde š/.
  \item[Shepard] Ad-hoc English $\rightarrow$ phoneme via context. Same as Wood
    but with higher ambiguity.
  \item[Eliot] Eliot orthography $\rightarrow$ phoneme via Grammar rules.
    p/b$\rightarrow$/p/ vs.\ /pʰ/, t/d$\rightarrow$/t/ vs.\ /tʰ/,
    k/g$\rightarrow$/k/ vs.\ /kʰ/, â$\rightarrow$/aː/, ô$\rightarrow$/oː\textasciitilde uː/.
  \item[Cotton] Same as Eliot + breve $\rightarrow$ extra-short vowel. Vowel
    quantity more reliable.
  \item[Trumbull] Modern normalized $\rightarrow$ phoneme directly. The cleanest
    mapping.
\end{description}

\subsubsection{Resolving Ambiguity}

Where a spelling could represent multiple phonemes (e.g., Wood's $\langle$ea$\rangle$
= /iː/ or /ɛː/), the parser uses source context and Algonquian comparative
evidence. The Proto-Algonquian reconstructions in Aubin (1975) provide the
external check needed to disambiguate competing analyses.

\subsection{Conversion Workflow}
\texttt{[Pending implementation --- see Issue \#1365]}

The conversion pipeline will be implemented as a Python module in
\texttt{src/commons/parsing/}. The module will:

\begin{itemize}
  \item Detect source fingerprint from orthographic patterns
  \item Apply source-specific normalizer functions
  \item Route through the unified Eliot parser
  \item Output standard modern transcription
  \item Handle the pre-Eliot vs.\ post-Eliot transition
  \item Resolve ambiguity using PA comparative evidence
\end{itemize}

\subsection{Linguist Feedback Template}

\texttt{[Template --- content pending review]}

The following questions are open for linguist review:

\begin{itemize}
  \item Are the source-specific vowel and consonant mappings correct?
  \item Does the proposed Eliot-consistent intermediate representation capture
    the phonemic distinctions correctly?
  \item Are there additional sources that should be included?
  \item Is the pidginization analysis (Goddard 1977) accepted, or are there
    alternative interpretations?
  \item How should the Pokanoket/Natick dialect distinction be handled in the
    unified output?
\end{itemize}

\subsection{Related Work: Colonial Recorder Perception Problem}

The core challenge of this research card --- English-speaking recorders
misperceiving and mis-transcribing Algonquian phonemes --- is a known
phenomenon documented across Eastern North America. Two researchers' work is
directly relevant:

\textbf{Dr.\ Keith Cunningham} (Linguist, Nanticoke Indian Tribe). His doctoral
dissertation \emph{A Phonological Analysis of Nanticoke With Practical
Applications for Language Revitalization} (Georgetown University) analyzes three
18th-century Nanticoke word lists (Heckewelder 1785) using the same methodology:
historical phonology from colonial recorders, with the same challenges of
multilingual informants and orthographic complexity. His 2024 Algonquian
Conference presentation \emph{A Revised Phonological Analysis of the Heckewelder
Vocabulary of Nanticoke} directly parallels the multi-source Wampanoag analysis
--- both involve multiple colonial recorders (Wood, Winslow, Shepard, Eliot for
Wampanoag; Heckewelder, Du Ponceau for Nanticoke) with different orthographic
conventions that must be cross-calibrated. Cunningham's work demonstrates that
the colonial recorder perception problem extends beyond SNEA into the broader
Eastern Algonquian family.

\textbf{Dr.\ Craig Kopris} (Independent scholar). His paper \emph{Les sons
wyandots perçus par des oreilles étrangères} (Wyandot sounds perceived by
foreign ears, \emph{Recherches amérindiennes au Québec}, 2014) is the exact
framing of the multi-source Wampanoag problem: how European recorders with
different linguistic backgrounds (17th c.\ English, 18th c.\ English, 19th c.\
German-influenced American English) systematically misperceived and
mis-transcribed indigenous phonemes. His \emph{Wyandot Phonology: Recovering the
Sound System of an Extinct Language} applies the same recovery methodology to an
unrelated language family (Iroquoian), demonstrating that the foreign-ear
perception problem is universal across colonial documentation of North American
languages.

The multi-source Wampanoag corpus (Wood 1634, Winslow 1624, Shepard 1647,
Eliot 1666) spans the full range of colonial recorder competence --- from raw
English-hearer perception (Wood, Winslow) through systematic missionary
orthography (Eliot). This range makes it the best case study for the recorder
calibration problem that both Cunningham and Kopris address in their respective
language families.

\subsection{Key References}

\begin{itemize}
  \item Aubin, George F. 1975. \emph{A Proto-Algonquian Dictionary}. Ottawa:
    National Museums of Canada.
  \item Aubin, George F. 1978. ``Toward the Linguistic History of an Algonquian
    Dialect: Observations on the Vocabulary.'' In William Cowan (ed.),
    \emph{Ninth Algonquian Conference}, pp.\ 1--9.
  \item Clark, Michael P. (ed.). 2003. \emph{The Eliot Tracts}. Westport, CT:
    Praeger.
  \item Cogley, Richard W. 1999. \emph{John Eliot's Mission to the Indians
    before King Philip's War}. Cambridge, MA: Harvard University Press.
  \item Cotton, Josiah. 1829. \emph{Vocabulary of the Massachusetts (or Natick)
    Indian Language}. Cambridge: Charles Folsom.
  \item Eliot, John. 1666. \emph{The Indian Grammar Begun}. Cambridge:
    Marmaduke Johnson.
  \item Goddard, Ives. 1977. ``The Pidginization of Massachusetts: Evidence
    from Wood's Vocabulary.'' \emph{International Journal of American
    Linguistics} 43: 222--226.
  \item Goddard, Ives \& Bragdon, Kathleen J. 1988. \emph{Native Writings in
    Massachusett}. Philadelphia: American Philosophical Society.
  \item Pilling, James C. 1891. \emph{Bibliography of the Algonquian Languages}.
    Washington, D.C.: Government Printing Office.
  \item Salvucci, Claudio (ed.). 2002. \emph{ALR Volume 27: Wood's Vocabulary of
    Massachusett}. Bristol, PA: Evolution Publishing.
  \item Shepard, Thomas (attrib.). 1647. \emph{The Day-breaking, if not the
    sun-rising of the Gospell with the Indians in New-England}. London: Rich.
    Cotes.
  \item Trumbull, J. Hammond. 1903. \emph{Natick Dictionary}. Washington, D.C.:
    Bureau of American Ethnology.
  \item Williams, Roger. 1643. \emph{A Key into the Language of America}.
    London: Gregory Dexter.
\end{itemize}

% Wampanoag (Winslow 1624) — Source Workflow
% Issue #1366
% Status: Research complete, implementation pending
% !TEX program = xelatex
% !TEX encoding = UTF-8 Unicode

\section{Wampanoag: Winslow 1624}
\label{sec:1366}

\subsection{Source Description}

\textbf{Recorder:} Edward Winslow (1595--1655), Mayflower passenger, Plymouth
Colony governor for three terms (1633--1634, 1636--1637, 1644--1645), and the
colony's primary negotiator and diplomat with the Native peoples of New England.

\textbf{Source:} \emph{Good Newes from New England} (1624), printed in London by
I.D.\ for William Bladen and John Bellamie. The earliest published source of
Wampanoag vocabulary in the SNEA corpus.

\textbf{Language:} Wampanoag, Pokanoket dialect [wam].

\textbf{Corpus:} 24 records (the smallest in the SNEA collection).

\textbf{Orthography:} English-hearer perception with zero regularization.
Winslow was unusually honest about his linguistic limitations, describing the
Massachusett language as ``very copious, large, and difficult. As yet we cannot
attain to any great measure thereof.'' This acknowledgment distinguishes
Winslow from many other colonial recorders who pretended to greater competence.
His transcriptions are therefore likely more reliable as raw phonetic
impressions --- he was not normalizing or regularizing what he heard.

\subsection{Recorder Biography}

Edward Winslow was among the 41 Mayflower passengers who signed the Mayflower
Compact in November 1620. Key events relevant to his linguistic data:

\begin{itemize}
  \item \textbf{March 1621:} First contact with Samoset (Abenaki sachem speaking
    English) and Massasoit.
  \item \textbf{March--July 1621:} Tisquantum (Squanto) serves as interpreter,
    teaching the Pilgrims Wampanoag vocabulary.
  \item \textbf{March 1623:} Winslow cures Massasoit of a life-threatening
    illness --- the encounter that produced his most extensive Wampanoag
    dialogue.
  \item \textbf{August 1623:} Hobbamock warns Plymouth of Massasoit's alleged
    conspiracy with the Massachusett; Standish's Wessagussett raid follows.
  \item \textbf{1624:} \emph{Good News from New England} published in London.
\end{itemize}

After returning to England in 1646, Winslow served as Oliver Cromwell's chief
agent in foreign missions. He died of yellow fever in 1655 en route to a
diplomatic mission in Spanish Hispaniola.

\subsection{Linguistic Intermediaries}

Winslow did not learn Wampanoag directly --- he worked through interpreters.
The identity and linguistic background of these intermediaries is crucial for
understanding potential biases in his transcriptions.

\subsubsection{Tisquantum (Squanto, c.~1585--1622)}

Tisquantum was a Patuxet Wampanoag with an extraordinary linguistic history:

\begin{itemize}
  \item \textbf{1614:} Kidnaped by Captain Thomas Hunt, along with 20+ other
    Patuxet men.
  \item \textbf{1614--1618:} Sold as a slave in Málaga, Spain; escaped with the
    help of Spanish friars.
  \item \textbf{1618--1619:} Lived in London with John Slany, a merchant of the
    Newfoundland Company; learned English.
  \item \textbf{1619:} Returned to Patuxet (modern Plymouth) --- only to find
    his entire village depopulated by the 1616--1619 epidemic.
  \item \textbf{1621:} Approached Plymouth Colony in March; served as
    interpreter and guide.
  \item \textbf{1622:} Died of ``Indian fever'' at Monomoit (Chatham, Cape Cod)
    while accompanying Winslow on a diplomatic mission.
\end{itemize}

Linguistic implications of Tisquantum's mediation:
\begin{itemize}
  \item His English was learned in London, not from the Pilgrims --- it was
    16th-century English, not colonial dialect.
  \item His Wampanoag was Patuxet, a sub-dialect of Wampanoag closely related
    to Pokanoket but not identical.
  \item His time in Spain and England (8+ years) may have affected his native
    language production (attrition).
\end{itemize}

\subsubsection{Hobbamock (Hobomok, d.~1642)}

Hobbamock was a Pokanoket \emph{pniese} --- a military champion and elite
warrior --- who served as Massasoit's trusted advisor and was assigned to live
among the English:

\begin{itemize}
  \item \textbf{March 1621:} Assigned by Massasoit to reside at Plymouth Colony
    as a liaison.
  \item \textbf{March 1623:} Revealed Massasoit's alleged conspiracy with the
    Massachusett, leading to Standish's raid.
  \item \textbf{Accompanied Winslow} on the mission to cure Massasoit --- his
    lament ``Neen womasu sagimus'' (``My loving sachem'') upon hearing of
    Massasoit's apparent death is one of Winslow's most famous transcriptions.
  \item \textbf{Died of European disease} in 1642.
\end{itemize}

Linguistic significance: Hobbamock likely spoke Pokanoket dialect natively.
He was not a language teacher but may have provided a more natural sample of
Wampanoag speech than Tisquantum. His English was likely limited --- Winslow
probably communicated with him through a mixture of Wampanoag and signs.

\subsection{Complete Vocabulary}

According to Salvucci (InteractiveALR Vol.\ 27) and confirmed by primary source
reading, Winslow's text contains 22 distinct words and phrases in Wampanoag:

\begin{table}[htbp]
\centering
\caption{Winslow's Wampanoag Lexicon}
\label{tab:1366-lexicon}
\begin{tabular}{lll}
\hline
Word/Phrase & English & Notes \\
\hline
Keen Winsnow? & ``Art thou Winslow?'' & Massasoit's greeting; l$\rightarrow$n \\
Ahhe & ``Yes'' & Winslow's reply \\
Matta neen wonckanet namen Winsnow & ``O Winslow, I shall never see thee again!'' & Deathbed lament \\
Winsnow & ``Winslow'' & /l/ $\rightarrow$ /n/; no /l/ in Wampanoag \\
Neen womasu sagimus & ``My loving sachem'' & Hobbamock's lament \\
Sachimo comaco & ``The sachem's house'' & Sachem + house at Pokanoket \\
Witeo & ``An ordinary house'' & Contrasts with comaco \\
Squasachim & ``The sachem's wife'' & Squa (woman) + sachim (chief) \\
Maskiet & ``Physic/Medicine'' & Healing context \\
Kiehtan & (The supreme God) & Explains Wampanoag theology \\
Chise & ``An old man'' & Winslow's etymology for Kiehtan \\
Kiehchise & ``A very old man'' & Comparative: chise $\rightarrow$ kiehchise \\
Hinnaim namen, hinnaim michen, matta cuts & ``By and by see, eat, but not speak'' & Complex sentence \\
Quatchet & ``Walk abroad'' & Imperative? \\
Sachim & ``King/Chief'' & General title \\
Pniese (also pinse, pinese) & ``Chief champion / man of valor'' & Military elite class \\
Powah & ``Shaman / healer'' & Calls upon the devil \\
Squaw & ``Woman'' & Occurring in compounds \\
Massasoit & ``Great Chief'' & Proper name / title \\
Canonickus & Massasoit's older brother & Proper name with -us? \\
Wamsutta & Massasoit's son & Proper name \\
Namasquemit & ``Powah'' or dreaming healer & One reference only \\
\hline
\end{tabular}
\end{table}

\subsection{Syntactic Complexity}

Despite the small corpus, Winslow records several sentences showing Wampanoag
syntax:

\begin{description}
  \item[Keen Winsnow?] Copular question: \emph{Keen} = 2sg pronoun ``thur'',
    predicate nominal.
  \item[Matta neen wonckanet namen Winsnow] Negative declarative: \emph{Matta}
    (neg) + \emph{neen} (1sg) + \emph{wonckanet} (again/future) + \emph{namen}
    (see) + vocative.
  \item[Hinnaim namen, hinnaim michen, matta cuts] Serial verb construction:
    three coordinate clauses with \emph{hinnaim} (by-and-by) + \emph{matta}
    (neg).
  \item[Neen womasu sagimus] Possessive + noun: \emph{neen} (I/my) +
    \emph{womasu} (loving?) + \emph{sagimus} (sachem + 1sg possessive).
\end{description}

\subsection{Orthography System}

Winslow wrote in Early Modern English orthography (pre-1625), before English
spelling standardization. Understanding what Winslow's spellings meant for
English is essential before mapping to Wampanoag phonemes.

\subsubsection{Vowel Grapheme-to-Phoneme Mapping}

\begin{table}[htbp]
\centering
\caption{Winslow Vowel Mapping}
\label{tab:1366-vowels}
\begin{tabular}{llll}
\hline
Grapheme & IPA & Example & Confidence \\
\hline
a (short) & /a/ $\sim$ /ə/ & matta & High \\
a (long) & /aː/ & ahhe (first a) & Moderate \\
e (short) & /ɛ/ & chise (final e) & Moderate \\
e (final) & /ə/ & ahhe, witeo, maskiet & High \\
i & /ɪ/ $\sim$ /i/ & kiehchise, chise & Moderate \\
o & /ɔ/ $\sim$ /ə/ & hobbamock, commaco & Moderate \\
u & /ʌ/ $\sim$ /ə/ & cuts & Low \\
ie & /iː/ & kiehtan, kiehchise & High \\
oo & /uː/ & askooke & High \\
ah & /ah/ & ahhe & Moderate \\
ee & /iː/ & keen (thou) & High \\
\hline
\end{tabular}
\end{table}

\subsubsection{Consonant Grapheme-to-Phoneme Mapping}

\begin{table}[htbp]
\centering
\caption{Winslow Consonant Mapping}
\label{tab:1366-consonants}
\begin{tabular}{llll}
\hline
Grapheme & IPA & Example & Confidence \\
\hline
h & /h/ & ahhe, hobbamock, kiehtan & High \\
k & /k/ & kiehtan, askooke, keen & High \\
c (before a/o) & /k/ & commaco, cuts, comaco & High \\
ch & /tʃ/ & chise, kiehchise & High \\
sh & /ʃ/ & Only in proper names & Low \\
m & /m/ & commaco, matta, namen & High \\
n & /n/ & neen, namen, Winsnow & High \\
s & /s/ & sachim, squaw, sagimus & High \\
t & /t/ & matta, kiehtan, quatchet & High \\
w & /w/ & winsnow, witeo & High \\
p & /p/ & pniese (initial pn-!) & Moderate \\
b & /b/ & hobbamock (medial) & Moderate \\
\hline
\end{tabular}
\end{table}

\subsubsection{Preaspiration /h/ in Winslow}

Winslow writes /h/ more consistently than any other early Wampanoag source
except Edwards (1787 Mahican). This is one of the most valuable aspects of his
small corpus.

\begin{table}[htbp]
\centering
\caption{Cross-Source Comparison for /h/}
\label{tab:1366-h}
\begin{tabular}{llll}
\hline
Source & hobbamock initial h & kiehtan medial ht & Overall /h/ Reliability \\
\hline
Winslow (1624) & YES & YES & Highest among early sources \\
Wood (1634) & NO (abamocho) & N/A & Unreliable for /h/ \\
Williams (1643) & Mixed (ábammocho) & N/A & Unreliable \\
Eliot (1663) & Different root & Likely kehtan & Partial \\
Edwards (1787) & N/A (Mahican) & YES & Most reliable overall \\
\hline
\end{tabular}
\end{table}

Interpretation: Winslow's relatively good /h/ transcription may reflect:
\begin{itemize}
  \item His intermediaries (Tisquantum, Hobbamock) had clearer /h/ articulation
  \item The Pokanoket dialect had more salient /h/ than Natick
  \item Winslow was a more careful listener --- his ``ignorance'' of the
    language made him more phonetically accurate (less normalization bias)
\end{itemize}

\subsection{Pidginization Analysis (Goddard 1977)}

Goddard's landmark analysis of early contact Algonquian data identified
Winslow's vocabulary as showing clear evidence of pidginization --- the English
and Wampanoag were communicating via a simplified contact variety, not the full
grammatical language.

\subsubsection{Key Evidence}

In the sentence \emph{Matta neen wonckanet namen Winsnow}:

\begin{enumerate}
  \item \textbf{\emph{namen} (transitive inanimate ``see it'')} is used instead
    of the grammatically correct transitive animate form (expected \emph{nunaw}
    ``I see him''). The verb agrees with the object for gender --- \emph{namen}
    treats ``Winslow'' as if he were an inanimate object.
  \item \textbf{Pro-drop is lost:} Full independent pronouns (\emph{neen} ``I'',
    \emph{keen} ``thou'') are used where verbal inflection alone would carry
    person in the full language. This is a universal pidgin feature ---
    morphology simplifies, syntax becomes analytic.
  \item \textbf{No verbal inflection:} The expected inflected verb
    (\emph{n-nam-en} = 1sg-see-TA) is replaced with what appears to be the
    citation form or a nominalized form.
\end{enumerate}

\subsubsection{The Massachusett Pidgin Hypothesis}

Goddard hypothesized that a Massachusett Pidgin (similar in function to Delaware
Pidgin or Mobilian Jargon) existed before English colonization, used for
inter-tribal trade. If so, what Winslow recorded may not have been ``pure''
Wampanoag at all, but rather a trade pidgin. This would explain:

\begin{itemize}
  \item The consistency of certain simplifications (pidgin features shared
    across sources)
  \item The spread of certain forms (e.g., \emph{matta} for negation appears in
    all colonial sources)
  \item The divergence between ``missionary'' Massachusett (Eliot 1663) and
    ``contact'' Massachusett (Winslow, Wood)
\end{itemize}

\subsection{Cross-Source Cognate Table}

\begin{table}[htbp]
\centering
\caption{Confirmed Cognates Across Sources}
\label{tab:1366-cognates}
\begin{tabular}{lllll}
\hline
Gloss & Winslow (1624) & Wood (1634) & Trumbull/Eliot & Williams (1643) \\
\hline
`no' & matta & matta & matta & matta \\
`woman' & squa(achim) & squaw & squa & squa \\
`devil' & hobbamock & abamocho & achm∞wonk & ábammocho \\
`yes' & ahhe & --- & --- & --- \\
`house' & witeo & wetu/weetu & wetu & --- \\
`snake' & askooke & askook & askkuhnk & askug \\
`king' & sachim & sachim & sachim & --- \\
`old man' & chise & --- & --- & --- \\
`speak' & cuts & --- & cuts & ketes \\
`walk' & quatchet & --- & --- & --- \\
`medicine' & maskiet & --- & maskiht & --- \\
`shaman' & powah & powaw & --- & powwáw \\
\hline
\end{tabular}
\end{table}

\subsection{Dialect Affiliation: Pokanoket Wampanoag}

Winslow's data was collected primarily at Plymouth (Patuxet) and Pokanoket
(modern Bristol, RI --- Massasoit's capital). The Wampanoag Confederacy in 1620
consisted of several bands: Pokanoket (Massasoit's people, primary), Patuxet
(Plymouth, Tisquantum's people, depopulated by 1620), Nauset (Cape Cod),
Monomoit (Chatham), Manomet (Sandwich/Bourne), and Mashpee (Cape Cod).

The key dialectal division within Massachusett/Wampanoag:

\begin{table}[htbp]
\centering
\caption{Pokanoket vs.\ Natick Features}
\label{tab:1366-dialect}
\begin{tabular}{llll}
\hline
Feature & Pokanoket (Winslow/Wood) & Natick (Eliot/Trumbull) & Notes \\
\hline
`house' & witeo (Winslow) & wetu (Eliot) & Real lexical difference? \\
`devil' & hobbamock (Winslow) & achm∞wonk (Eliot) & Different roots entirely \\
`man of valor' & pniese (Winslow) & Not attested & Nativized English? \\
Initial /h/ & Written (Winslow) & Often omitted (Eliot) & Possibly orthographic \\
/l/ substitution & n only (Winslow) & n only (Eliot) & Consistent across dialects \\
\hline
\end{tabular}
\end{table}

Per Salvucci's assessment: ``Winslow's data has not been systematically
analyzed for dialect affiliation.'' The Pokanoket sub-corpus is so small (22
items) that confident dialect assignment is impossible without external
evidence.

\subsection{Problems Encountered}

\begin{itemize}
  \item \textbf{Tiny corpus:} 24 records is insufficient for any statistical
    approach. Every entry must be treated as a unique case. Winslow is a
    corroborative source only.
  \item \textbf{Pidginization:} Goddard (1977) identifies ``a good deal of
    pidginization'' in Winslow's data. This affects whether the data can be
    used for phonological reconstruction at all. If the English were learning a
    pidgin variety, word order may not reflect full Wampanoag syntax, verb forms
    may be bare stems, and vocabulary may be mixed from multiple Algonquian
    languages.
  \item \textbf{Intermediary filter:} Tisquantum (Squanto) and Hobbamock served
    as interpreters --- their own linguistic backgrounds may have affected the
    transcriptions. Tisquantum's English was learned in London, not from the
    Pilgrims, and his Wampanoag was Patuxet, a sub-dialect closely related to
    but not identical with Pokanoket.
  \item \textbf{Dialect uncertainty:} Pokanoket vs.\ Natick differences are
    unresolved. The distinction may represent real dialect differences, register
    differences (pidgin vs.\ full language), chronological differences (1620s
    vs.\ 1660s), or orthographic convergence.
  \item \textbf{Proper name contamination:} Some entries include English names
    (Winsnow = Winslow), which follow English orthography and must be excluded
    from phonemic analysis.
  \item \textbf{Multi-word entries:} Several entries are entire sentences
    requiring syntactic segmentation before lexical analysis.
  \item \textbf{Uncertain dialect:} Pokanoket dialect vs.\ Patuxet (via
    Tisquantum) vs.\ generic pidgin. Not resolvable from Winslow's data alone.
\end{itemize}

\subsection{Research Approach}

Five-step corroboration pipeline using Trumbull and Wood as primary sources.
Winslow data is treated as corroborative, not primary. Each entry is verified
against expected Eliot/Trumbull forms.

\subsubsection{Step 1: Corroboration-Only}

Treat Winslow as corroborative evidence for rules derived from larger sources.
For each Winslow entry, find the cognate in larger Wampanoag sources (Trumbull
4,491 records, Wood 339 records). Use the cross-source consensus for the
phonemic mapping. Flag Winslow-only words for manual review.

\subsubsection{Step 2: Key Cognate Mappings}

\begin{table}[htbp]
\centering
\caption{Key Cognate Mappings}
\label{tab:1366-key-cognates}
\begin{tabular}{llll}
\hline
Winslow & Wood & Trumbull/Eliot (expected) & PA Source \\
\hline
hobbamock (devil) & abamocho & Different root & See §11 --- two different roots? \\
matta (no) & matta & matta & PA *mati- \\
kiehtan (God) & --- & kehtan & PA *ki·ht-? \\
sachim (chief) & sachim & sachim & PA *sa·kima·wa \\
squa (woman) & squaw & squa & PA *e·θkwe·wa \\
askooke (snake) & askook & askkuhnk & PA *aθkw-? \\
\hline
\end{tabular}
\end{table}

\subsubsection{Step 3: h-Transcription as Secondary Evidence}

Winslow's relatively consistent /h/ transcriptions provide secondary evidence
for where preaspiration occurs:
\begin{itemize}
  \item When Winslow writes \emph{h} and other sources also confirm it
    $\rightarrow$ well-attested
  \item When Winslow writes \emph{h} and other sources omit it $\rightarrow$
    possible but needs external confirmation
  \item When Winslow omits \emph{h} and Edwards (Mahican) writes it
    $\rightarrow$ probable preaspiration unrecorded
\end{itemize}

\subsubsection{Step 4: Segment k- Prefix}

The 2nd person prefix \emph{k-} is seen in \emph{keen} (thou, 2sg pronoun) and
should be segmented before conversion. Note that \emph{k-} with vowel-initial
stems would produce different surface forms than \emph{k-} with consonant-initial
stems.

\subsubsection{Step 5: Verify Against Trumbull}

For each converted Winslow entry, verify against the expected form in the
Eliot/Trumbull orthography. Flag mismatches for manual review --- especially:
\begin{itemize}
  \item Where Winslow writes /h/ and Trumbull does not (possible preaspiration
    evidence)
  \item Where Winslow and Trumbull use different lexical roots (possible dialect
    divergence)
  \item Where Winslow's form contradicts known PA correspondences
\end{itemize}

\subsection{Conversion Workflow}
\texttt{[Pending implementation --- see Issue \#1366]}

The conversion pipeline will be implemented as a Python module in
\texttt{src/commons/parsing/}. The module will:

\begin{itemize}
  \item Treat Winslow as corroborative only --- no regex-based conversion
  \item For each Winslow entry, find the cognate in larger Wampanoag sources
  \item Use cross-source consensus for phonemic mapping
  \item Flag Winslow-only words for manual review
  \item Segment \emph{k-} prefix before conversion
  \item Verify each converted entry against expected Eliot/Trumbull forms
\end{itemize}

\subsection{Linguist Feedback Template}

\texttt{[Template --- content pending review]}

The following questions are open for linguist review:

\begin{enumerate}
  \item \textbf{Pidgin or genuine dialect?} Are the simplifications in Winslow's
    data (e.g., \emph{namen} for ``see'' used as transitive animate) evidence
    of a pre-existing Massachusett Pidgin (Goddard 1977) or simply a learner's
    imperfect recording? This affects whether Winslow's data can be used for
    phonological reconstruction at all.
  \item \textbf{\emph{witeo} vs.\ \emph{wetu}:} Is the \emph{witeo} (Winslow) /
    \emph{wetu} (Eliot) split evidence of Pokanoket vs.\ Natick dialect
    difference, or are both from PA *wi·kiwa·mi with different English
    ear-perceptions? Williams records Narragansett \emph{wetu} (agreeing with
    Eliot), suggesting \emph{witeo} may be uniquely Pokanoket.
  \item \textbf{\emph{hobbamock} volcanism:} Is \emph{hobbamock} a genuine
    Wampanoag lexical item or a portmanteau of Wampanoag + English
    ``hobgoblin''? The initial \emph{h} and the \emph{-ock} suffix appear in
    both the proper name and the common noun ``devil'', suggesting folk
    etymology contamination is possible.
  \item \textbf{Phonological reliability of pidgin data:} If pidgin phonology
    approximates the substrate, then Winslow's transcriptions are phonologically
    useful. If the pidgin introduced phonological simplifications (cluster
    reduction, final devoicing, vowel merger), then even the phonology is
    unreliable.
  \item \textbf{\emph{pniese} etymology:} Is this the same root as PA
    *-epye·wa (``to be powerful, to be a man'')? The initial cluster /pn-/ is
    unusual for Massachusett and may reflect a specific morphotactic condition
    or epenthesis.
\end{enumerate}

\subsection{Related Work: Colonial Recorder Perception Problem}

The core challenge of this research card --- English-speaking recorders
misperceiving and mis-transcribing Algonquian phonemes --- is a known
phenomenon documented across Eastern North America. Two researchers' work is
directly relevant:

\textbf{Dr.\ Keith Cunningham} (Linguist, Nanticoke Indian Tribe). His doctoral
dissertation \emph{A Phonological Analysis of Nanticoke With Practical
Applications for Language Revitalization} (Georgetown University) analyzes three
18th-century Nanticoke word lists (Heckewelder 1785) using the same methodology:
historical phonology from colonial recorders, with the same challenges of
multilingual informants and orthographic complexity. His 2024 Algonquian
Conference presentation \emph{A Revised Phonological Analysis of the Heckewelder
Vocabulary of Nanticoke} directly parallels the Winslow analysis --- both
involve tiny corpora (Winslow 24 records, Heckewelder \textasciitilde200)
recorded by English speakers with no linguistic training, both require PA
comparative calibration, and both must account for intermediary effects
(Tisquantum/Hobbamock for Winslow, multilingual informants for Heckewelder).
Cunningham's work demonstrates that the colonial recorder perception problem
extends beyond SNEA into the broader Eastern Algonquian family.

\textbf{Dr.\ Craig Kopris} (Independent scholar). His paper \emph{Les sons
wyandots perçus par des oreilles étrangères} (Wyandot sounds perceived by
foreign ears, \emph{Recherches amérindiennes au Québec}, 2014) is the exact
framing of the Winslow problem: how European recorders systematically
misperceived and mis-transcribed indigenous phonemes. His \emph{Wyandot
Phonology: Recovering the Sound System of an Extinct Language} applies the same
recovery methodology to an unrelated language family (Iroquoian), demonstrating
that the foreign-ear perception problem is universal across colonial
documentation of North American languages.

The Winslow corpus (24 records, the smallest in the SNEA collection) represents
the most extreme case of the tiny-corpus problem. The methodological question
--- whether such a small dataset can yield reliable phonological data --- is
directly addressed by both Cunningham's work on small Nanticoke word lists and
Kopris's work on recovering phonology from limited documentary sources.

\subsection{Key References}

\begin{itemize}
  \item Aubin, George F. 1975. \emph{A Proto-Algonquian Dictionary}. Ottawa:
    National Museums of Canada.
  \item Goddard, Ives. 1977. ``Massachusett Pidgin.'' \emph{Language} 53(1):
    147--184.
  \item Goddard, Ives \& Bragdon, Kathleen. 1988. \emph{Native Writings in
    Massachusett}. Transactions of the American Philosophical Society, Vol.\ 78.
  \item O'Brien, Frank Waobu. \emph{Guide to Historical Spellings}. --- Master
    tables for spelling-to-phoneme mapping across SNEA languages.
  \item Salvucci, Claudio. \emph{InteractiveALR Vol.\ 27: Massachusetts}.
    Evolution Publishing.
  \item Trumbull, J. Hammond. 1903. \emph{Natick Dictionary}. Washington, D.C.:
    Bureau of American Ethnology.
  \item Williams, Roger. 1643. \emph{A Key into the Language of America}.
    London: Gregory Dexter.
  \item Winslow, Edward. 1624. \emph{Good Newes from New England}. London: I.D.
    for William Bladen and John Bellamie.
  \item Winslow, Edward. 2014. \emph{Good News from New England}. Ed.\ Kelly
    Wisecup. University of Massachusetts Press. ISBN 9781625340832.
\end{itemize}

% Wampanoag (Wood 1634) — Source Workflow
% Issue #1367
% Status: Research complete, implementation pending
% !TEX program = xelatex
% !TEX encoding = UTF-8 Unicode

\section{Wampanoag: Wood 1634}
\label{sec:1367}

\subsection{Source Description}

\textbf{Recorder:} William Wood, English colonist living in the Massachusetts
Bay Colony (1629--1633). He returned to London to publish his observations.
Wood had no linguistic training; his orthography is purely English-based ad-hoc
spelling.

\textbf{Source:} \emph{New Englands Prospect: A true, lively, and experimentall
description of that part of America, commonly called Nevv England} (1634),
printed by Tho.\ Cotes for John Bellamie, London. Small 4to, contains a ``small
Nomenclator'' of \textasciitilde262 Massachusett words organized by semantic
category (body parts, animals, trade goods, place names, month names,
day-name classifiers).

\textbf{Language:} Wampanoag [wam].

\textbf{Corpus:} 339 records (the second-largest pre-Eliot Wampanoag source).

\textbf{Orthography:} Pre-Eliot English ad-hoc, no linguistic training. The
Eliot Massachusett writing system was still decades away. This makes Wood's
transcriptions reflect raw English-hearer perception of Wampanoag sounds,
without the regularization that Eliot later introduced.

\subsection{Recorder Biography}

William Wood was an English colonist who lived in the Massachusetts Bay Colony
from 1629 to 1633. Little is known of his life beyond his publication. He
returned to London and published \emph{New Englands Prospect} in 1634, which
went through multiple editions:

\begin{table}[htbp]
\centering
\caption{Publication History of \emph{New Englands Prospect}}
\label{tab:1367-editions}
\begin{tabular}{llll}
\hline
Edition & Year & Printer & Place \\
\hline
1st & 1634 & Tho.\ Cotes for John Bellamie & London \\
2nd & 1635 & --- & London \\
3rd & 1639 & John Dawson for John Bellamie & London \\
4th (Boston) & 1764 & Thomas and John Fleet & Boston \\
Prince Society & 1865 & --- & Boston \\
\hline
\end{tabular}
\end{table}

The 1865 Prince Society edition is the most accessible for modern researchers.
Wood's Wampanoag vocabulary appears in all editions (pp.\ 123--128 in the 1764
edition, pp.\ 94--100 in the 1865 reprint), suggesting the word list was stable
across editions with no significant revision.

\subsection{Orthography System}

Wood's orthography follows English sound-spelling conventions of the 1630s.
It is a pre-Eliot system --- the Eliot Massachusett writing system was still
decades away. Key characteristics:

\begin{itemize}
  \item \textbf{Initial /w/ retained:} \emph{Wigwam}, \emph{Web}, \emph{Wompey}
  \item \textbf{$\langle$ea$\rangle$ for /iː/:} \emph{Kean} (cf.\ Eliot's
    \emph{keen})
  \item \textbf{$\langle$s$\rangle$ for /ʃ/ inconsistently:} \emph{Sawup} where
    later sources might write \emph{sh-}
  \item \textbf{$\langle$k$\rangle$ preferred over $\langle$c$\rangle$:}
    Initial /k/ written \emph{K-} even before front vowels (\emph{Kean},
    \emph{Ksitta})
  \item \textbf{No diacritics:} Wood uses plain English letters only
  \item \textbf{Vowel length not marked:} Short vs.\ long vowel distinction
    invisible
  \item \textbf{Pidginization noted:} Goddard (1977) observes a ``good deal of
    pidginization'' in Wood's data (Salvucci 2002)
\end{itemize}

\subsection{Vowel Tables}

\begin{table}[htbp]
\centering
\caption{Wood Vowel Rules}
\label{tab:1367-vowels}
\begin{tabular}{llll}
\hline
Wood Grapheme & Likely Phoneme & Example & Notes \\
\hline
a & /a/ $\sim$ /ə/ & abbona (five), matta (no), abamocho & Short a dominant \\
e & /ɛ/ $\sim$ /i/ & aberginian (an Indian), apponees (twelve) & \\
i & /ɪ/ $\sim$ /i/ & aberginian & \\
o & /ɔ/ $\sim$ /o/ $\sim$ /ə/ & abbona, abamocho, abbomocho & Broad range \\
u & /ə/ $\sim$ /ʌ/ & sucqunnocquock (sleeps), nugge (sieve) & \\
a followed by a & /aː/ (implied) & abbona, abamacho & Double letters often for emphasis \\
oo & /uː/ $\sim$ /ɔ/ & askooke (in Winslow), week (in Wood) & \\
au & /ɔː/ $\sim$ /aw/ & pausawniscosu (half fathom), aug (I) & \\
ee & /iː/ & apponees (twelve), nees (two) & \\
ai & /eː/ $\sim$ /aj/ & pausawniscosu, appause (morning) & \\
\hline
\end{tabular}
\end{table}

\subsection{Consonant Tables}

\begin{table}[htbp]
\centering
\caption{Wood Consonant Rules}
\label{tab:1367-consonants}
\begin{tabular}{llll}
\hline
Wood Grapheme & Likely Phoneme & Example & Notes \\
\hline
b & /p/ $\sim$ /b/ & abbona (five), abamocho (devil) & Wood uses b where later sources use p \\
p & /p/ $\sim$ /b/ & apponees, appause & \\
m & /m/ & matta (no), abamocho & \\
n & /n/ & abbona, an nu ocke (a bed) & \\
t & /t/ $\sim$ /d/ & matta, abbona quit & \\
s & /s/ & apponees, pausawniscosu & \\
ch & /tʃ/ & & \\
sh & /ʃ/ & & \\
k & /k/ $\sim$ /g/ & nugge (sieve), asquoquin & \\
g & /g/ $\sim$ /k/ & nugge & \\
qu & /kʷ/ & sucqunnocquock, asquoquin & Labialized k \\
w & /w/ & naw aug (I will tell), yoaw & \\
ck & /k/ & asquoquin, an nu ocke & \\
\hline
\end{tabular}
\end{table}

\subsection{Key Orthographic Observations}

\subsubsection{b vs.\ p Alternation}

Wood writes \emph{abbona} (five) and \emph{abamocho} (devil) with \emph{b}
where the Eliot/Trumbull orthography would use \emph{p}. This shows Wood
perceived the unaspirated Wampanoag /p/ closer to English /b/ --- direct
evidence of the voicing neutralization.

\subsubsection{Geminate Consonants}

Wood frequently doubles consonants: \emph{bb}, \emph{pp}, \emph{nn}, \emph{ck},
\emph{gg}. This probably indicates the preceding vowel is short/checked rather
than true gemination.

\subsubsection{qu Digraph}

The \emph{qu} digraph consistently represents /kʷ/ --- seen in
\emph{sucqunnocquock} (sleeps), \emph{asquoquin}.

\subsubsection{Word-Initial a- Prefix}

Word-initial \emph{a-} is extremely common: \emph{abbona}, \emph{abamocho},
\emph{aberginian}, \emph{abbomocho}, \emph{abonetta}. This may be a
demonstrative/indefinite prefix that Wood transcribed as part of the root.
Roughly 25--30\% of nouns and adjectives in Wood's 339-entry word list begin
with \emph{a-}.

\subsubsection{Preaspiration /h/}

Wood does \textbf{not} write preaspirated /h/ before stops. Like Williams,
English-trained ears miss this feature. Example: \emph{matta} (no) --- modern
Wampanoag /maht-a/ or similar; \emph{sucqunnocquock} --- no \emph{h} written
despite probable /h/ in the medial cluster.

\subsection{Number System}

Wood has unusually good coverage of Wampanoag numbers (10+ distinct number
entries). This is valuable because numerals are relatively easy to verify
cross-linguistically:

\begin{table}[htbp]
\centering
\caption{Wood's Number System}
\label{tab:1367-numbers}
\begin{tabular}{lll}
\hline
Wood & Meaning & Expected Wampanoag \\
\hline
abbona & five & (p)àpàna? \\
nawhut & maybe one & nequt \\
nees & two & nees \\
quín & three & qush, shwa \\
yoâw & four & yaw, yau \\
apponees & twelve & --- \\
apponabonna & fifteen & --- \\
apponaquinta & sixteen & --- \\
apponasquoquin & nineteen & --- \\
apponenotta & seventeen & --- \\
\hline
\end{tabular}
\end{table}

\subsection{Comparison: Wood vs.\ Eliot Orthography}

\begin{table}[htbp]
\centering
\caption{Wood vs.\ Eliot Grapheme Mapping}
\label{tab:1367-vs-eliot}
\begin{tabular}{llll}
\hline
Wood & Eliot & Phoneme & Reliability \\
\hline
a & a, â & /a/, /ã/ & Low (Wood conflates length/nasal) \\
e & e & /ɛ/ & Moderate \\
i & i, u & /ɪ/, /ə/ & Low (often confused) \\
o & o, ô & /ɔ/, /ã/ & Low \\
u & u & /ə/ & Low \\
oo & oo & /uː/ & High \\
ee & ee & /iː/ & High \\
au & au, âu & /aw/, /ɔː/ & Moderate \\
b & p (never b) & /p/ & Context-dependent (normalize to p) \\
p & p & /p/ & High \\
t & t & /t/ & High \\
k/c & k & /k/ & High \\
ch & ch & /č/ & High \\
sh & sh & /š/ & High \\
qu & q & /kʷ/ & High \\
m & m & /m/ & High \\
n & n & /n/ & High \\
w & w & /w/ & High \\
y & y & /y/ & High \\
s & s & /s/ & High \\
h & h & /h/ (including preaspiration) & Missing in Wood \\
g & k (never g) & /k/ & Normalize to k \\
\hline
\end{tabular}
\end{table}

Key difference: Eliot uses a systematic vowel-length distinction (â, ô vs.\ a, o)
which Wood entirely lacks. Eliot also notes preaspiration (h) where Wood never
writes it. Eliot never uses \emph{b} or \emph{g}, always \emph{p} and \emph{k},
confirming the voicing-neutralization analysis.

\subsection{Triple-Spelling Deep Dive: ``Devil''}

The devil word is a case study in Wood's orthographic instability:

\begin{table}[htbp]
\centering
\caption{Wood's Spellings for ``Devil''}
\label{tab:1367-devil}
\begin{tabular}{ll}
\hline
Spelling & Source \\
\hline
abbomocho & Primary Wood spelling \\
abomocho & Variant (single b) \\
abamacho & Variant (a for o in second position) \\
abamocho & Variant \\
\hline
\end{tabular}
\end{table}

Proposed reconstruction: /apomaːčə/ or /apomaːčo/.

\begin{itemize}
  \item \emph{a-} prefix (demonstrative/article) $\rightarrow$ segmentable
  \item \emph{bb} $\rightarrow$ /p/ (geminate or short-vowel marker)
  \item \emph{o} $\sim$ \emph{a} in second syllable $\rightarrow$ /a/ or /ə/
    (vowel reduction)
  \item \emph{m} $\rightarrow$ /m/ (stable across all variants)
  \item \emph{ch} $\rightarrow$ /č/ (English $\langle$ch$\rangle$ = /tʃ/)
  \item final \emph{o} $\rightarrow$ /o/ or /ə/
\end{itemize}

The fact that Wood uses the same word for ``devil'' that appears in other
contexts (without Christian gloss) suggests he is mapping the Christian concept
onto an indigenous spirit-being category. This triple spelling is also useful
diagnostically: if the same word appears in Wood with multiple spellings, the
stable consonants (m, ch) are reliable while the unstable vowels (o/a, o/a)
indicate reduced/unstressed vowels that English ears struggled to categorize.

\subsection{A-Prefix Analysis}

The \emph{a-} prefix is one of the most distinctive features of Wood's
orthography. In the 339-entry word list, roughly 25--30\% of nouns and
adjectives begin with \emph{a-}. Wood treats \emph{a-} as part of the root (no
word-boundary space), suggesting either:

\begin{enumerate}
  \item \textbf{English hearer error:} Wood heard the prefixed form as a single
    phonological word and wrote it as such
  \item \textbf{Demonstrative prefix:} Wampanoag had a demonstrative/determiner
    prefix \emph{a-} (cf.\ Massachusett \emph{a-} ``that, yonder'') that was
    frequently attached to nouns in connected speech
  \item \textbf{Indefinite prefix:} Prefixed to indefinite/unspecified referents
    (a man, not the man)
  \item \textbf{Nominalizing prefix:} The \emph{a-} may be a nominalizer on verb
    stems
\end{enumerate}

Cross-source comparison: The \emph{a-} is absent in Eliot --- suggesting either
Eliot was more systematic about morphological segmentation, the \emph{a-} was
a casual-speech feature that Eliot deliberately excluded from his written
standard, or Wood heard and recorded a discourse feature (article-like prefix)
that Eliot's educated linguistic eye removed.

\subsection{Pidginization Theory}

Several scholars have suggested that Wood's word list may represent not fluent
Wampanoag but a pidginized or foreigner-talk register.

\subsubsection{Goddard's Analysis (1977, 1981)}

Ives Goddard noted that early colonial word lists often show features of
reduced morphology --- the characteristic of pidgin/foreigner-talk registers:

\begin{itemize}
  \item \textbf{Missing inflectional endings:} Wood's entries notably lack the
    inflectional suffixes that Eliot's grammar requires
  \item \textbf{Reduced pronominal marking:} Missing the complex subject/object
    marking that Eliot documents
  \item \textbf{Lexical rather than grammatical:} Wood records content words
    (nouns, verbs in citation form) rather than inflected forms
  \item \textbf{Simplified phonology:} Some consonant clusters that Eliot writes
    are simplified in Wood's transcriptions
\end{itemize}

\subsubsection{Evidence in Wood's List}

\begin{table}[htbp]
\centering
\caption{Pidginization Evidence in Wood}
\label{tab:1367-pidgin}
\begin{tabular}{lll}
\hline
Feature & Example & Pidgin Characteristic \\
\hline
Citation-form verbs & sucqunnocquock (sleeps) --- appears as-is, not inflected & Foreigner-talk citation \\
Missing possessive prefixes & an nu ocke (a bed) --- literal ``a bed'' not ``my bed'' & Reduced morphology \\
Paratactic phrases & abonetta ta sucqunnocquock (six sleeps) --- juxtaposed numerals+noun & No numeral agreement \\
Unusual word order & Some entries show English word order with Wampanoag words & Code interference \\
\hline
\end{tabular}
\end{table}

\subsubsection{Counter-Argument}

The numeral system argues against full pidginization: Wood's number series
\emph{nawhut}-\emph{nees}-\emph{quín}-\emph{yoâw}-\emph{abbona} is a close
match to the attested Wampanoag number system. Pidgins typically simplify or
restructure numeral systems; Wood's closely matches the standard language.

\subsection{Problems Encountered}

\begin{itemize}
  \item \textbf{b/p alternation:} Wood writes \emph{b} where Eliot/Trumbull use
    \emph{p} --- direct evidence of voicing neutralization. The same Wampanoag
    phoneme (unaspirated /p/) is variously written as \emph{b} or \emph{p}
    with no consistent rule.
  \item \textbf{Triple spelling:} \emph{abbomocho}/\emph{abomocho}/
    \emph{abamacho} for ``devil'' --- strongest evidence of colonial vowel
    instability in the corpus. The stable consonants (m, ch) are reliable while
    the unstable vowels (o/a, o/a) indicate reduced/unstressed vowels.
  \item \textbf{Word-initial \emph{a-} prefix:} Many entries begin with
    \emph{a-} that may be a demonstrative/indefinite prefix --- must be
    segmented before conversion. Roughly 25--30\% of nouns and adjectives are
    affected.
  \item \textbf{Geminate consonants:} Frequent doubling (\emph{bb, pp, nn, ck,
    gg}) likely marks preceding vowel as short rather than true gemination.
  \item \textbf{Preaspiration /h/ unreported:} Like all pre-Eliot sources, Wood
    does not write preaspirated /h/ before stops.
  \item \textbf{Vowel quality instability:} Multiple spellings for the same word
    show Wood had trouble distinguishing Wampanoag vowels.
  \item \textbf{Multi-word entries:} Many entries are phrases, not single
    lexemes (e.g., \emph{abonetta ta sucqunnocquock} = ``five sleeps'',
    \emph{appepes naw aug} = ``when I see it I will tell you'').
  \item \textbf{No phonetic guide:} No \emph{\textbackslash ph} entries for
    validation.
\end{itemize}

\subsection{Research Approach}

Treat Wood as unreliable but corroborative. When Wood agrees with
Trumbull/Eliot, it confirms the reading. When disagreeing, prefer Trumbull.
Normalize \emph{b} to /p/ universally. FSM with vowel averaging across Wood's
multiple spellings.

\subsubsection{Conversion Pipeline Recommendations}

\textbf{Priority 1: Normalize Consonants (High Confidence)}

\begin{itemize}
  \item \emph{b} $\rightarrow$ \emph{p} (always)
  \item \emph{g} $\rightarrow$ \emph{k} (always, but check if voiced /g/ in
    context)
  \item \emph{ck} $\rightarrow$ \emph{k} (always)
  \item \emph{qu} $\rightarrow$ \emph{kʷ} (keep as digraph or convert to q for
    phonemic)
\end{itemize}

\textbf{Priority 2: Normalize Vowels (Moderate Confidence)}

\begin{itemize}
  \item All short-vowel candidates treated as underspecified /ə/ $\sim$ /a/
    unless confirmed by Eliot
  \item Use Eliot/Trumbull/Winslow cross-reference as tiebreaker
  \item Double vowels (oo, ee) are more stable and map to /uː/, /iː/
\end{itemize}

\textbf{Priority 3: Morphological Segmentation}

\begin{itemize}
  \item Segment \emph{a-} prefix where cross-reference supports it
  \item Do not segment otherwise
  \item Flag ambiguous cases
\end{itemize}

\textbf{Priority 4: Preaspiration /h/}

\begin{itemize}
  \item Insert /h/ between vowels and following stops where Eliot confirms
  \item Default rule: when a short vowel is followed by a stop consonant in a
    stressed syllable, check Eliot for /h/
  \item Without Eliot confirmation, do NOT insert /h/ speculatively
\end{itemize}

\textbf{Priority 5: Multi-Word Entries}

\begin{itemize}
  \item Tokenize by whitespace
  \item Convert each token independently
  \item Reconstruct the phrase from converted tokens
\end{itemize}

\subsection{Conversion Workflow}
\texttt{[Pending implementation --- see Issue \#1367]}

The conversion pipeline will be implemented as a Python module in
\texttt{src/commons/parsing/}. The module will:

\begin{itemize}
  \item Treat Wood as an unreliable but corroborative source
  \item Normalize \emph{b} $\rightarrow$ /p/ universally across all Wampanoag
    entries
  \item Use FSM with vowel averaging across Wood's multiple spellings
  \item Segment \emph{a-} prefix on a word-by-word basis
  \item Handle multi-word entries by whitespace tokenization
  \item Insert preaspiration /h/ only where Eliot confirms it
  \item Verify each converted entry against expected Eliot/Trumbull forms
\end{itemize}

\subsection{Linguist Feedback Template}

\texttt{[Template --- content pending review]}

The following questions are open for linguist review:

\begin{enumerate}
  \item \textbf{b/p normalization:} Is the universal \emph{b} $\rightarrow$
    /p/ normalization correct, or are there contexts where Wood's \emph{b}
    represents a different phoneme?
  \item \textbf{a- prefix segmentation:} Is the \emph{a-} prefix a genuine
    Wampanoag morpheme (demonstrative/indefinite) or an English hearer error?
    Should it be segmented universally or on a case-by-case basis?
  \item \textbf{Pidginization:} Does the numeral system's accuracy argue against
    full pidginization, or can a language have accurate numerals alongside
    reduced morphology?
  \item \textbf{Preaspiration insertion:} Is the conservative approach (insert
    /h/ only where Eliot confirms) correct, or are there systematic contexts
    where /h/ can be reliably reconstructed?
  \item \textbf{``Devil'' reconstruction:} Is /apomaːčə/ a plausible
    reconstruction, or does the root derive from a different PA source?
\end{enumerate}

\subsection{Related Work: Colonial Recorder Perception Problem}

The core challenge of this research card --- English-speaking recorders
misperceiving and mis-transcribing Algonquian phonemes --- is a known
phenomenon documented across Eastern North America. Two researchers' work is
directly relevant:

\textbf{Dr.\ Keith Cunningham} (Linguist, Nanticoke Indian Tribe). His doctoral
dissertation \emph{A Phonological Analysis of Nanticoke With Practical
Applications for Language Revitalization} (Georgetown University) analyzes three
18th-century Nanticoke word lists (Heckewelder 1785) using the same methodology:
historical phonology from colonial recorders, with the same challenges of
multilingual informants and orthographic complexity. His 2024 Algonquian
Conference presentation \emph{A Revised Phonological Analysis of the Heckewelder
Vocabulary of Nanticoke} directly parallels the Wood analysis --- both involve
pre-missionary English recorders with no linguistic training, both show b/p
alternation as evidence of voicing neutralization, and both require PA
comparative calibration. Cunningham's work demonstrates that the colonial
recorder perception problem extends beyond SNEA into the broader Eastern
Algonquian family.

\textbf{Dr.\ Craig Kopris} (Independent scholar). His paper \emph{Les sons
wyandots perçus par des oreilles étrangères} (Wyandot sounds perceived by
foreign ears, \emph{Recherches amérindiennes au Québec}, 2014) is the exact
framing of the Wood problem: how European recorders systematically misperceived
and mis-transcribed indigenous phonemes. His \emph{Wyandot Phonology: Recovering
the Sound System of an Extinct Language} applies the same recovery methodology
to an unrelated language family (Iroquoian), demonstrating that the foreign-ear
perception problem is universal across colonial documentation of North American
languages.

The Wood corpus (339 records, the second-largest pre-Eliot Wampanoag source) is
the best case study for the pre-orthographic recorder problem. Wood's triple
spelling for a single word (\emph{abbomocho}/\emph{abomocho}/\emph{abamacho}) is
the strongest evidence in the entire SNEA corpus for colonial vowel instability
--- a phenomenon that both Cunningham (Nanticoke) and Kopris (Wyandot) document
in their respective language families.

\subsection{Key References}

\begin{itemize}
  \item Aubin, George F. 1975. \emph{A Proto-Algonquian Dictionary}. Ottawa:
    National Museums of Canada.
  \item Eliot, John. 1666. \emph{The Indian Grammar Begun}. Cambridge, MA:
    Marmaduke Johnson.
  \item Goddard, Ives. 1977. ``Some Observations on the Massachusetts
    Language.'' \emph{Papers of the 8th Algonquian Conference}.
  \item Goddard, Ives. 1981. ``Massachusetts Phonology.'' \emph{International
    Journal of American Linguistics} 47(4): 283--305.
  \item Goddard, Ives \& Bragdon, Kathleen. 1988. \emph{Native Writings in
    Massachusett}. 2 vols.\ Philadelphia: APS.
  \item Salvucci, Claudio. 2008. \emph{The Vocabulary of Massachusett}.
    Evolution Publishing.
  \item Salisbury, Neal. 1982. \emph{Manitou and Providence: Indians, Europeans,
    and the Making of New England, 1500--1643}. Oxford University Press.
  \item Trumbull, James Hammond. 1903. \emph{Natick Dictionary}. Bureau of
    American Ethnology Bulletin 25.
  \item Vaughan, Alden T. 1965. \emph{New England Frontier: Puritans and
    Indians 1620--1675}. Little, Brown.
  \item Wood, William. 1634. \emph{New Englands Prospect}. London: Tho.\ Cotes
    for John Bellamie. STC 25957.
  \item Wood, William. 1865. \emph{Wood's New England's Prospect}. Boston:
    Prince Society. (Reprint of 1634 with notes.)
\end{itemize}

% Mohegan-Pequot (Prince-Speck 1904) — Source Workflow
% Issue #1368
% Status: Research complete, implementation pending

\section{Mohegan-Pequot: Prince-Speck 1904}
\label{sec:1368}

\subsection{Source Description}

\subsubsection{Publication History and Citation Clarification}

The glossary was published twice in slightly different venues, which causes
persistent citation confusion. The first printing appeared in \emph{American
Anthropologist} n.s.\ 6(1): 18--45 (January--March 1904). A typographically
identical reprint appeared in the \emph{Boas Anniversary Volume} (G.\ E.\
Stechert, New York, 1906), pp.\ 18--45. Because the glossary appeared first
in \emph{American Anthropologist} in 1904, that is the correct year for
scholarly citation. The 1906 reprint is typographically identical---same page
numbering, same text---which is why later references often cite the Boas
volume without noticing the earlier \emph{AA} publication.

The exact title is ``Glossary of the Mohegan-Pequot Language'' (not ``The
Mohegan-Pequot Language'' as sometimes mis-cited). The work is a glossary
with brief ethnographic introduction, not a grammar.

\subsubsection{J. Dyneley Prince (1868--1945): The Transcriber}

John Dyneley Prince was a remarkable polymath: diplomat, politician, and
academic linguist. Born in New York City in 1868, he earned his A.B.\ (1888)
and Ph.D.\ (1892) from Columbia University and also studied in Europe. He
served as Professor of Germanic Languages at Columbia (later Dean), as a New
Jersey State Senator, and as U.S.\ Minister to Denmark (1921--1926) and
Yugoslavia (1926--1932). He was reported to know approximately 30 languages
including Estonian, Turkish, Arabic, Slavic languages, and Hungarian.

Prince was trained in the \textbf{Neogrammarian (Junggrammatiker)} tradition,
which dominated Germanic linguistics in the late 19th century. This background
has specific consequences for his transcription:

\begin{itemize}
  \item \textbf{Phonetic impressionism over phonemic analysis:}
    Neogrammarians believed in recording ``exact'' phonetic detail rather
    than phonemic contrasts. Prince's transcription therefore varies between
    \textit{a}, \textit{ah}, \textit{au}, \textit{o} for what is likely a
    single phoneme /aː/ or /ɔː/, writes vowel length inconsistently, and
    records allophonic variation as if it were phonemic.
  \item \textbf{Germanic spelling conventions:} Prince unconsciously applied
    German orthographic habits even when intending to use English spelling
    (see \S\ref{sec:1368-german}).
  \item \textbf{No IPA training:} The glossary predates the International
    Phonetic Alphabet (1888/1900). Prince had no standardized phonetic
    alphabet---he used a blended English-German ad-hoc system.
\end{itemize}

\textbf{Critical distinction:} Prince was NOT a fieldworker. He transcribed
from Speck's field notes. He never heard the language spoken himself except
possibly in brief demonstration. This means every transcription is two steps
removed from the speaker: Fidelia Fielding $\rightarrow$ Speck's ears
$\rightarrow$ Speck's field notation $\rightarrow$ Prince's desk transcription
$\rightarrow$ published form.

\subsubsection{Frank G. Speck (1881--1950): The Fieldworker}

Frank G. Speck was a Boasian anthropologist whose fieldwork methodology was
ethnographic rather than strictly linguistic. Born in New York in 1881, he
studied at Columbia University under Franz Boas (Ph.D.\ 1905) and later
taught at the University of Pennsylvania. He was a prolific ethnographer of
Eastern Algonquian peoples (Penobscot, Mohegan, Nanticoke, Tutelo, etc.).

Speck conducted fieldwork with \textbf{Fidelia Fielding (Aiji, 1827--1908)},
the last fluent speaker of Mohegan-Pequot. Fielding was born in 1827 on the
Mohegan reservation in Uncasville, Connecticut, a direct descendant of the
Mohegan chief Uncas. She was the keeper of Mohegan traditions, maintained a
diary in Mohegan using her own spelling system, was a monolingual Mohegan
speaker as a child (learning English later in life), and was known for her
fierce preservation of the language---she refused to speak English with other
Mohegans.

Fielding was literate in English and developed her own orthography for
recording Mohegan, primarily in her personal diaries. Her spelling system
was based on English conventions as she understood them, with considerable
idiosyncrasy.

\subsubsection{The Speck--Prince Handoff}

The exact chain of custody for the data is:

\begin{enumerate}
  \item Fidelia Fielding produces utterances in Mohegan
  \item Frank Speck hears and records them in field notation
  \item Speck provides his field notes to J.\ Dyneley Prince
  \item Prince converts Speck's notation into the published orthography
  \item Prince writes the introductory notes and publishes under both names
\end{enumerate}

This two-stage recording (fieldworker $\rightarrow$ transcriber) introduces
compounding transcription errors: Speck's hearing of Fielding's speech
(accent, speed, dialect), Speck's field notation choices (ad-hoc English-based
spelling), Prince's interpretation of Speck's notation (he never heard the
original), Prince's application of German orthographic conventions, and
typesetting errors in the 1904 publication.

\subsubsection{Fidelia Fielding's Orthography and Its Influence}

Prince and Speck explicitly state that they sometimes deferred to Fielding's
spelling for ``sentimental reasons''---i.e., they kept her English-based
spelling even when a more phonetic transcription would have been technically
better. From analyses of her diaries, Fielding's orthography shows:

\begin{itemize}
  \item English vowel letters for approximate sounds (\textit{a, e, i, o, u}
    as in English orthography) with no consistent length marking
  \item \textit{gh} for /k/ or /x/, following English convention, which may
    reflect actual velar fricative
  \item \textit{w} for /w/ and as glide, consistent with English
  \item Doubled consonants for occasional geminate marking, which may reflect
    actual length or stress
  \item No systematic nasalization marking---nasal vowels were likely present
    but unwritten
\end{itemize}

The Prince-Speck preface indicates they intentionally preserved some of
Fielding's spellings because she was the last speaker and her written form
had cultural value. This means some words in the glossary may be Fielding's
\textbf{written} spelling rather than Prince's \textbf{hearing-based}
transcription, and the distinction between ``written by Fielding'' and
``recorded by Speck/Prince'' is not marked in the glossary.

\subsubsection{Orthography System}

The glossary presents a three-way orthographic competition:

\begin{table}[h]
\centering
\begin{tabular}{llll}
\toprule
\textbf{Source} & \textbf{Convention} & \textbf{Example} & \textbf{Interpretation} \\
\midrule
English colonial & \textit{ch} for /tʃ/, \textit{sh} for /ʃ/ & ``cheese'' & Familiar, predictable \\
English vernacular (Fielding) & Ad-hoc English-based & Variable spelling & Unpredictable \\
German linguistic (Prince) & \textit{tsh} for /tʃ/, \textit{z} for /ts/, Dehnungs-\textit{h} & Specific forms & Detectable pattern \\
\bottomrule
\end{tabular}
\caption{Orthographic confusion index}
\label{tab:1368-orthographic}
\end{table}

\subsection{Phonemic Reconstruction}
\label{sec:1368-phonemic}

\subsubsection{Consonant Inventory (Reconstructed)}

\begin{table}[h]
\centering
\begin{tabular}{lccccc}
\toprule
 & Bilabial & Alveolar & Palatal & Velar & Glottal \\
\midrule
Stops & p & t & & k & \\
Voiced stops & b? & d? & & g? & \\
Affricates & & & tʃ & & \\
Fricatives & & s & ʃ & (x?) & h \\
Nasals & m & n & & & \\
Approximants & w & & j & & \\
\bottomrule
\end{tabular}
\caption{Reconstructed consonant inventory for Mohegan-Pequot}
\label{tab:1368-consonants}
\end{table}

Notes: \textit{b, d, g} are likely English loan phenomena or allophones of
/p, t, k/. The velar fricative /x/ may exist but is not clearly distinguished
from /h/ or /k/ in the transcription. Preaspiration /ʰp, ʰt, ʰk/ is probably
a suprasegmental feature of the syllable coda, not a phoneme sequence.

\subsubsection{Vowel Inventory (Reconstructed)}

\begin{table}[h]
\centering
\begin{tabular}{lccc}
\toprule
 & Front & Central & Back \\
\midrule
High & i /iː/ & & u /uː/ \\
Mid & e /ɛ/ & ə /ə/ & o /o/ \\
Low & & a /a/ & ɔ /ɔː/ \\
\bottomrule
\end{tabular}
\caption{Reconstructed vowel inventory for Mohegan-Pequot}
\label{tab:1368-vowels}
\end{table}

Diphthongs: /aj/ (spelled \textit{ai}), /aw/ (spelled \textit{au}), /oj/
(if present).

\subsubsection{Preaspiration Phonology}

Preaspiration in SNEA languages is a debated feature. The Prince-Speck
evidence suggests that preaspiration occurred before fortis stops in coda
position: \textit{ahk}, \textit{ahp}, \textit{aht}. It may have been a
syllable-level prosodic feature (fortis/lenis distinction) rather than a
segment. Prince-Speck writes it as \textit{hC} sequences, but these may
represent single preaspirated segments /ʰC/. The fact that it appears in
Prince-Speck and not in earlier colonial sources may mean: (a) it developed
later in Mohegan-Pequot, (b) earlier recorders missed it, or (c) Prince's
Germanic ear was better tuned to it.

\subsubsection{Nasalization}

Mohegan-Pequot had nasal vowels (inherited from Proto-Algonquian *ē, *ā,
etc.). Prince-Speck does not systematically mark nasalization, may
occasionally indicate it through vowel quality choices (e.g., \textit{o} for
nasalized /ã/), and may miss it entirely, mapping nasal vowels to oral
vowels. Fielding's Mohegan has lost nasal vowels, so the modern source is no
help for this feature.

\subsection{German Orthographic Influence}
\label{sec:1368-german}

Prince's Germanistic training is visible throughout the transcription.

\subsubsection{Affricate Representations}

\begin{table}[h]
\centering
\begin{tabular}{llll}
\toprule
\textbf{Prince-Speck} & \textbf{German convention} & \textbf{English equivalent} & \textbf{Phoneme} \\
\midrule
\textit{tsh} & German \textit{tsch} (as in \textit{deutsch}) & Rare, used in ``batsha'' & /tʃ/ \\
\textit{z} & German \textit{z} = /ts/ & English \textit{z} = /z/ & ambiguous \\
\bottomrule
\end{tabular}
\caption{Affricate representations}
\label{tab:1368-affricates}
\end{table}

The use of \textit{tsh} for the /tʃ/ affricate is a distinctly Germanic
choice. English colonial recorders typically used \textit{ch} or \textit{tch}.
The \textit{ts} prefix represents a German-trained linguist's desire to show
the affricate's stop onset.

\subsubsection{Sibilant Representations}

\begin{itemize}
  \item \textit{sch} (if present): German \textit{sch} = /ʃ/ (English
    \textit{sh} would be more typical of American recorders)
  \item \textit{s} before vowels: German \textit{s} = /z/ in some positions
    (English reads this as /s/)
\end{itemize}

\subsubsection{Vowel Length Marking}

German linguistics of the period marked vowel length with vowel doubling
(\textit{ee}, \textit{oo})---shared with English, so the source is
ambiguous---and \textit{h} after vowel (\textit{ah}, \textit{eh})---German
uses \textit{Dehnungs-h} (lengthening h), which is NOT the same as English
``ah'' for /aː/. Prince's \textit{ah} spelling could derive from either
English ``ah'' (indicating quality as much as length) or German
\textit{Dehnungs-h} (indicating length primarily). The ambiguity cannot be
resolved without acoustic evidence.

\subsubsection{Consonant Fortition}

German-influenced recorders sometimes write initial /b, d, g/ as \textit{b,
d, g} but intend fortis articulation, apply final devoicing even when
hearing voiced finals---writing \textit{d} but hearing /t/, and use
unaspirated \textit{p, t, k} where English recorders would write \textit{b,
d, g}.

\subsubsection{Umlaut-Like Notation}

German linguists sometimes used \textit{ü}, \textit{ö}, \textit{ä} when
hearing sounds not in English or German. Whether Prince does this in the
glossary requires a full search, but it is a known pattern in
German-trained linguist transcriptions.

\subsection{Preaspiration Evidence}

\subsubsection{Explicit /h/ in Prince-Speck}

One of Prince-Speck's most valuable features for SNEA historical phonology
is its treatment of /h/:

\begin{itemize}
  \item \textbf{Word-initial h:} \textit{ahupanun} (come here)---English
    recorders often drop word-initial h
  \item \textbf{Intervocalic h:} \textit{bahduntah} (rising)---often missed
    by colonial recorders
  \item \textbf{Cluster hC:} \textit{bahkeder} (bitter?)---PREASPIRATION,
    critical for SNEA reconstruction
\end{itemize}

\subsubsection{Preaspiration (/hk/, /hp/, /ht/)}

The spelling \textit{hk} in \textit{bahkeder} is significant because it
represents a phonetic feature that English-colonial recorders systematically
miss. Roger Williams (Narragansett) writes plain \textit{k} for both /k/ and
/hk/. John Eliot (Wampanoag) uses doubled vowels to indicate syllable weight
but does not write preaspiration. Ezra Stiles (Mohegan-Pequot) sometimes
writes \textit{hk} but inconsistently. \textbf{Prince writes \textit{hk},
\textit{hp}, \textit{ht} clusters explicitly}---this is unique among the
pre-modern SNEA sources.

\subsubsection{Why Prince-Speck Catches Preaspiration}

Prince's Germanic training gave him an ear for aspiration contrasts. German
distinguishes aspirated /pʰ, tʰ, kʰ/ in certain positions from unaspirated
/p, t, k/ in clusters, and the \textit{Dehnungs-h} convention made him
conscious of h-like elements. Additionally, Mohegan-Pequot may have preserved
preaspiration longer than other SNEA languages, or Speck's field notes may
have captured it from Fidelia Fielding's careful speech.

\subsubsection{Comparison with Other Sources}

\begin{table}[h]
\centering
\begin{tabular}{llll}
\toprule
\textbf{Source} & \textbf{Preaspiration /hk/} & \textbf{Intervocalic /h/} & \textbf{Word-initial /h/} \\
\midrule
Williams 1643 & Not written & Sometimes & Rarely \\
Eliot 1660s & Not written & Rarely & Rarely \\
Stiles 1760s--90s & Inconsistent & Inconsistent & Inconsistent \\
\textbf{Prince-Speck 1904} & \textbf{Written} & \textbf{Written} & \textbf{Written} \\
Fielding 2013 & Not phonemic & Not present & Not present \\
\bottomrule
\end{tabular}
\caption{Preaspiration and /h/ recording across SNEA sources}
\label{tab:1368-preaspiration}
\end{table}

This makes Prince-Speck a \textbf{key witness for preaspiration} in SNEA
historical phonology, on par with Mayhew and Bourne (Wampanoag) for the same
feature.

\subsection{English Loanword Typology}

The Prince-Speck glossary contains a very high proportion of English
loanwords, reflecting the contact situation of early 20th-century Mohegan.

\subsubsection{Loanword Categories}

\textbf{Category 1: Cultural/Technological (new concepts)}

\begin{itemize}
  \item \textit{beed} (bed), \textit{beeddunk} (bedstead), \textit{beesh}
    (peas), \textit{beetkuz} (lady's dress), \textit{manshetshe} (matches),
    \textit{tchiken} (chicken), \textit{breeches} (breeches), \textit{cock}
    (rooster), \textit{nute} (nut), \textit{rug} (rug), \textit{soyler}
    (soldier), \textit{spuner} (spoon), \textit{shtoo} (stew), \textit{tappe}
    (tap/faucet), \textit{tool} (tool), \textit{wamp} (wampum, earlier
    borrowing)
\end{itemize}

\textbf{Category 2: Extended/Replaced Domains}

\begin{itemize}
  \item \textit{lif} (leaf---English ``leaf'' replacing native word),
    \textit{appece} (apple---meaning extended to any fruit), \textit{cheese}
    (cheese), \textit{beksees} (pig---English ``pigs'' with Mohegan suffix?),
    \textit{square} (square)
\end{itemize}

\textbf{Category 3: Calendric/Temporal}

\begin{itemize}
  \item \textit{beitar} (Friday---possibly from English ``Friday'' via Dutch
    ``Vrijdag''?)
\end{itemize}

\textbf{Category 4: Phonological Nativization}

\begin{itemize}
  \item \textit{manshetshe} (matches): /tʃ/ $\rightarrow$ \textit{tsh};
    final vowel epenthesis
  \item \textit{soyler} (soldier): r $\rightarrow$ /l/ (no native /r/);
    final /ər/ $\rightarrow$ /er/
  \item \textit{spuner} (spoon): no initial cluster simplification;
    final /n/ $\rightarrow$ /nər/
  \item \textit{appece} (apple): /l/ $\rightarrow$ ?; final schwa
    $\rightarrow$ \textit{e}
\end{itemize}

\subsubsection{Identifying Loanwords}

Loanwords in the glossary can be identified by phonotactic cues (presence of
/r/, /l/, initial /b,d,g/ without alternation, consonant clusters not typical
of Algonquian), semantic cues (European cultural artifacts, introduced
concepts, calendric terms), and orthographic cues (English spelling patterns
that diverge from Mohegan-Pequot patterns).

\subsubsection{Handling in Conversion Pipeline}

Loanwords should be kept with identity mapping (loanword $\rightarrow$ same
word) rather than pushed through Mohegan-Pequot phoneme conversion. The
detection FSM should flag suspected loanwords at runtime based on phonotactic
patterns and semantic context.

\subsection{Comparison with Other SNEA Recording Traditions}

\subsubsection{Williams 1643 (Narragansett)}

\begin{itemize}
  \item /h/ recording: Williams poor, Prince-Speck good---Prince better for
    preaspiration
  \item Vowel system: Williams English short/long, Prince-Speck more
    detailed---both need normalization
  \item Loanword handling: Williams marked as loans, Prince-Speck
    unmarked---Prince needs loanword detection
  \item Transcription consistency: Williams moderate, Prince-Speck low
    (two-stage)---Prince less consistent
  \item Grammar coverage: Williams extensive, Prince-Speck minimal---Prince
    is glossary only
\end{itemize}

\subsubsection{Eliot 1660s (Wampanoag Mass Bible)}

\begin{itemize}
  \item Orthography: Eliot semi-standardized, Prince-Speck ad-hoc
  \item Corpus size: Eliot massive (~1000pp), Prince-Speck small (446 words)
  \item Phonemic value: Eliot low (English-colonial), Prince-Speck
    potentially high (German training)
  \item Preaspiration: Eliot not written, Prince-Speck written
\end{itemize}

\subsubsection{Stiles 1760s--90s (Mohegan-Pequot)}

Ezra Stiles (1727--1795), president of Yale, collected Mohegan-Pequot
vocabulary from several speakers (including Occom family members) over
decades.

\begin{itemize}
  \item Time depth: Stiles ~140 years earlier, Prince-Speck later, closer to
    extinction
  \item Number of speakers: Stiles multiple (including Occom), Prince-Speck
    single (Fidelia Fielding)
  \item Transcription: Stiles ad-hoc English, Prince-Speck Germanic-influenced
  \item /h/ and preaspiration: Stiles inconsistent, Prince-Speck explicit
  \item Dialect: both Mohegan-Pequot proper
  \item Loanword rate: Stiles lower, Prince-Speck higher (more contact)
\end{itemize}

Stiles and Prince-Speck together provide diachronic depth---comparing the
same language approximately 140 years apart.

\subsubsection{Comparison with Stephanie Fielding (2013) Modern Mohegan}

Stephanie Fielding (Yôpôwi Yoht, ``Morning Fire''), a Mohegan linguist
educated at MIT, published the \emph{Modern Mohegan Dictionary} in 2013
(revised 2015). This is the only modern comprehensive dictionary of Mohegan.

\begin{table}[h]
\centering
\begin{tabular}{lll}
\toprule
\textbf{Feature} & \textbf{Fielding (2013)} & \textbf{Prince-Speck (1904)} \\
\midrule
Base & Phonemic (designed) & Phonetic (impressionistic) \\
Vowel length & Explicit & Inconsistent \\
/h/ & Not phonemic & Orthographic \\
Preaspiration & Not marked & Marked \\
Nasalization & Not marked & Not marked \\
Loanwords & Identified & Unidentified \\
Consonant system & Simplified (modern) & More complex (archaic) \\
\bottomrule
\end{tabular}
\caption{Orthography comparison: Fielding vs.\ Prince-Speck}
\label{tab:1368-fielding-comparison}
\end{table}

The pipeline of Prince-Speck $\rightarrow$ Fielding $\rightarrow$ Proto-SNEA
is powerful for reconstruction. Fielding gives the modern phonemic form (what
the language sounds like now); Prince-Speck gives the older phonetic form
(what the language sounded like ~110 years ago); differences between them
indicate sound changes in progress. However, Fielding's dictionary is
reconstructive---many entries are not from fluent speakers but from archival
sources (including Prince-Speck). This introduces circularity: Fielding may
have reconstructed a word FROM Prince-Speck, so using Fielding to ``verify''
Prince-Speck is invalid for that word.

\subsection{Problems Encountered}

\begin{itemize}
  \item \textbf{German orthographic influence:} Prince's Neogrammarian
    training imposed German conventions: \textit{tsh} for /tʃ/, \textit{z}
    for /ts/, \textit{Dehnungs-h} for vowel length, final devoicing, and
    possible umlaut-like notation for non-English vowels.
  \item \textbf{English loanword density:} Approximately 20--30\% of entries
    are English loanwords (bed, matches, chicken, spoon, soldier). These must
    be detected and excluded from phoneme conversion.
  \item \textbf{Fielding dictionary circularity:} Many Fielding (2013) entries
    are reconstructed FROM Prince-Speck---using Fielding to verify
    Prince-Speck is circular for those entries. A cross-reference is needed
    to identify which Fielding entries derive from Prince-Speck
    (self-source) vs.\ from other sources (independent verification).
  \item \textbf{Record count correction:} 446 entries, not 465 as previously
    stated in some references. The actual published glossary contains 446
    numbered entries (some are phrases rather than single words).
  \item \textbf{Two-stage recording chain:} Speck collected, Prince
    transcribed---Prince never heard the language spoken. This introduces
    compounding transcription errors at each stage.
  \item \textbf{``Sentimental reasons'' spellings:} Prince and Speck
    intentionally preserved some of Fielding's English-based spellings for
    cultural value, making it impossible to distinguish Fielding's written
    spelling from Prince's hearing-based transcription.
  \item \textbf{Inconsistent vowel length marking:} Prince writes vowel
    length sometimes with \textit{h}, sometimes with doubling, sometimes not
    at all---the source of the length marking (English ``ah'' vs.\ German
    \textit{Dehnungs-h}) is ambiguous.
  \item \textbf{Multi-word entries:} The glossary includes multi-word phrases,
    which complicates conversion because phrases span different syntactic
    contexts.
\end{itemize}

\subsection{Research Approach}

The conversion pipeline follows this sequence:

\begin{enumerate}
  \item \textbf{Loanword detection FSM:} Automated phonotactic scan of all
    446 entries to identify and catalog English loanwords. Loanwords are
    flagged for identity mapping (excluded from phoneme conversion).
  \item \textbf{German orthography normalizer:} Identify all words where
    German conventions (\textit{tsh}, \textit{z}, \textit{Dehnungs-h})
    appear and apply the correct phoneme mapping.
  \item \textbf{Fielding cross-reference with circularity flag:} Map each
    Prince-Speck entry against Fielding (2013) with source tracing---identify
    which Fielding entries derive from Prince-Speck (circular) vs.\
    independent sources.
  \item \textbf{Preaspiration capture:} Extract all \textit{hC} clusters
    from the glossary and classify by position (initial, medial, final).
    Prince-Speck is one of the few sources that writes preaspiration
    (/hC/ clusters) explicitly---this must be preserved.
  \item \textbf{Stiles comparison:} Cross-reference Prince-Speck entries
    against Stiles's Mohegan-Pequot vocabulary for diachronic comparison.
  \item \textbf{Phoneme-to-orthography mapping table:} Full conversion table
    from Prince-Speck spelling to phonemes, covering all 446 entries.
\end{enumerate}

\subsection{Conversion Workflow}
\label{sec:1368-workflow}

\texttt{[Pending implementation --- see Issue \#1368]}

The conversion workflow will implement the research approach described above
as a multi-stage pipeline. Each stage will be implemented as a separate
module with its own test suite.

\subsection{Linguist Feedback}
\label{sec:1368-feedback}

\texttt{[Template --- content pending review]}

The following aspects require linguist review:

\begin{itemize}
  \item German orthography pattern extraction: confirmation of detected
    German conventions
  \item Preaspiration audit: verification of /hC/ cluster classification
  \item Loanword catalog: review of flagged loanwords for accuracy
  \item Fielding cross-reference: validation of circularity detection
  \item Phoneme mapping: review of the conversion table for accuracy
\end{itemize}

\subsection{Related Work: Colonial Recorder Perception Problem}

The core challenge of this research card---English-speaking recorders
misperceiving and mis-transcribing Algonquian phonemes---is a known
phenomenon documented across Eastern North America. Two researchers' work
is directly relevant.

\textbf{Dr.\ Keith Cunningham} (Linguist, Nanticoke Indian Tribe). His
doctoral dissertation \emph{A Phonological Analysis of Nanticoke With
Practical Applications for Language Revitalization} (Georgetown University)
analyzes three 18th-century Nanticoke word lists (Heckewelder 1785) using
the same methodology: historical phonology from colonial recorders, with the
same challenges of multilingual informants and orthographic complexity. His
2024 Algonquian Conference presentation \emph{A Revised Phonological Analysis
of the Heckewelder Vocabulary of Nanticoke} directly parallels the
Prince-Speck analysis---both involve a two-stage recorder chain (fieldworker
$\rightarrow$ desk transcriber) that introduces compounding transcription
errors, and both require identifying the transcriber's native-language
orthographic conventions (German for Prince, English for Heckewelder).
Cunningham's work demonstrates that the colonial recorder perception problem
extends beyond SNEA into the broader Eastern Algonquian family.

\textbf{Dr.\ Craig Kopris} (Independent scholar). His paper \emph{Les sons
wyandots perçus par des oreilles étrangères} (Wyandot sounds perceived by
foreign ears, \emph{Recherches amérindiennes au Québec}, 2014) is the exact
framing of the Prince-Speck problem: how European recorders systematically
misperceived and mis-transcribed indigenous phonemes. His \emph{Wyandot
Phonology: Recovering the Sound System of an Extinct Language} applies the
same recovery methodology to an unrelated language family (Iroquoian),
demonstrating that the foreign-ear perception problem is universal across
colonial documentation of North American languages. His Dictionary of
Wyandot project (NEH award FN-255576-17) shows the long-term application of
these methods to language revitalization.

The Prince-Speck glossary (446 entries) is unique among the SNEA sources in
having a documented two-stage fieldworker-transcriber chain (Speck collected,
Prince transcribed) with known linguistic biases (Prince's Neogrammarian
training). This makes it the most directly comparable case to the
recorder-chain problems that both Cunningham and Kopris address in their
respective language families.

\subsection{References}

\begin{itemize}
  \item Prince, J.\ Dyneley \& Speck, Frank G.\ (1904). ``Glossary of the
    Mohegan-Pequot Language.'' \emph{American Anthropologist} n.s.\ 6(1):
    18--45.
  \item Prince, J.\ Dyneley \& Speck, Frank G.\ (1906). ``Glossary of the
    Mohegan-Pequot Language.'' In \emph{Boas Anniversary Volume}, pp.\
    18--45. New York: G.\ E.\ Stechert. [Reprint, typographically identical]
  \item Prince, J.\ Dyneley. (1907). ``A Singular Mohegan-Pequot Form.''
    \emph{American Anthropologist} n.s.\ 9: 496--497.
  \item Fielding, Stephanie. (2013/2015). \emph{Modern Mohegan Dictionary.}
    Mohegan Tribe.
  \item Cowan, William. (1973). ``Pequot from Stiles to Speck.''
    \emph{International Journal of American Linguistics} 39(3): 164--172.
  \item Speck, Frank G.\ (1928). ``Native Tribes and Dialects of Connecticut:
    A Mohegan-Pequot Diary.'' \emph{Annual Report of the Bureau of American
    Ethnology} 43: 199--287.
  \item Goddard, Ives. (1978). ``Eastern Algonquian Languages.'' In
    \emph{Handbook of North American Indians}, Vol.\ 15: Northeast, pp.\
    70--77. Washington: Smithsonian Institution.
  \item Costa, David J.\ (2007). ``The Dialectology of Southern New England
    Algonquian.'' In \emph{Papers of the 38th Algonquian Conference}, pp.\
    81--127.
\end{itemize}

% Mahican (Edwards 1787) — Source Workflow
% Issue #1369
% Status: Research complete, implementation pending

\section{Mahican: Edwards 1787}
\label{sec:1369}

\subsection{Source Description}

\subsubsection{Jonathan Edwards Jr.\ --- Biographical Context}

Jonathan Edwards Jr.\ (May 26, 1745 -- August 1, 1801) was a
Congregationalist minister, theologian, college president, and linguist. He
was the son of the theologian Jonathan Edwards Sr.\ (Princeton's 3rd
president). He graduated from Princeton University in 1765 and served as
president of Union College in Schenectady, New York from 1799 until his death
in office.

Edwards grew up in Stockbridge, Massachusetts, a ``praying town'' where
Mahican people were the majority population (approximately 150 Mahican
families versus approximately 12 white families). He learned Mahican as a
child from playmates, becoming fluent. At age 10 (1755), his father sent him
to live in the Iroquois settlement of Onaquaga (Onohoquaga) to learn Oneida,
giving him first-hand knowledge of both Algonquian and Iroquoian language
families.

Unlike his father, who refused to learn Mahican (calling it ``barbarous''),
Edwards Jr.\ defended the language's complexity, demonstrated that it had
abstract terms (love, hatred, malice, religion), and translated the Lord's
Prayer into Mahican. He consulted a Mahican elder, ``Captain Yohum,'' for
revisions before publication. George Washington, James Madison, and Thomas
Jefferson all received copies of his treatise.

\subsubsection{The 1788 Publication}

\textbf{Full title:} \emph{Observations on the Language of the Muhhekaneew
Indians; in which the Extent of that Language in North-America is shewn; its
Genius Grammatically Traced; some of its Peculiarities, and some Instances
of Analogy between that and the Hebrew are Pointed Out.}

Communicated to the Connecticut Society of Arts and Sciences on October 23,
1787. Published 1788 by Josiah Meigs, New Haven. A second edition with notes
by John Pickering appeared in 1823 (in Massachusetts Historical Collections).

The treatise includes:

\begin{itemize}
  \item A grammatical sketch of Mahican (parts of speech, person marking,
    tense)
  \item A comparative glossary: English--Mahican--Shawnee
  \item A comparative glossary: English--Mahican--Chippeway/Ojibwa
  \item Numerals 1--10 in Mahican and Mohawk
  \item The Lord's Prayer in Mahican and in the ``language of the Six
    Nations''
  \item Word lists organized by semantic domain
  \item Observations on the geographic extent of Algonquian languages
  \item Speculative comparison with Hebrew (a common 18th-century trope)
\end{itemize}

Edwards uses the term ``Mohegan'' interchangeably with ``Muhhekaneew'' for
the language, though he distinguishes it clearly from Mohegan-Pequot
(Southern New England). His dialect is specifically Stockbridge Mahican.

\subsubsection{Orthography System}

Edwards uses an English-based ad-hoc orthography typical of 18th-century
colonial sources. He provides a pronunciation key indicating his intended
sound values for vowel letters.

\textbf{Consonant values (relatively clear):}

\begin{table}[h]
\centering
\begin{tabular}{lll}
\toprule
\textbf{Grapheme} & \textbf{Likely value} & \textbf{Evidence} \\
\midrule
\textit{p, t, k} & /p, t, k/ (unaspirated or lightly aspirated) & PA *p, *t, *k $>$ Mahican /p, t, k/ \\
\textit{b, d, g} & Probably /p, t, k/ after nasals or in voiced contexts & Unclear; English-orthography convention \\
\textit{m, n} & /m, n/ & Straightforward \\
\textit{s} & /s/ & PA *s, *š merged to /s/ in Mahican by modern period \\
\textit{sh} & /š/ & English convention \\
\textit{h} & /h/ & Explicitly written; Edwards is unusually consistent here \\
\textit{w, y} & /w, j/ & Glides \\
\textit{ch} & /č/ & English convention \\
\textit{gh} & \textbf{UNKNOWN} & See \S\ref{sec:1369-gh} \\
\bottomrule
\end{tabular}
\caption{Edwards' consonant graphemes and likely values}
\label{tab:1369-consonants}
\end{table}

\textbf{Vowel values (more ambiguous):}

The key insight (Goddard 2008, Masthay 1991, Warne) is that Edwards uses
\textit{e} to represent /iː/ in many cases. This is confirmed by his
pronunciation key and by cognate comparison.

\begin{table}[h]
\centering
\begin{tabular}{llll}
\toprule
\textbf{Edwards grapheme} & \textbf{Likely Mahican phoneme} & \textbf{PA source} & \textbf{Example} \\
\midrule
\textit{ee} & /iː/ & PA *iː & \textit{neesoh} `two' (PA *niːšwi) \\
\textit{e} & /iː/ or /i/ & PA *i(ː) & \textit{seepoo} `river' (PA *siːpoːwi) \\
\textit{a} & /a/ or /aː/ & PA *a, *eː $\rightarrow$ PEA *aː & \textit{nannah} `five' (PA *nyaːlanwi) \\
\textit{o, oo} & /o/ or /u/ & PA *o(ː) & \\
\textit{u} & /ə/ (schwa) in unstressed; possibly /o/ elsewhere & PA *e $\rightarrow$ PEA *ə & \textit{muhhekaneew} (with mə h-) \\
\textit{au, aw} & /a/ + /w/ or /aː/ & & \\
\textit{uh} & /ə h/ or /ə/ & & Frequent in final syllables \\
\bottomrule
\end{tabular}
\caption{Edwards' vowel graphemes and likely values}
\label{tab:1369-vowels}
\end{table}

\subsection{Phonemic Reconstruction}

\subsubsection{Key Phonological Correspondences (PA $\rightarrow$ Mahican)}

From Warne's analysis of the Schmick manuscript, confirmed across sources:

\begin{table}[h]
\centering
\begin{tabular}{llll}
\toprule
\textbf{PA} & \textbf{Old Mahican} & \textbf{Modern Mahican} & \textbf{Notes} \\
\midrule
*p & /p/ & /p/ & No change \\
*t & /t/ & /t/ & No change \\
*k & /k/ & /k/ & No change \\
*č & /č/ & /č/ & No change \\
*θ{} & /n/ & /n/ & \textbf{Merged} with *n by Old Mahican period \\
*l & /n/ & /n/ & \textbf{Merged} with *n by Old Mahican period \\
*n & /n/ & /n/ & \\
*s & /s/ & /s/ & \\
*š & /s/ & /s/ & Merged to /s/ by Modern Mahican \\
*m & /m/ & /m/ & \\
*h & /h/ & /h/ & \\
*iː & /iː/ & /iː/ & \\
*i & /i/ & /i/ & \\
*eː & /a/ (lax low) & /a/ & \textbf{Retracted/lowered} from *eː to /a/ in PEA \\
*e & /ə/ (schwa) & /ə/ & Reduced to schwa (pan-EA development) \\
*aː & /aː\~ãː/ & /ãː/ & Nasalized in modern Mahican \\
*a & /a/ & /a/ & \\
*oː & /uː/ & /uː/ & \\
*o & /o/ & /o/ & Marginal phoneme even in PA \\
\bottomrule
\end{tabular}
\caption{PA $\rightarrow$ Mahican sound correspondences}
\label{tab:1369-correspondences}
\end{table}

\subsubsection{The /iː/ vs.\ /eː/ Issue}

Recent work (documented in \emph{ImportantMahicanVowel2.pdf} from
munseedelaware.com, with consultation by Goddard and Masthay) has established
that Edwards' \textit{e} frequently represents /iː/, NOT /eː/. Goddard
confirmed (c.\ 2023 email, cited in the Mahican vowel paper) that the second
vowel of the tribal name was absolutely not /eː/ and is properly the ``lower
high'' version of the vowel /iː/, produced by setting the throat and tongue
in a position to say the English word ``hit'' and then producing the sound.

Evidence for this mapping:

\begin{itemize}
  \item Edwards writes \textit{Muhhekaneew}---the second vowel is /iː/
  \item \textit{Neesoh} `two'---cf.\ PA *niːšwi, Munsee /niːšə/
  \item \textit{Seepoo} `river'---cf.\ PA *siːpoːwi, Munsee /siːpə w/
  \item Goddard's 2008 paper transcribes Dick's /kīsōx/ `sun' and /θ{}īpə w/
    `river' with /iː/---even where the English spelling suggests ``e''
  \item Goddard's 2010 phoneme inventory for Mahican lists NO /eː/ phoneme
\end{itemize}

\textbf{Impact on orthography conversion:} Edwards' \textit{e} should NOT be
mapped to /e/ or /ɛ/ by default. The default mapping should be /i(ː)/ unless
cognate evidence specifically indicates a different value.

\subsection{The \textit{gh} Grapheme Mystery}
\label{sec:1369-gh}

Edwards uses \textit{gh} in several words where the phonological value is
uncertain. Cross-linguistic comparison suggests multiple possible values:

\begin{table}[h]
\centering
\begin{tabular}{lll}
\toprule
\textbf{Candidate} & \textbf{Evidence For} & \textbf{Evidence Against} \\
\midrule
/x/ (voiceless velar fricative) & Schmick uses German \textit{ch} for /x/ in cognate forms; PA *h before consonant sometimes appears as /x/ in daughter languages & \\
/h/ & Simple /h/ is the default reflex of PA *h & Edwards also writes \textit{h} separately, so \textit{gh} may be a distinct phoneme \\
/k/ (lenis/voiced allophone) & English \textit{gh} occasionally used for /k/ in colonial sources & Unlikely given that Edwards has \textit{k} for /k/ \\
\bottomrule
\end{tabular}
\caption{Candidate values for the \textit{gh} grapheme}
\label{tab:1369-gh}
\end{table}

\textbf{Recommended resolution:} Use comparative evidence. If a cognate from
Schmick, Munsee, or Wampanoag shows /x/, the mapping is /x/. If it shows /h/,
the mapping is /h/. If no cognate exists, flag as ambiguous.

\subsection{Phonological Processes Affecting Edwards' Orthography}

\subsubsection{Vowel Reduction and \textit{uh}}

Edwards frequently writes \textit{u} and \textit{uh} in unstressed syllables.
The PA *e regularly became schwa /ə/ in Proto-Eastern Algonquian. In
Edwards' orthography, unstressed \textit{u} most likely = /ə/;
\textit{uh} may represent /ə h/ or breathy schwa in a checked syllable.
Munsee shows parallel reduction: PA *e $\rightarrow$ PEA *ə
$\rightarrow$ Munsee /ə/ (raised to /i/ or /u/ before certain clusters).

\subsubsection{Nasalization of *aː}

PA *aː developed a nasal quality in Mahican, marked by Schmick with a tilde
(ã). Warne considers this a ``secondary feature'' that was optional in Old
Mahican but generalized to ALL instances of /aː/ by the Modern Mahican
period. Edwards does NOT mark nasalization, so Edwards' \textit{a} may
represent either oral /aː/ or nasal /ãː/ depending on cognate evidence.

\subsubsection{Word-Final Vowel Deletion}

PA word-final vowels were deleted in Mahican (a pan-Eastern Algonquian
process). Edwards sometimes writes these (reflecting an older/more careful
pronunciation) and sometimes omits them. For example, PA *-i (plural marker)
sometimes appears as \textit{-eh} or \textit{-uh} in Edwards, sometimes
absent.

\subsubsection{Schwa Lowering Before /h, x/}

Before /h/ and /x/, schwa /ə/ is lowered to /a/ in Mahican (Warne).
This accounts for alternations where Edwards writes \textit{a} where PA
reconstruction predicts *e. Example: \textit{ntah} `my heart' $<$ PA *-teːh-
(with *eː $\rightarrow$ *a and schwa-lowering in certain environments).

\subsubsection{Rounding Assimilation}

PA *e rounds to Mahican /o/ when adjacent to *w. This explains forms like
the word for `tree' and `beaver' where Edwards' \textit{o} corresponds to PA
*e. (Warne's rule ordering: rounding before laxing.)

\subsection{Morphological Patterns in Edwards' Data}

Edwards' data reveals several regular morphological patterns that can aid
word segmentation before orthography conversion.

\subsubsection{Person Prefixes}

\begin{table}[h]
\centering
\begin{tabular}{llll}
\toprule
\textbf{Person} & \textbf{Prefix} & \textbf{Edwards example} & \textbf{Gloss} \\
\midrule
1st sg. & n- & \textit{ncmeetseh, ncmeetfeh} & `I eat' \\
2nd sg. & k- & \textit{kcmeetseh, ka meetjeh} & `you eat' \\
3rd sg. & w-/∅- & \textit{pumiffoo} & `he walks' \\
\bottomrule
\end{tabular}
\caption{Person prefixes in Edwards' data}
\label{tab:1369-prefixes}
\end{table}

Edwards himself noted these correspondences and compared them to Hebrew
prefixes, which was one of his major observations.

\subsubsection{Verbal Suffixes}

Edwards records frequent verbal suffixes: \textit{-uh, -eh, -seh} (likely
aspect/mood markers), \textit{-oo} (3rd person singular independent
indicative, e.g., \textit{pumiffoo}), and \textit{-nuh} (1st person plural,
e.g., \textit{ncmeetfchnuh} `we eat').

\subsubsection{Noun Inflection}

Possessive prefixes precede nouns (e.g., \textit{noghkuh} `my father').
Plural formation occurs via suffixation (e.g., \textit{penumpaufoo} `boy'
$\rightarrow$ \textit{penumpaufoouk} `boys').

\subsection{Comparative Sources for Orthography Interpretation}

\subsubsection{The Schmick Manuscript (Primary Internal Source)}

\begin{itemize}
  \item \textbf{Author:} Johann Jacob Schmick, Moravian missionary
  \item \textbf{Period:} c.\ 1753--1767 (before Edwards)
  \item \textbf{Dialect:} Moravian/Western Mahican (Ohio/PA branch)
  \item \textbf{Format:} Two volumes (322 pp.), German-based orthography
  \item \textbf{Publication:} Masthay (ed.), \emph{Schmick's Mahican
    Dictionary}, APS Memoir 197 (1991)
  \item \textbf{Contains:} Mahican historical phonology by David Pentland
\end{itemize}

Schmick's manuscript uses a German-based orthography which is in some ways
more precise than Edwards' English-based system: German \textit{ch} for /x/
(vs.\ Edwards' ambiguous \textit{gh}), more consistent vowel representation,
and extensive vocabulary (much larger than Edwards' 144 records). However,
Schmick shows inconsistency with vowels---especially the reflex of PA *aː
which he writes variously as \textit{a, aa, ã, ã, a} (Warne). The tilde
likely indicates a secondary nasalization feature.

\textbf{Key finding:} Schmick is much more consistent with consonants than
vowels. Identical forms appear written with different vowel symbols or
sequences (Warne, \emph{Time-Depth in Mahican Diachronic Phonology}).

Pentland's historical phonology (in Masthay 1991) provides the connecting
framework between PA reconstructions and Mahican forms. Goddard (2007--2008)
provided important corrections to Pentland's phonology, particularly in the
treatment of vowel length and the PA *e reflexes.

\subsubsection{Zeisberger's Delaware Materials}

David Zeisberger (1721--1808), Moravian missionary, produced extensive
documentation of Munsee Delaware (a close relative of Mahican within the
Delawaran subgroup): \emph{Zeisberger's Delaware Vocabulary of 1772} (earliest
extensive Munsee vocabulary), \emph{Grammar of the Language of the Lenni
Lenape or Delaware Indians} (1827), and the \emph{Diary of David Zeisberger}
(several volumes). These are important for comparative cognate work because
Munsee and Mahican share a common ancestor (Delawaran, per Goddard) and often
show parallel developments.

\subsubsection{Later Mahican Sources}

\begin{itemize}
  \item \textbf{Hendrick Aupaumut writings} (c.\ 1790s--1820s): Stockbridge
    Mahican letters and documents by a fluent Mahican speaker
  \item \textbf{Wm.\ Dick (Michelson notes)} (c.\ 1910s): Modern Mahican
    (Wisconsin) phonetic transcriptions
  \item \textbf{Bernice Robinson recordings} (mid-20th century): Audio
    recordings of the last speakers
\end{itemize}

\subsection{Problems Encountered}

\begin{itemize}
  \item \textbf{The \textit{gh} grapheme mystery:} \textit{ghusooh} (eight)
    contains \textit{gh}---the most unusual grapheme across all eight sources.
    Candidates include /x/ (voiceless velar fricative), /h/, or /k/. Must be
    resolved per cognate on a case-by-case basis. This is the single biggest
    ambiguity.
  \item \textbf{\textit{uh} reduction:} Edwards' orthography is dominated by
    \textit{u} and \textit{uh} sequences representing reduced /ə/. May
    represent /ə/, /ə h/, or /uh/ depending on context.
  \item \textbf{Word-initial consonant clusters:} Highly unusual by English
    standards (\textit{wnissoo}, \textit{hkeesque}, \textit{hpoon},
    \textit{kmeetseh})---indicate rich Mahican morphology. Need morphological
    analysis to determine whether initial CC sequences are prefix+root or
    root-internal.
  \item \textbf{Sentence-level entries:} Many entries are full sentences, not
    lexemes---require morphological segmentation before conversion.
  \item \textbf{Vowel quality ambiguity:} Edwards' English-based vowel letters
    are ambiguous. The /iː/ = \textit{e} mapping (Goddard) is the most
    important correction. \textit{a} could be /a/, /aː/, or /ãː/.
  \item \textbf{No phonetic transcriptions:} Edwards did not use a systematic
    phonetic alphabet. No phonetic entries exist.
  \item \textbf{Small corpus:} Approximately 144 records limits statistical
    approaches.
  \item \textbf{Dialect differences:} Stockbridge Mahican (Edwards) vs.\
    Moravian Mahican (Schmick) may have phonological differences.
  \item \textbf{18th-century conventions:} Edwards' spellings reflect his own
    English pronunciation as much as Mahican sounds; e.g., \textit{oo} may be
    /uː/ or /oː/.
\end{itemize}

\subsection{Research Approach}

The conversion pipeline follows this sequence:

\begin{enumerate}
  \item \textbf{Preprocessing:} Split multi-word entries into individual
    tokens. Flag entries with \textit{gh} for comparative cognate resolution.
    Normalize Edwards' typographic conventions (long s, etc.).
  \item \textbf{Phoneme mapping with cognate-based ambiguity resolution:}
    Apply the default mapping table with overrides per cognate. The
    \textit{e} = /iː/ default (Goddard) is the most important rule.
    \textit{gh} is resolved per cognate (Schmick /x/ or /h/).
  \item \textbf{Morphological segmentation:} Identify and strip person
    prefixes (n-, k-, w-). Segment verbal suffixes (-uh, -eh, -seh, -oo,
    -nuh). Produce root forms for lexicon building.
  \item \textbf{Cross-validation:} Compare converted forms against Schmick's
    phonemicized forms, against Munsee cognates (Zeisberger, Goddard), and
    against PA reconstructions (Pentland, Siebert). Flag entries where
    converted phonemes mismatch expected reflexes.
\end{enumerate}

\subsubsection{Reference Cognate Set for Validation}

The following Edwards forms have well-understood cognates and should be used
as the validation baseline for the conversion:

\begin{table}[h]
\centering
\begin{tabular}{llll}
\toprule
\textbf{Edwards form} & \textbf{Gloss} & \textbf{PA reconstruction} & \textbf{Expected Mahican} \\
\midrule
\textit{neesoh} & `two' & *niːšwi & /niːsə/ \\
\textit{seepoo} & `river' & *siːpoːwi & /siːpə w/ \\
\textit{keesogh} & `sun' & *kiːšoʔθ a & /kiːsox/? \\
\textit{hpoon} & `winter' & *pepoon- & /pə puːn/? \\
\textit{hkeesque} & `eye' & & /kiːskw/? \\
\textit{nannah} & `five' & *nyaːlanwi & /naːnə w/? \\
\textit{penumpaufoo} & `boy' & & \\
\textit{pumiffoo} & `he walks' & *pemihθ eːwa & /pə mihseːwa/? \\
\textit{ncmeetseh} & `I eat' & *miːči- & /nə miːčə/? \\
\textit{ghusooh} & `eight' & *nešwaːšika & \\
\textit{noghkah} & `my father' & & /nə xk-/ \\
\textit{muhhekaneew} & `Mahican' & & /mə hiːkaniːw/ \\
\bottomrule
\end{tabular}
\caption{Reference cognate set for validation}
\label{tab:1369-cognates}
\end{table}

\subsubsection{Software Architecture}

Given the ambiguities and the need for comparative-evidence-based resolution,
the conversion system should be implemented as:

\begin{itemize}
  \item \textbf{FSM core} for the unambiguous mappings (consonants, clear
    vowels)
  \item \textbf{Cognate database} providing priors for ambiguous mappings
    (\textit{gh}, \textit{e} values, vowel length)
  \item \textbf{Resolver module} that queries the cognate database before
    applying ambiguous rules; falls back to a default with a confidence flag
    if no cognate found
  \item \textbf{Output format:} phonemic transcription with confidence
    markers: \texttt{[high]} = confirmed by cognate, \texttt{[med]} = default
    rule no contradicting cognate, \texttt{[low]} = guess no cognate
    available, \texttt{[FLAG]} = contradictory evidence manual review needed
\end{itemize}

\subsection{Conversion Workflow}
\label{sec:1369-workflow}

\texttt{[Pending implementation --- see Issue \#1369]}

The conversion workflow will implement the research approach described above
as a multi-stage pipeline. Each stage will be implemented as a separate
module with its own test suite.

\subsection{Linguist Feedback}
\label{sec:1369-feedback}

\texttt{[Template --- content pending review]}

The following aspects require linguist review:

\begin{itemize}
  \item \textit{gh} grapheme resolution: confirmation of /x/ vs.\ /h/
    mapping per cognate
  \item \textit{e} = /iː/ default: verification of the Goddard finding
    across all entries
  \item Morphological segmentation: review of prefix/suffix identification
  \item Cognate database: validation of reference cognate set
  \item Confidence markers: review of \texttt{[low]} and \texttt{[FLAG]}
    entries
\end{itemize}

\subsection{Related Work: Colonial Recorder Perception Problem}

The core challenge of this research card---English-speaking recorders
misperceiving and mis-transcribing Algonquian phonemes---is a known
phenomenon documented across Eastern North America. Two researchers' work
is directly relevant.

\textbf{Dr.\ Keith Cunningham} (Linguist, Nanticoke Indian Tribe). His
doctoral dissertation \emph{A Phonological Analysis of Nanticoke With
Practical Applications for Language Revitalization} (Georgetown University)
analyzes three 18th-century Nanticoke word lists (Heckewelder 1785) using
the same methodology: historical phonology from colonial recorders, with the
same challenges of multilingual informants and orthographic complexity. His
2024 Algonquian Conference presentation \emph{A Revised Phonological
Analysis of the Heckewelder Vocabulary of Nanticoke} directly parallels the
Edwards analysis---both involve 18th-century recorders with greater linguistic
awareness than their 17th-century predecessors, both document preaspiration
more reliably than earlier sources, and both require PA comparative
calibration. Cunningham's work demonstrates that the colonial recorder
perception problem extends beyond SNEA into the broader Eastern Algonquian
family.

\textbf{Dr.\ Craig Kopris} (Independent scholar). His paper \emph{Les sons
wyandots perçus par des oreilles étrangères} (Wyandot sounds perceived by
foreign ears, \emph{Recherches amérindiennes au Québec}, 2014) is the exact
framing of the Edwards problem: how European recorders systematically
misperceived and mis-transcribed indigenous phonemes. His \emph{Wyandot
Phonology: Recovering the Sound System of an Extinct Language} applies the
same recovery methodology to an unrelated language family (Iroquoian),
demonstrating that the foreign-ear perception problem is universal across
colonial documentation of North American languages. His Dictionary of
Wyandot project (NEH award FN-255576-17) shows the long-term application of
these methods to language revitalization.

The Edwards Mahican corpus (144 records) is unique among the SNEA sources in
being recorded by a fluent second-language speaker of the language (Edwards
learned Mahican as a child from playmates). This makes Edwards the most
phonetically reliable colonial recorder in the corpus---but the \textit{gh}
grapheme mystery and \textit{uh} reduction patterns documented in this card
show that even a fluent L2 speaker's orthography is filtered through English
phonological categories, the same problem that both Cunningham and Kopris
address in their respective language families.

\subsection{References}

\begin{itemize}
  \item Brasser, T.J.\ (1978). ``Mahican.'' In \emph{Handbook of North
    American Indians}, Vol.\ 15, \emph{Northeast}, pp.\ 198--212. Washington:
    Smithsonian Institution.
  \item Edwards, Jonathan. (1788). \emph{Observations on the Language of the
    Muhhekaneew Indians}. New Haven: Josiah Meigs. [Reprinted 1823 with notes
    by John Pickering, in Massachusetts Historical Collections.]
  \item Goddard, Ives. (1979). ``Munsee Delaware.'' In \emph{Handbook of
    North American Indians}, Vol.\ 15, pp.\ 228--234.
  \item Goddard, Ives. (1982). \emph{The Historical Phonology of Munsee.}
    International Journal of American Linguistics 48:16--48.
  \item Goddard, Ives. (2007--2008). Corrections to Pentland's Mahican
    historical phonology. Papers of the Algonquian Conference.
  \item Goddard, Ives. (2010). Handout on Mahican phoneme inventory.
    Conference on Language Reconstruction, Ann Arbor, MI.
  \item Masthay, Carl (ed.). (1991). \emph{Schmick's Mahican Dictionary,
    with a Mahican Historical Phonology by David Pentland}. Memoirs of the
    American Philosophical Society, Vol.\ 197. Philadelphia.
  \item Pentland, David. (1991). ``A Mahican Historical Phonology.'' In
    Masthay 1991.
  \item Warne, Janet. n.d. ``Time-Depth in Mahican Diachronic Phonology:
    Evidence from the Schmick Manuscript.'' McGill University.
  \item \emph{Important Mahican Vowel Information} (PDF). 2024.
    munseedelaware.com. Consultation with Carl Masthay and Ives Goddard.
  \item Zeisberger, David. (1772). \emph{Zeisberger's Delaware Vocabulary of
    1772}.
  \item Zeisberger, David. (1827). \emph{Grammar of the Language of the
    Lenni Lenape or Delaware Indians}. Philadelphia.
\end{itemize}


% ============================================================
% PART II — Thematic Synthesis
% ============================================================
\part{Thematic Synthesis}

% Preaspiration Across Sources — Thematic Synthesis
% !TEX program = xelatex
% !TEX encoding = UTF-8 Unicode

\section{Preaspiration Across Sources}
\label{sec:preaspiration}

Preaspiration --- the occurrence of /h/ before stops (/hp/, /ht/, /hk/) --- is
one of the most debated features in SNEA historical phonology. The eight
sources in this study show dramatically different rates of preaspiration
reporting, which tells us more about the recorders than about the language.
English has no phonemic preaspiration, and English-speaking recorders
routinely failed to hear this feature. The result is a systematic gap in the
documentary record that must be filled through comparative and morphotactic
evidence.

\subsection{Cross-Source Comparison of Preaspiration Reporting}

Table~\ref{tab:preaspiration-comparison} summarizes how each source treats
preaspiration. The contrast between Edwards (who writes it consistently) and
Williams (who almost never writes it) is the most striking.

\begin{table}[h]
\centering
\caption{Preaspiration reporting across SNEA sources}
\label{tab:preaspiration-comparison}
\begin{tabular}{llll}
\toprule
\textbf{Source} & \textbf{Year} & \textbf{/h/ written?} & \textbf{Reliability} \\
\midrule
Edwards (Mahican) & 1787 & YES --- consistent & Most reliable \\
Prince-Speck (Mohegan-Pequot) & 1904 & YES --- explicit /hC/ clusters & High \\
Eliot (Wampanoag) & 1666 & Partial --- ~15\% of expected positions & Moderate \\
Winslow (Wampanoag) & 1624 & Sometimes --- \textit{kiehtan}, \textit{hobbamock} & Partial \\
Williams (Narragansett) & 1643 & Rarely --- 0 of 2,134 records & Unreliable \\
Wood (Wampanoag) & 1634 & Never & Unreliable \\
Trumbull (Wampanoag) & 1903 & Inherits Eliot's inconsistency & Moderate \\
Fielding (Mohegan) & 2013 & Not phonemic (modern loss) & N/A \\
\bottomrule
\end{tabular}
\end{table}

\subsubsection{Edwards (1787) --- The Gold Standard}

Jonathan Edwards Jr.\ is the only colonial source who can be considered
reliable for /h/ detection. He was a fluent L2 speaker of Mahican, having
learned the language as a child from playmates in Stockbridge, Massachusetts.
His transcriptions consistently write /h/ in positions where PA
reconstruction predicts it: \textit{hkeesque} 'eye', \textit{hpoon} 'winter',
\textit{keesogh} 'sun' (with /h/ in the final fricative). Edwards' training
in Hebrew and Greek, combined with his childhood fluency, gave him
phonological categories that included /h/ as a salient segment.

\subsubsection{Prince-Speck (1904) --- The Trained Linguist}

J.\ Dyneley Prince, trained in the Neogrammarian tradition, explicitly writes
/hC/ clusters where all earlier SNEA sources miss them. The spelling
\textit{bahkeder} 'bitter' shows /hk/ written; Williams or Eliot would write
only \textit{k}. Prince's Germanic training gave him an ear for aspiration
contrasts: German distinguishes aspirated /p\textsuperscript{h}, t\textsuperscript{h},
k\textsuperscript{h}/ in certain positions from unaspirated /p, t, k/ in
clusters, and the \textit{Dehnungs-h} (lengthening h) convention made him
conscious of h-like elements. Additionally, Mohegan-Pequot may have preserved
preaspiration longer than other SNEA languages, or Speck's field notes may
have captured it from Fidelia Fielding's careful speech.

\subsubsection{Eliot (1666) --- Partial and Inconsistent}

John Eliot writes \textit{hk}, \textit{hp}, \textit{ht} clusters in
approximately 15\% of expected positions. His \textit{Indian Grammar Begun}
(1666) distinguishes aspirated from unaspirated stops using English
voiced/voiceless pairs (b/p, d/t, k/g), but he does not apply this
consistently in the Bible translation. The apostrophe in Eliot's orthography
may also encode preaspiration or glottal stop --- the function is ambiguous
without comparative evidence.

\subsubsection{Winslow (1624) --- Partial but Notable}

Edward Winslow writes /h/ more consistently than any other early Wampanoag
source except Edwards. His \textit{kiehtan} (the supreme God) shows the /ht/
cluster, and \textit{hobbamock} (devil) shows word-initial /h/ where Wood's
cognate \textit{abamocho} omits it. This relatively good /h/ transcription
may reflect: (1) his intermediaries (Tisquantum, Hobbamock) had clearer /h/
articulation, (2) the Pokanoket dialect had more salient /h/ than Natick, or
(3) Winslow's ``ignorance'' of the language made him more phonetically
accurate (less normalization bias).

\subsubsection{Williams (1643) --- The Invisible /h/}

Roger Williams almost never wrote preaspirated /h/ before stops. This is the
single most important finding for the Narragansett source. In Narragansett
(and other SNEA languages), /h/ commonly precedes /p/, /t/, /k/ in
word-medial position. English speakers perceive these as single stops /k/,
/p/, /t/ because preaspiration is not a phonemic distinction in English.
Williams wrote \textit{k} where the actual pronunciation was /hk/
(preaspirated k).

The database contains only 3 records with phonetic (\textbackslash ph) fields
out of 2,134 total. All three show the same pattern:

\begin{table}[h]
\centering
\caption{The three \textbackslash ph entries in the Williams corpus}
\label{tab:williams-ph}
\begin{tabular}{llll}
\toprule
\textbf{ID} & \textbf{Williams spelling} & \textbf{\textbackslash ph (phonetic)} & \textbf{Pattern} \\
\midrule
3834 & Aaunchemókaw & ąːčimoːhkaːw & \textit{k} $\rightarrow$ /hk/ \\
4286 & Anakáusichick & anoːhkaːsiːčiːk & \textit{k} $\rightarrow$ /hk/ \\
4285 & Anakáusu & anoːhkaːsəw & \textit{k} $\rightarrow$ /hk/ \\
\bottomrule
\end{tabular}
\end{table}

All three show Williams writing \textit{k} where the reconstruction has /hk/.
The /h/ is invisible to English ears. If this pattern generalizes --- and
comparative evidence from Massachusett and Mohegan-Pequot suggests it does ---
then more than 50\% of Williams' stops may be missing a preceding /h/.

\subsection{Cowan (1973): Indirect Evidence for Preaspiration}

William Cowan's landmark study ``The Development of PA *-k- and *-hk- in
Narragansett'' (1973) demonstrated that PA *-k- (simple velar) and *-hk-
(preaspirated velar) merged phonetically in Narragansett as /hk/ medially,
but can sometimes be distinguished through the behavior of preceding vowels.
PA *-k- shortens a preceding long vowel; PA *-hk- does not. This provides an
indirect method for reconstructing preaspiration from vowel length patterns
--- even though Williams never writes the /h/.

Cowan's finding means that vowel length in Williams' orthography is not just
a vowel quality marker; it carries information about the following consonant's
preaspiration status. A long vowel before a stop suggests the stop was
simple (PA *-k-); a short vowel before a stop suggests the stop was
preaspirated (PA *-hk-). This interaction between vowel length and
preaspiration is a critical input to the conversion workflow.

\subsection{Implications for Conversion Workflows}

The pattern is clear: preaspiration was a genuine phonological feature of
SNEA languages, but English-speaking recorders routinely failed to hear it.
The recorders who wrote it most consistently (Edwards, Prince-Speck) were
either fluent L2 speakers or trained linguists with non-English phonological
categories. This has direct implications for the conversion workflows:

\begin{enumerate}
  \item \textbf{Preaspiration must be inferred, not read.} For Williams, Wood,
    and most of Eliot, preaspiration /h/ must be inferred from morphotactic
    and comparative evidence, not from the orthography.

  \item \textbf{Morphotactic position is the primary cue.} Preaspiration is
    most common before stops in the coda of stressed syllables. A
    stop preceded by a short vowel in a stressed syllable is a candidate for
    preaspiration.

  \item \textbf{Comparative evidence from Edwards is the best guide.}
    Edwards' consistent /h/ marking in Mahican provides a model for
    predicting preaspiration in sources where it is unwritten. Where Edwards
    writes /h/ and Williams omits it, the preaspiration is probable.

  \item \textbf{Cowan's vowel-length diagnostic.} For Narragansett
    specifically, vowel length patterns can serve as an indirect cue: a short
    vowel before a stop suggests preaspiration (PA *-hk-), while a long vowel
    suggests a simple stop (PA *-k-).

  \item \textbf{PA reconstructions provide the baseline.} Where PA
    reconstruction shows */h/ in a corresponding position, preaspiration is
    expected in the daughter language unless a specific sound change removed
    it.

  \item \textbf{Prince-Speck as a key witness.} The Prince-Speck glossary
    (1904) is unique among the pre-modern SNEA sources in writing /hC/
    clusters explicitly. This makes it a key witness for preaspiration in
    SNEA historical phonology, on par with Mayhew and Bourne (Wampanoag) for
    the same feature.
\end{enumerate}

Without these inference strategies, the conversion pipeline would silently
drop a phonemic /h/ from every word where it was present but unwritten ---
producing a systematic underdifferentiation that would propagate through all
downstream phonological analysis.

% Colonial English Phonology as a Filter — Thematic Synthesis
% !TEX program = xelatex
% !TEX encoding = UTF-8 Unicode

\section{Colonial English Phonology as a Filter}
\label{sec:colonial-phonology}

Each recorder's English orthography is a lens through which Algonquian
phonemes were filtered. Understanding the recorder's native phonological
system is essential to recovering the target language's sound system. The
colonial recorders did not write what they heard --- they wrote what their
English phonological categories allowed them to perceive. This section
establishes the Early Modern English baseline for each recorder and
documents how their native systems systematically distorted their
transcriptions.

\subsection{Early Modern English Phonology Baseline (c.\ 1600--1620)}

The Great Vowel Shift (GVS) was complete by approximately 1600--1620 for
most English dialects. This means the colonial recorders' vowel system was
significantly different from both Middle English and Modern English.
Table~\ref{tab:eme-vowels} shows the vowel values that the earliest
recorders (Winslow, Wood, Williams) would have used.

\begin{table}[h]
\centering
\caption{Early Modern English vowel values (c.\ 1600--1620)}
\label{tab:eme-vowels}
\begin{tabular}{llll}
\toprule
\textbf{Spelling} & \textbf{ME value} & \textbf{c.\ 1600 value} & \textbf{Modern value} \\
\midrule
\textit{a} (short) & /a/ & /a/ & /æ/ \\
\textit{e} (short) & /ɛ/ & /ɛ/ & /ɛ/ \\
\textit{i} (short) & /ɪ/ & /ɪ/ & /ɪ/ \\
\textit{o} (short) & /ɔ/ & /ɔ/ & /ɒ/~/ɑ/ \\
\textit{u} (short) & /ʊ/ & /ʌ/ & /ʌ/ \\
\textit{a} (long) & /aː/ & /ɛː/~/eː/ & /eɪ/ \\
\textit{ee} (long) & /eː/ & /iː/ & /iː/ \\
\textit{ie} (long) & /ɛː/ & /iː/ & /iː/ \\
\textit{oo} (long) & /oː/ & /uː/ & /uː/ \\
\textit{ou} (long) & /uː/ & /aʊ/ & /aʊ/ \\
\textit{ea} (digraph) & /ɛː/ & /ɛː/~/eː/ & /iː/ \\
\textit{ai/ay} & /aɪ/ & /ɛː/~/eː/ & /eɪ/ \\
\bottomrule
\end{tabular}
\end{table}

The critical point for Algonquian transcription is that the \textit{ea}
digraph in 1620--1643 represented /ɛː/ or /eː/, NOT the modern /iː/. When
Williams writes \textit{wean} for a Narragansett word, the \textit{ea}
represents a mid front vowel, not a high front vowel. Similarly,
\textit{oo} in 1620 represented /oː/ (still raising toward /uː/), not the
modern /uː/. These differences matter for mapping colonial spellings to
Algonquian phonemes.

\subsection{Per-Recorder Linguistic Background}

Each recorder brought a different linguistic background to the task of
transcribing an unfamiliar language. Table~\ref{tab:recorder-background}
summarizes their training and its impact.

\begin{table}[h]
\centering
\caption{Recorder linguistic background and transcription impact}
\label{tab:recorder-background}
\begin{tabular}{llll}
\toprule
\textbf{Recorder} & \textbf{Training} & \textbf{Key bias} & \textbf{Impact} \\
\midrule
Winslow (1624) & None (colonist) & English-hearer perception & Raw phonetic impression \\
Wood (1634) & None (colonist) & English-hearer perception & b/p alternation, unstable vowels \\
Williams (1643) & Theologian, self-taught linguist & English categories + stress awareness & Diacritics for stress, no /h/ \\
Eliot (1666) & Missionary, Hebrew/Greek training & Systematic but English-filtered & Aspiration via voiced/voiceless \\
Edwards (1787) & Theologian, fluent L2 speaker & Hebrew/Greek + childhood fluency & Most reliable /h/ detection \\
Prince (1904) & Neogrammarian linguist & German orthographic conventions & \textit{tsh} for /tʃ/, \textit{Dehnungs-h} \\
Speck (1904) & Boasian anthropologist & Ethnographic, minimal linguistic training & Two-stage recording chain \\
Fielding (2013) & Modern linguist (MIT) & Phonemic orthography & 1:1 phoneme-to-grapheme \\
\bottomrule
\end{tabular}
\end{table}

\subsubsection{Winslow (1624) --- The Honest Amateur}

Edward Winslow was unusually honest about his linguistic limitations,
describing the Massachusett language as ``very copious, large, and
difficult.'' His transcriptions are raw phonetic impressions --- he was not
normalizing or regularizing what he heard, unlike Eliot (who developed a
systematic orthography) or Williams (who, though linguistically gifted,
worked through an English lens). Winslow's data passed through two
intermediaries: Tisquantum (Squanto), a Patuxet Wampanoag who learned
English in London and may have experienced language attrition, and
Hobbamock, a Pokanoket \textit{pniese} (military champion) with limited
English. This intermediary filter means every transcription is at least two
steps removed from the speaker.

\subsubsection{Wood (1634) --- The Pre-Orthographic Recorder}

William Wood's orthography is pre-Eliot --- the Massachusett writing system
was still decades away. His transcriptions reflect raw English-hearer
perception without regularization. The most diagnostic feature is his b/p
alternation: Wood writes \textit{abbona} (five) and \textit{abamocho} (devil)
with \textit{b} where the Eliot/Trumbull orthography uses \textit{p}. This
shows Wood perceived the unaspirated Wampanoag /p/ as closer to English /b/
--- direct evidence of the voicing neutralization that characterizes SNEA
stop systems. Wood's triple spelling for a single word
(\textit{abbomocho}/\textit{abomocho}/\textit{abamacho}) is the strongest
evidence in the entire SNEA corpus for colonial vowel instability.

\subsubsection{Williams (1643) --- The Linguistically Gifted Amateur}

Roger Williams was a theologian with no formal linguistic training, but his
orthography is the most phonetically sophisticated of the 17th-century
sources. His unique contribution is the systematic use of diacritics for
stress: acute (á) for primary stress, grave (à) for low/falling tone, and
circumflex (â) for rising-falling tone. This is a major advantage over
Eliot's Massachusett Bible (1663), which uses no accent marks, making stress
reconstruction for Massachusett far more difficult. However, Williams'
English phonological filter still produced systematic gaps: he almost never
wrote preaspirated /h/, and his five-vowel system (a, e, i, o, u) must
represent 8--10 Narragansett vowel phonemes.

\subsubsection{Eliot (1666) --- The Systematic Missionary}

John Eliot developed the first standardized writing system for an Algonquian
language, working with Massachusett collaborators including Job Nesutan and
John Sassamon. His orthography was phonemically oriented for its time --- he
attempted to represent the sounds of Massachusett systematically, not just
through English ear-approximations. His key innovation was distinguishing
aspirated from unaspirated stops using English voiced/voiceless pairs (b/p,
d/t, k/g). However, this system was still filtered through English
phonological categories: Eliot's \textit{b, d, g} represent unaspirated lenis
stops (not voiced), and his \textit{p, t, k} represent aspirated fortis
stops. The distinction is phonetically real in Massachusett but mapped onto
an English voicing contrast that does not align perfectly.

\subsubsection{Edwards (1787) --- The Fluent L2 Speaker}

Jonathan Edwards Jr.\ is unique among the colonial recorders in being a
fluent L2 speaker of Mahican, having learned the language as a child in
Stockbridge, Massachusetts, where Mahican families were the majority
population. His training in Hebrew and Greek gave him explicit phonological
categories, and his childhood fluency gave him native-like perception. The
result is the most reliable colonial transcription in the entire SNEA corpus.
Edwards consistently writes /h/ where PA reconstruction predicts it, and his
vowel transcriptions are more stable than any other colonial source. His
\textit{Observations on the Language of the Muhhekaneew Indians} (1788)
includes comparative glossaries with Shawnee and Ojibwa, demonstrating a
comparative method that was far ahead of its time.

\subsubsection{Prince-Speck (1904) --- The Two-Stage Chain}

J.\ Dyneley Prince and Frank Speck represent a two-stage recording chain
(fieldworker $\rightarrow$ desk transcriber) that introduces compounding
transcription errors. Speck, a Boasian anthropologist with minimal linguistic
training, collected field notes from Fidelia Fielding, the last fluent
speaker of Mohegan-Pequot. Prince, a Neogrammarian linguist who never heard
the language spoken, converted Speck's notes into the published orthography.
Prince's Germanic training introduced specific biases: \textit{tsh} for /tʃ/
(following German \textit{tsch}), \textit{z} for /ts/ (German \textit{z}),
and \textit{Dehnungs-h} (lengthening h) for vowel length. The glossary also
preserves some of Fielding's own English-based spellings for ``sentimental
reasons'' (Prince \& Speck 1904:19), creating a three-way orthographic
competition between English colonial conventions, Fielding's vernacular
spelling, and Prince's German linguistic conventions.

\subsection{The Kopris (2014) Framework}

Craig Kopris' paper ``Les sons wyandots perçus par des oreilles étrangères''
(Wyandot sounds perceived by foreign ears, 2014) provides the exact
theoretical framework for the colonial recorder problem. Kopris demonstrates
that European recorders systematically misperceived and mis-transcribed
indigenous phonemes because their native phonological systems lacked the
necessary categories. His key insight is that the recorder's L1 is not a
source of random error but a systematic filter: the errors are predictable
given the recorder's native phoneme inventory.

Applied to the SNEA corpus, the Kopris framework predicts:

\begin{enumerate}
  \item \textbf{English recorders will merge Algonquian stop series.} English
    has a two-way laryngeal contrast (voiced vs.\ voiceless aspirated);
    Algonquian languages have a different contrast (plain vs.\ preaspirated,
    or fortis vs.\ lenis). The mapping is systematically distorted.

  \item \textbf{Vowel quality will be underdifferentiated.} English has more
    vowel qualities than the five-letter Latin alphabet can represent, but
    colonial recorders still used only five vowel letters. The mapping from
    Algonquian vowels to English letters is systematically ambiguous.

  \item \textbf{Stress and length will be confused.} English stress is
    dynamic (loudness-based); Algonquian stress may be tonal or length-based.
    Recorders who mark stress (Williams) may be indirectly marking length,
    and vice versa.

  \item \textbf{Preaspiration will be invisible.} English has no phonemic
    preaspiration, so English-speaking recorders will systematically fail to
    hear it --- unless they have training (Prince) or L2 fluency (Edwards)
    that provides the necessary category.
\end{enumerate}

The Kopris framework transforms the colonial recorder problem from a
liability into an asset: once we know the recorder's L1 system, we can
predict the systematic distortions and reverse them. This is the theoretical
foundation for the cross-source calibration methodology described in
Section~\ref{sec:calibration}.

% Vowel Ambiguity Patterns — Thematic Synthesis
% !TEX program = xelatex
% !TEX encoding = UTF-8 Unicode

\section{Vowel Ambiguity Patterns}
\label{sec:vowel-ambiguity}

Every colonial source in this study exhibits vowel underdetermination, but the
patterns differ systematically. Understanding what is systematic vs.\ random
in the vowel ambiguity is essential for designing conversion workflows. The
core problem is that each recorder used a five-letter Latin alphabet (a, e, i,
o, u) to represent an Algonquian vowel system of 8--10 phonemes, including
length distinctions and nasalization that the Latin alphabet was not designed
to encode.

\subsection{Per-Source Vowel Ambiguity}

Table~\ref{tab:vowel-ambiguity} shows how many Algonquian phonemes each
grapheme maps to in each source. The number of possible mappings per
grapheme ranges from 1 (for the most consistent digraphs) to 4 (for the
most ambiguous single letters).

\begin{table}[h]
\centering
\caption{Per-source vowel ambiguity: number of Algonquian phonemes per grapheme}
\label{tab:vowel-ambiguity}
\small
\begin{tabular}{lcccccc}
\toprule
\textbf{Grapheme} & \textbf{Williams} & \textbf{Eliot} & \textbf{Wood} & \textbf{Winslow} & \textbf{Edwards} & \textbf{Fielding} \\
\midrule
\textit{a} & 3 (/ə, ɔː, aː/) & 2 (/a, ə/) & 2 (/a, ə/) & 2 (/a, ə/) & 2 (/a, aː/) & 1 (/a/) \\
\textit{e} & 4 (/iː, ɛ, ə, ∅/) & 3 (/iː, ɛ, ə/) & 2 (/ɛ, i/) & 2 (/ɛ, ə/) & 2 (/iː, i/) & 1 (/iː/) \\
\textit{i} & 3 (/ə, iː, ɪ/) & 2 (/iː, ɪ/) & 2 (/ɪ, i/) & 2 (/ɪ, i/) & 1 (/i/) & 1 (/i/) \\
\textit{o} & 4 (/ə, aː, uː, ɔː/) & 3 (/uː, ɔ, ə/) & 3 (/ɔ, o, ə/) & 2 (/ɔ, ə/) & 2 (/o, u/) & --- \\
\textit{u} & 2 (/ə, a/) & 2 (/ə, ʌ/) & 2 (/ə, ʌ/) & 1 (/ə/) & 2 (/ə, o/) & 1 (/u/) \\
\textit{ee} & 1 (/iː/) & 1 (/iː/) & 1 (/iː/) & 1 (/iː/) & 1 (/iː/) & --- \\
\textit{oo} & 1 (/uː/) & 1 (/uː/) & 1 (/uː/) & 1 (/uː/) & 1 (/uː/) & --- \\
\midrule
\textbf{Total phonemes} & 8--10 & 8--10 & 6--8 & 6--7 & 7--8 & 6 \\
\bottomrule
\end{tabular}
\end{table}

The table reveals a clear hierarchy of ambiguity. The most reliable
transcriptions are the English digraphs \textit{ee} and \textit{oo}, which
consistently map to /iː/ and /uː/ respectively across all sources. The most
ambiguous are Williams' \textit{e} (4 possible values) and \textit{o} (4
possible values), reflecting the fact that Narragansett had a larger vowel
inventory than the five Latin letters could represent.

\subsubsection{Williams' Vowel System in Detail}

Williams' orthography is the most ambiguous because Narragansett had the
largest vowel inventory among the SNEA languages documented. His simple
vowel mappings (from O'Brien's Chart B) illustrate the problem:

\begin{itemize}
  \item \textit{a} $\rightarrow$ /ə/, /ɔː/, /aː/ --- three-way ambiguity
    covering schwa, open-o, and long a
  \item \textit{ah} $\rightarrow$ /aː/, /ɔː/ --- long variants
  \item \textit{e} $\rightarrow$ /iː/, /ɛ/, /ə/, ∅ --- four-way ambiguity;
    silent word-final
  \item \textit{ê} $\rightarrow$ /iː/ --- circumflex = long i (consistent)
  \item \textit{i} $\rightarrow$ /ə/, /iː/, /ɪ/ --- three-way; often = schwa
    in unstressed
  \item \textit{o} $\rightarrow$ /ə/, /aː/, /uː/; after \textit{w}: /aː/,
    /ɔː/ --- broad range
  \item \textit{oo}, \textit{ô} $\rightarrow$ /uː/ --- consistent long u
  \item \textit{ie}, \textit{ee} $\rightarrow$ /iː/ --- consistent
  \item \textit{ei} $\rightarrow$ /iː/, /ə/ --- two-way
  \item \textit{ea} $\rightarrow$ /iː/, /ɛə/, /aː/ --- three-way
\end{itemize}

\subsubsection{Eliot's Vowel System}

Eliot's orthography is more systematic but still underdetermined. His five
vowel letters (a, e, i, o, u) must represent 8--10 Massachusett phonemes.
The key disambiguating feature is that Eliot sometimes uses vowel doubling
to mark length: \textit{ee} = /iː/, \textit{oo} = /uː/, \textit{aa} = /aː/.
However, this doubling is not applied consistently across the Bible
translation. Eliot's distinctive ⟨ꝏ⟩ (oo ligature, often mis-encoded as ∞)
represents long /uː/, contrasting with \textit{oo} which may represent a
different vowel quality or length.

\subsubsection{Wood's Vowel Instability}

Wood's vowel system is the least stable. His triple spelling for 'devil'
(\textit{abbomocho}/\textit{abomocho}/\textit{abamacho}) shows that the
same recorder, recording the same word, produced three different vowel
transcriptions. The stable consonants (\textit{m}, \textit{ch}) are reliable,
while the unstable vowels (\textit{o}/\textit{a}) indicate reduced or
unstressed vowels that English ears struggled to categorize. This pattern is
diagnostic: when a word appears with multiple spellings in Wood, the stable
consonants are reliable while the unstable vowels indicate schwa or
unstressed reduction.

\subsubsection{Edwards' Vowel System --- The /iː/ Correction}

Recent work (Goddard 2007--2008, confirmed by Masthay 2024) has established
that Edwards' \textit{e} frequently represents /iː/, NOT /eː/ or /ɛ/. This
is the most important correction to the interpretation of Edwards'
orthography. Evidence includes:

\begin{itemize}
  \item \textit{Muhhekaneew} --- the second vowel is /iː/, not /eː/
  \item \textit{Neesoh} 'two' --- cf.\ PA *niːšwi, Munsee /niːšə/
  \item \textit{Seepoo} 'river' --- cf.\ PA *siːpoːwi, Munsee /siːpəw/
  \item Goddard's 2010 phoneme inventory for Mahican lists NO /eː/ phoneme
\end{itemize}

The default mapping for Edwards' \textit{e} should be /i(ː)/ unless cognate
evidence specifically indicates a different value.

\subsection{Systematic vs.\ Random Variation}

A critical distinction for the conversion workflow is whether vowel variation
is systematic (reflecting real phonological or dialect differences) or random
(reflecting the recorder's inability to consistently categorize the same
sound).

\subsubsection{Random Variation: Wood's Triple Spelling}

Wood's \textit{abbomocho}/\textit{abomocho}/\textit{abamacho} for 'devil' is
the clearest example of random variation. The same recorder, in the same
publication, produces three different transcriptions of the same word. The
vowel in the second syllable alternates between \textit{o} and \textit{a},
suggesting a reduced vowel that Wood could not consistently categorize. This
type of variation can be averaged out: the stable reading across multiple
attestations is likely the correct vowel quality.

\subsubsection{Systematic Variation: Williams vs.\ Winslow}

The consistent difference between Williams' \textit{wétu} and Winslow's
\textit{witeo} for 'house' represents systematic variation that may reflect
real dialect differences. Williams records Narragansett (Rhode Island), while
Winslow records Pokanoket Wampanoag (Plymouth area). The \textit{e}/\textit{i}
alternation in the first syllable and the final \textit{u}/\textit{o}
alternation are consistent across multiple attestations, suggesting a real
phonological difference between the dialects rather than random transcription
error.

\subsubsection{Diagnostic Criteria}

To distinguish systematic from random variation, we apply the following
criteria:

\begin{enumerate}
  \item \textbf{Consistency across attestations:} If the same source
    consistently writes the same vowel for the same word across multiple
    occurrences, the variation is likely systematic. If the vowel varies
    unpredictably, it is likely random.

  \item \textbf{Cross-source agreement:} If two independent sources agree on
    a vowel value, the variation is likely systematic. If they disagree, the
    disagreement may reflect dialect differences (systematic) or recorder
    error (random).

  \item \textbf{PA correspondence:} If the PA reconstruction predicts a
    specific vowel reflex, the variation can be evaluated against the
    prediction. Variation that matches the prediction is systematic;
    variation that does not is likely random.

  \item \textbf{Phonological environment:} Variation in unstressed syllables
    is more likely to be random (schwa reduction). Variation in stressed
    syllables is more likely to be systematic.
\end{enumerate}

\subsection{Schwa Reduction Patterns}

Schwa (/ə/) is the most variably transcribed vowel across all sources.
Colonial English recorders had no consistent way to represent schwa, so it
appears as \textit{a}, \textit{e}, \textit{i}, \textit{o}, or \textit{u}
depending on context and recorder. The following patterns emerge:

\begin{itemize}
  \item \textbf{Williams:} Schwa is most often written as \textit{i} or
    \textit{u} in unstressed syllables. Word-final schwa is often written as
    \textit{e} (which may be silent).

  \item \textbf{Eliot:} Schwa is most often written as \textit{u} in
    unstressed syllables. Eliot's grammar describes \textit{u} as the
    ``obscure'' vowel, corresponding to schwa.

  \item \textbf{Wood:} Schwa alternates unpredictably between \textit{a},
    \textit{e}, \textit{i}, \textit{o}, and \textit{u}. The triple spelling
    for 'devil' shows \textit{o} and \textit{a} alternating for what is
    likely a schwa.

  \item \textbf{Winslow:} Schwa is most often written as \textit{e} in final
    position (e.g., \textit{ahhe}, \textit{witeo}).

  \item \textbf{Edwards:} Schwa is most often written as \textit{u} in
    unstressed syllables. The sequence \textit{uh} may represent /əh/ or
    breathy schwa in a checked syllable.

  \item \textbf{Prince-Speck:} Schwa is written as \textit{'} (apostrophe)
    in their phonetic key, but in the orthographic forms it appears as
    \textit{a}, \textit{e}, or \textit{u}.
\end{itemize}

The schwa reduction pattern is so consistent across sources that it can be
used as a diagnostic: when a vowel in an unstressed syllable varies across
attestations, the default interpretation should be schva.

\subsection{Nasalization Ambiguity}

Nasalization is one of the most poorly documented features in the colonial
sources. Proto-Algonquian *aː and *oː developed nasalized reflexes in SNEA
languages before certain consonants, but no colonial source marks
nasalization consistently.

\begin{itemize}
  \item \textbf{Williams:} Nasalization is implied but not marked. He uses
    nasal coda consonant digraphs (\textit{an}, \textit{aun}, \textit{om},
    \textit{on}) to imply nasal vowels, but these same sequences also
    represent oral /an/, /aun/, /om/, /on/. Disambiguating requires
    checking the Massachusett cognate (Eliot marks nasal /ô/ explicitly) or
    the PA reconstruction.

  \item \textbf{Eliot:} The only colonial source that marks nasalization
    explicitly, using \textit{ô}, \textit{û}, or an overline. However, this
    marking is not applied consistently across the Bible translation.

  \item \textbf{Wood, Winslow:} Nasalization is entirely unmarked.

  \item \textbf{Edwards:} Nasalization is not marked. PA *aː developed a
    nasal quality in Mahican (marked by Schmick with a tilde), but Edwards
    does not indicate this.

  \item \textbf{Prince-Speck:} Nasalization is not systematically marked.
    May occasionally indicate it through vowel quality choices (e.g.,
    \textit{o} for nasalized /ã/).

  \item \textbf{Fielding:} Modern Mohegan has lost nasal vowels, so the
    modern source provides no help for this feature.
\end{itemize}

The nasalization ambiguity is consequential because it is phonemic in SNEA
languages. In Narragansett, \textit{n now} (yes) and \textit{núnow} (no) are
distinguished only by nasality, and Williams collapses this distinction
orthographically in related forms. The conversion workflow must infer
nasalization from comparative evidence rather than reading it from the
orthography.

% Cross-Source Calibration Methodology — Thematic Synthesis
% !TEX program = xelatex
% !TEX encoding = UTF-8 Unicode

\section{Cross-Source Calibration Methodology}
\label{sec:calibration}

One of the key findings of this project is that data from one source can
resolve ambiguity in another. This cross-source calibration is the most
powerful tool in the conversion toolkit. Because each recorder had different
strengths and weaknesses --- Edwards was reliable for /h/ but not for vowel
quality, Williams was reliable for stress but not for preaspiration, Eliot
was systematic but English-filtered --- the sources complement each other.
By calibrating across sources, we can recover phonological information that
no single source provides.

\subsection{The Calibration Pipeline}

The calibration pipeline proceeds from most reliable to least reliable, with
each source used to calibrate the next. The ordering is determined by the
recorder's linguistic training, fluency, and orthographic system:

\begin{enumerate}
  \item \textbf{Fielding (2013)} --- Modern phonemic orthography with 1:1
    phoneme-to-grapheme correspondence. The most phonologically reliable
    source, but limited to Mohegan-Pequot and representing the modern
    revived language.

  \item \textbf{Trumbull/Eliot (1903/1666)} --- Systematic missionary
    orthography with the largest corpus (4,491 records). Eliot's grammar
    provides explicit phonological descriptions, and the large corpus enables
    statistical approaches.

  \item \textbf{Edwards (1787)} --- Fluent L2 speaker with the most reliable
    /h/ detection. Small corpus (144 records) but high per-record
    reliability.

  \item \textbf{Prince-Speck (1904)} --- Trained linguist with explicit /hC/
    cluster notation. Medium corpus (446 records) but two-stage recording
    chain introduces compounding errors.

  \item \textbf{Williams (1643)} --- Linguistically gifted amateur with
    unique stress diacritics. Large corpus (2,134 records) but high vowel
    ambiguity and no /h/ detection.

  \item \textbf{Wood/Winslow (1634/1624)} --- Raw English perception with no
    linguistic training. Small corpora (339 and 24 records) but valuable as
    the earliest attestations.
\end{enumerate}

\subsection{Cross-Source Disambiguation: Specific Examples}

\subsubsection{Massachusett Disambiguates Narragansett Vowels}

Williams' vowel ambiguity in Narragansett is partially resolved by comparing
against Eliot's Massachusett cognates, which have a more systematic
orthography. For example, Williams' \textit{a} can represent /ə/, /ɔː/, or
/aː/. When the Massachusett cognate shows \textit{â} (Eliot's long a marker),
the Narragansett vowel is likely /aː/. When the Massachusett cognate shows
\textit{u} (Eliot's schwa marker), the Narragansett vowel is likely /ə/.

The cognate pair for 'dog' illustrates this: Williams writes \textit{ayim},
Eliot writes \textit{arum}. The PA reconstruction *aɬemwa confirms the
vowel quality and shows the /ɬ/ $\rightarrow$ /j/ (Narragansett) vs.\ /ɬ/
$\rightarrow$ /r/ (Massachusett) reflex difference. The cross-source
comparison resolves both the vowel quality (first vowel /a/) and the
consonant reflex.

\subsubsection{Edwards Disambiguates Preaspiration}

Edwards' consistent /h/ marking in Mahican provides a model for predicting
preaspiration in Williams and Wood, where it is unwritten. When Edwards
writes \textit{hkeesque} 'eye' with initial /hk/, and the PA reconstruction
shows */h/ in a corresponding position, we can infer that the cognate
Narragansett form (which Williams writes without /h/) also had preaspiration.
This inference is strengthened when the morphotactic position matches: a stop
in the coda of a stressed syllable is a candidate for preaspiration.

\subsubsection{Prince-Speck Confirms Fielding}

Where Prince-Speck (1904) and Fielding (2013) agree on a Mohegan-Pequot
form, the confidence is high. For example, both record \textit{wetu} for
'house', showing remarkable lexical stability across 110 years. Where they
disagree, the circularity risk must be assessed: Fielding's dictionary is
partially reconstructive, and some entries may derive from Prince-Speck
itself. Using Fielding to ``verify'' Prince-Speck for such entries would be
circular. The conversion workflow must track which Fielding entries derive
from Prince-Speck (self-source) vs.\ from independent sources.

\subsubsection{Wood and Winslow Corroborate Each Other}

Wood (339 records) and Winslow (24 records) are the two earliest Wampanoag
sources, separated by only 10 years. Where they agree on a form, the
reading is likely correct. For example, both record \textit{matta} for 'no',
and both record \textit{sachim} for 'chief'. Where they disagree --- as with
Wood's \textit{abamocho} vs.\ Winslow's \textit{hobbamock} for 'devil' ---
the disagreement may reflect real dialect differences (Wood records
Massachusetts Bay, Winslow records Pokanoket) or different lexical roots
entirely.

\subsection{The Cognate Set Approach}

The most powerful calibration tool is the cognate set: a word with the same
meaning attested in multiple SNEA languages, with PA reconstruction as the
anchor. Table~\ref{tab:cognate-set} shows a representative cognate set for
basic vocabulary items.

\begin{table}[h]
\centering
\caption{Cross-source cognate set for basic vocabulary}
\label{tab:cognate-set}
\small
\begin{tabular}{llllll}
\toprule
\textbf{Gloss} & \textbf{Williams} & \textbf{Eliot} & \textbf{Edwards} & \textbf{Fielding} & \textbf{PA} \\
 & \textbf{(NRG)} & \textbf{(WMP)} & \textbf{(MAH)} & \textbf{(MPQ)} & \\
\midrule
'house' & wétu & wetu & --- & wetu & *wiːkiwaːmi \\
'water' & nipp & nippe & --- & nip & *nepyi \\
'dog' & ayim & arum & --- & ayum & *aɬemwa \\
'man' & weenawash & wénawaš & --- & --- & --- \\
'fire' & apeisuwau & apeisuwâ & --- & --- & --- \\
'I' & neen & neen & neen & neen & *niːra \\
'two' & nees & nees & neesoh & nees & *niːšwi \\
'no' & matta & matta & --- & matta & *mati- \\
'woman' & squa & squa & --- & squa & *eːθkweːwa \\
'five' & --- & pàpàna & nannah & --- & *nyaːlanwi \\
\bottomrule
\end{tabular}
\end{table}

The cognate set reveals several patterns:

\begin{enumerate}
  \item \textbf{Lexical stability:} Basic vocabulary items like 'house',
    'water', 'I', and 'two' are nearly identical across all SNEA languages
    and across 260+ years (Williams 1643 to Fielding 2013). This confirms
    the reliability of the core vocabulary for cross-source calibration.

  \item \textbf{Regular sound correspondences:} The PA *r reflex distinguishes
    N-dialects (Massachusett \textit{neen} from *niːra) from Y-dialects
    (Mohegan \textit{neen} with /n/ from a different PA source). The PA *ɬ
    reflex distinguishes Narragansett /j/ (\textit{ayim}) from Massachusett
    /r/ (\textit{arum}).

  \item \textbf{Cross-source vowel confirmation:} The vowel in 'house' is
    consistently \textit{e} across Williams, Eliot, and Fielding, confirming
    the /i/ quality (from PA *iː). The vowel in 'water' is consistently
    \textit{i} across all sources, confirming the /ɪ/ quality (from PA *e).

  \item \textbf{Edwards' Mahican as an outlier:} Edwards' \textit{neesoh}
    'two' shows the Mahican final vowel (\textit{-oh}) that is absent in the
    other SNEA languages, reflecting Mahican's transitional status between
    SNEA and Delaware.
\end{enumerate}

\subsection{How PA Reconstructions Anchor Vowel Quality}

Proto-Algonquian reconstructions (Bloomfield 1925, 1946; Aubin 1975;
Pentland 1979) provide the comparative evidence needed to choose among
ambiguous vowel mappings. The logic is:

\begin{enumerate}
  \item For a given colonial spelling, identify the set of possible phonemic
    interpretations (from the grapheme-to-phoneme mapping table).

  \item Find the PA reconstruction for the cognate set.

  \item Apply the known sound correspondences (PA $\rightarrow$ daughter
    language) to predict the expected reflex.

  \item Select the phonemic interpretation that matches the expected reflex.
\end{enumerate}

For example, Williams' \textit{wétu} 'house' has \textit{e} in the first
syllable. The possible phonemic interpretations of Williams' \textit{e} are
/iː/, /ɛ/, /ə/, or ∅. The PA reconstruction is *wiːkiwaːmi. The known
correspondence is PA *iː $\rightarrow$ Narragansett /iː/. Therefore, the
\textit{e} in \textit{wétu} represents /iː/, not /ɛ/ or /ə/. This is
confirmed by the consistent \textit{ee} spelling in Eliot's Massachusett
cognate \textit{wetu} (where \textit{ee} = /iː/).

This logic applies to every ambiguous vowel in the corpus. The PA
reconstruction is the anchor that resolves the ambiguity.

\subsection{Calibration Confidence Levels}

Each converted form is assigned a confidence level based on the strength of
the calibration evidence:

\begin{itemize}
  \item \textbf{[high]} --- Confirmed by cognate in a more reliable source or
    by PA reconstruction. No contradictory evidence.

  \item \textbf{[med]} --- Default rule applies, no contradicting cognate,
    but no confirming cognate either.

  \item \textbf{[low]} --- Guess; no cognate available, default rule may not
    apply.

  \item \textbf{[FLAG]} --- Contradictory evidence; manual review needed.
\end{itemize}

The calibration pipeline produces a confidence-tagged output that allows
downstream users to assess the reliability of each converted form. This
transparency is essential for both scholarly analysis and revitalization
applications, where the confidence level determines how the form can be
used.

% Dialect Classification Evidence — Thematic Synthesis
% !TEX program = xelatex
% !TEX encoding = UTF-8 Unicode

\section{Dialect Classification Evidence}
\label{sec:dialect}

The SNEA languages belong to different dialect groups within the Eastern
Algonquian family. The corpus data provides evidence for these
classifications, and understanding the dialect relationships is essential
for cross-source calibration. A form that looks like a transcription error
in one dialect may be a regular phonological reflex in another.

\subsection{The Four-Way Dialect Classification}

Southern New England Algonquian languages are traditionally classified by
the reflex of Proto-Algonquian *r (sometimes reconstructed as *l or *θ).
Four dialect types are recognized:

\begin{table}[h]
\centering
\caption{SNEA dialect classification by PA *r reflex}
\label{tab:dialect-classification}
\begin{tabular}{llll}
\toprule
\textbf{Dialect type} & \textbf{Reflex of PA *r} & \textbf{Languages} & \textbf{Sources} \\
\midrule
\textbf{N-dialect} & /n/ & Massachusett (Wampanoag), & Eliot, Trumbull, \\
 & & Narragansett, Nipmuc & Williams \\
\textbf{Y-dialect} & /j/ & Mohegan-Pequot, Montauk & Prince-Speck, Fielding \\
\textbf{R-dialect} & /r/ & Western Connecticut (Loup A) & (limited) \\
\textbf{L-dialect} & /l/ & Munsee Delaware, Unami & Zeisberger \\
\bottomrule
\end{tabular}
\end{table}

The PA *r reflex is the most diagnostic isogloss for SNEA dialect
classification, but it is not the only one. Costa (2007) identifies a
network of isoglosses that collectively define the SNEA dialect continuum.

\subsection{Costa (2007) Isoglosses}

David Costa's definitive study ``The Dialectology of Southern New England
Algonquian'' (2007) identifies the following key isoglosses that distinguish
the SNEA languages:

\begin{enumerate}
  \item \textbf{PA *r reflex:} /n/ (N-dialect) vs.\ /j/ (Y-dialect) vs.\
    /r/ (R-dialect) vs.\ /l/ (L-dialect). This is the primary
    classification criterion.

  \item \textbf{PA *θ reflex:} In N-dialects, PA *θ merges with /n/ (e.g.,
    Massachusett \textit{neen} 'I' from PA *niːra, where *r $\rightarrow$ /n/
    and *θ $\rightarrow$ /n/). In Y-dialects, PA *θ $\rightarrow$ /j/ in some
    environments.

  \item \textbf{PA *š reflex:} In Mahican, PA *š merges with /s/ (a
    distinguishing feature from the other SNEA languages, where /š/ is
    retained as a separate phoneme).

  \item \textbf{Word-final vowel deletion:} Mahican shows more advanced
    word-final vowel deletion than the other SNEA languages, reflecting its
    transitional status toward Delaware.

  \item \textbf{Preaspiration retention:} Mohegan-Pequot (as recorded by
    Prince-Speck) appears to have retained preaspiration longer than
    Narragansett or Massachusett, though this may be a recording artifact.

  \item \textbf{Stress patterns:} Narragansett stress patterns (as recorded
    by Williams' diacritics) appear similar to Loup A and Massachusett,
    with a preference for stress on long syllables and alternating rhythm on
    short syllables.
\end{enumerate}

\subsection{Mohegan-Pequot as a Y-Dialect: The Evidence}

Mohegan-Pequot is classified as a Y-dialect based on the reflex of PA *r
$\rightarrow$ /j/. The evidence from the corpus includes:

\begin{itemize}
  \item \textbf{PA *r $\rightarrow$ /j/ in cognate sets:} Where Massachusett
    (N-dialect) has /n/, Mohegan-Pequot has /j/. Stephanie Fielding noted
    this correspondence operating even in modern Connecticut English
    pronunciation (e.g., ``Sullivan'' $\rightarrow$ ``Suh-yuh-van'').

  \item \textbf{The Y-dialect isogloss bundle:} Costa (2007) shows that the
    Y-dialect area (Mohegan-Pequot, Montauk) forms a coherent isogloss
    bundle that separates it from the N-dialect area (Massachusett,
    Narragansett, Nipmuc).

  \item \textbf{Prince-Speck evidence:} The Prince-Speck glossary (1904)
    confirms the Y-dialect classification for early 20th-century
    Mohegan-Pequot, showing /j/ reflexes where N-dialects have /n/.

  \item \textbf{Fielding evidence:} The modern Mohegan dictionary (2013)
    confirms the Y-dialect classification for the revived language.
\end{itemize}

Granberry (2003:13) explicitly documents the core correspondences: Pequot
\textit{p, t, k, ch, s} $\rightarrow$ Mohegan \textit{b, d, g, j, z}

\subsection{Central Algonquian Reference Points}

The dialect classification can be further verified against Central Algonquian
languages, which preserve the PA *r reflex distinctions more clearly than
the SNEA colonial sources:

\begin{itemize}
  \item \textbf{Cree (Y-dialect):} Plains Cree (\url{https://itwewina.altlab.app})
    is a Y-dialect (PA *r $\rightarrow$ /y/), providing a reference for the
    Mohegan-Pequot Y-dialect classification. Cree also preserves the PA
    vowel system more faithfully than any SNEA source.
  \item \textbf{Ojibwe (N-dialect):} The Ojibwe People's Dictionary
    (\url{https://ojibwe.lib.umn.edu}) represents an N-dialect (PA *r
    $\rightarrow$ /n/), providing a reference for the Wampanoag and
    Narragansett N-dialect classification.
  \item \textbf{Munsee Delaware (L-dialect):} The Lenape Talking Dictionary
    (\url{https://www.talk-lenape.org}) represents an L-dialect (PA *r
    $\rightarrow$ /l/), providing a reference for the Mahican transitional
    status between SNEA and Delaware.
\end{itemize}

These Central Algonquian reference points confirm that the SNEA dialect
classifications are consistent with the broader Algonquian family patterns.
The Y/N/L/R classification is not unique to SNEA --- it reflects a
continent-wide pattern of PA *r reflexes across the Algonquian family.
form-initially, and final Pequot \textit{p, t, k, ch, s} $\rightarrow$
Mohegan form-final \textit{b, d, g, j, z}. This lenis-fortis alternation is
a distinctive feature of the Mohegan-Pequot dialect complex.

\subsection{Mahican as Transitional}

Mahican occupies a transitional position between the SNEA languages and the
Delaware languages. Evidence for this transitional status includes:

\begin{itemize}
  \item \textbf{PA *š $\rightarrow$ /s/:} Mahican merges PA *š with /s/, a
    feature shared with Delaware but not with the other SNEA languages
    (which retain /š/ as a separate phoneme).

  \item \textbf{PA *θ $\rightarrow$ /n/:} Mahican merges PA *θ with /n/,
    like the N-dialects of SNEA, but unlike Delaware (which retains /θ/ or
    shifts it to /s/).

  \item \textbf{Word-final vowel deletion:} Mahican shows more advanced
    word-final vowel deletion than the other SNEA languages, approaching the
    Delaware pattern.

  \item \textbf{PA *eː $\rightarrow$ /a/:} Mahican retracts/lowers PA *eː to
    /a/, a development shared with some Delaware dialects but not with the
    other SNEA languages (where PA *eː $\rightarrow$ /iː/).

  \item \textbf{Edwards' data:} Edwards' Mahican forms show features of
    both groups. His \textit{neesoh} 'two' shows the SNEA /n/ reflex of PA
    *r, but the final \textit{-oh} suffix is more characteristic of Delaware.
\end{itemize}

This transitional status means that Mahican cannot be used as a direct model
for the other SNEA languages without accounting for the differences. The
conversion workflow must apply different sound correspondences for Mahican
than for the N-dialect or Y-dialect languages.

\subsection{Cross-Linguistic Cognate Evidence Table}

Table~\ref{tab:cognate-dialect} shows how the same PA root is reflected
across the four dialect types, illustrating the regular sound
correspondences that define the classification.

\begin{table}[h]
\centering
\caption{Cross-dialect cognate evidence for PA *r reflex}
\label{tab:cognate-dialect}
\small
\begin{tabular}{llllll}
\toprule
\textbf{Gloss} & \textbf{PA} & \textbf{N-dialect} & \textbf{Y-dialect} & \textbf{Mahican} & \textbf{L-dialect} \\
 & & \textbf{(Mass.)} & \textbf{(Mohegan)} & \textbf{(Edwards)} & \textbf{(Munsee)} \\
\midrule
'I' & *niːra & neen & neen & neen & niːl \\
'thou' & *kiːra & keen & keen & keen & kiːl \\
'man' & *naːpeːwa & napp & nap & --- & náːpeːw \\
'house' & *wiːkiwaːmi & wetu & wetu & --- & wiːkiwaːm \\
'water' & *nepyi & nippe & nip & --- & mbíː \\
'dog' & *aɬemwa & arum & ayum & --- & --- \\
'three' & *neʔθwi & shwa & shwa & --- & nəxá \\
'four' & *nyeːwi & yaw & yaw & --- & néːwi \\
'five' & *nyaːlanwi & pàpàna & --- & nannah & náːlan \\
\midrule
\textbf{PA *r reflex} & & \textbf{/n/} & \textbf{/j/} & \textbf{/n/} & \textbf{/l/} \\
\bottomrule
\end{tabular}
\end{table}

The table reveals several important patterns:

\begin{enumerate}
  \item \textbf{The PA *r reflex is the primary classifier.} Massachusett
    (N-dialect) and Mohegan (Y-dialect) both show /n/ in 'I' and 'thou',
    but the underlying PA source differs: Massachusett /n/ comes from PA *r,
    while Mohegan /n/ comes from a different source (PA *n). The Y-dialect
    reflex of PA *r is /j/, which appears in other lexical items.

  \item \textbf{Mahican patterns with N-dialects for *r.} Edwards' Mahican
    shows /n/ for PA *r, aligning it with the N-dialect group for this
    isogloss.

  \item \textbf{Munsee (L-dialect) is clearly distinct.} The L-dialect
    reflex /l/ for PA *r separates Munsee from all the SNEA languages.

  \item \textbf{The PA *ɬ reflex distinguishes Narragansett from
    Massachusett.} The 'dog' cognate shows Narragansett /j/ (\textit{ayim})
    vs.\ Massachusett /r/ (\textit{arum}), both from PA *aɬemwa. This is a
    secondary isogloss within the N-dialect group.

  \item \textbf{Lexical stability across dialect boundaries.} Basic
    vocabulary items like 'house', 'water', and 'three' are remarkably
    similar across all four dialect types, confirming the reliability of
    these items for cross-source calibration.
\end{enumerate}

\subsection{Implications for the Conversion Workflow}

The dialect classification has direct implications for the conversion
workflow:

\begin{enumerate}
  \item \textbf{Source-specific sound correspondences.} Each source language
    requires its own set of PA $\rightarrow$ daughter sound correspondences.
    Applying Massachusett correspondences to Mohegan-Pequot data would
    produce systematic errors.

  \item \textbf{Cross-dialect calibration is valid only for shared
    isoglosses.} When using Massachusett to disambiguate Narragansett (both
    N-dialects), the calibration is valid for shared features. When using
    Massachusett to disambiguate Mohegan-Pequot (N-dialect vs.\ Y-dialect),
    the calibration must account for the isogloss differences.

  \item \textbf{Mahican requires separate treatment.} Mahican's transitional
    status means it cannot be grouped with either the N-dialect or Y-dialect
    languages for conversion purposes. A separate set of sound
    correspondences is required.

  \item \textbf{The cognate set approach is the most reliable calibration
    method.} By anchoring each conversion in a known PA reconstruction and
    applying the correct dialect-specific sound correspondences, the
    conversion workflow can produce reliable phonemic transcriptions even
    from highly ambiguous colonial orthographies.
\end{enumerate}

\subsection{Cross-Linguistic Reassessment of Dialect Classification}

The newly acquired reference dictionaries provide independent evidence that
both confirms and complicates the traditional four-way dialect classification:

\begin{itemize}
  \item \textbf{Watkins Cree dialect chain model:} Watkins (1865) documents
    a l/n/th/y dialect chain across Cree, where the same PA segment surfaces
    as /l/, /n/, /θ/, or /j/ depending on dialect. This is a continuum, not a
    discrete classification. The SNEA Y/N/R/L classification may similarly
    represent a continuum rather than discrete groups --- the apparent
    binary between Mohegan (Y-dialect) and Wampanoag (N-dialect) may be an
    artifact of limited documentation rather than a genuine linguistic
    boundary. The Cree model suggests that intermediate dialects (R-dialect
    Loup A, L-dialect Munsee) are not separate types but points on a
    continuum.

  \item \textbf{Brinton Lenape Unami/Minsi split:} Brinton (1888) documents
    a dialect split within Lenape (Unami vs.\ Minsi) that parallels the
    SNEA N-dialect/Y-dialect split. The Unami/Minsi distinction involves the
    same PA *r reflex (Unami /l/ vs.\ Minsi /n/), confirming that the
    Y/N/R/L classification is a recurring pattern across Eastern Algonquian,
    not unique to SNEA.

  \item \textbf{Sauk vowel system as PA model:} The Sauk 4-vowel length
    system (Whittaker 2005) provides a template for the PA vowel system that
    underlies all dialect classifications. The fact that Sauk (Central
    Algonquian) preserves a clean 4-way length distinction suggests that
    the SNEA languages' original vowel systems were similarly structured,
    and that the ambiguity in colonial transcriptions reflects recording
    limitations rather than phonological merger.

  \item \textbf{Arapaho divergence as boundary condition:} The Arapaho
    Dictionary (Cowell et al.\ 2012) represents the extreme end of Algonquian
    phonological divergence. Arapaho's radical changes (PA *p $\rightarrow$
    /b/, PA *k $\rightarrow$ /g/ or /k/, development of /θ/, /x/, triple
    vowels) provide the outer bounds for what is possible in Algonquian sound
    change. The SNEA languages, by contrast, are relatively conservative ---
    their sound changes are minor compared to Arapaho's transformations.
\end{itemize}

These comparative data suggest that the traditional four-way classification
is essentially correct but should be understood as a continuum rather than
discrete categories. The conversion workflows should treat the dialect
classification as a gradient feature, with each source occupying a position
on the continuum rather than belonging to a discrete type.


% ============================================================
% PART III — External Resources and Open Questions
% ============================================================
\part{External Resources and Open Questions}

% Non-SNEA Algonquian Resources — External Resources
% Draftable now. Full details in research cards at .issues/research-cards/

\section{Non-SNEA Algonquian Resources}
\label{sec:external}

This project draws on a range of Algonquian resources beyond the SNEA languages
directly represented in the corpus. These are organized by language family group
below. Full details for each resource are documented in the research cards at
\texttt{.issues/research-cards/}.

\subsection{Proto-Algonquian Reconstructions}

\begin{itemize}
  \item \textbf{Bloomfield (1925, 1946):} Foundational PA reconstruction.
  \item \textbf{Aubin (1975):} \emph{A Proto-Algonquian Dictionary.} Mercury
    Series, Canadian Ethnology Service Paper No. 39. \textasciitilde2,300
    reconstructions. Analyzed Narragansett as substantially accurate relative
    to PA.
  \item \textbf{Hewson (1993):} \emph{A Computer-Generated Dictionary of
    Proto-Algonquian.} 4,066 entries. Online at
    \url{https://protoalgonquian.atlas-ling.ca}. Typos corrected 2015--2017.
  \item \textbf{Pentland (1979):} \emph{Algonquian Historical Phonology.}
    Comprehensive PA-to-daughter phonological developments.
\end{itemize}

\subsection{Eastern Algonquian (Closest Genetic Relatives)}

Eastern Algonquian is the closest genetic relative of the SNEA languages.
These resources provide the best comparative data for resolving SNEA
phonological questions.

\begin{itemize}
  \item \textbf{Mi'kmaq Talking Dictionary} (\url{https://mikmaqonline.org})
    --- 6,000+ words, each recorded by 3+ speakers. Preserves PA cluster
    distinctions collapsed in SNEA colonial sources.
  \item \textbf{Passamaquoddy-Maliseet Dictionary} (Francis \& Leavitt,
    \url{https://pmportal.org}) --- 18,000+ words with video archive. 50+
    years of collaboration among speakers, educators, and linguists.
  \item \textbf{Western Abenaki Dictionary} (Day 1994, 2 vols, 538 pp) ---
    Documents the language as spoken in the last half of the 20th century.
    Abenaki is the closest living relative of the SNEA languages.
  \item \textbf{Nulhegan Abenaki Dictionary} (2025,
    \url{https://abenakionline.com}) --- First publicly shareable English-Abenaki
    dictionary, released with National Archives grant support.
  \item \textbf{Lenape Talking Dictionary} (\url{https://www.talk-lenape.org})
    --- 17,000+ words with audio, stories, usage examples. Lenape is
    transitional between SNEA and Central Algonquian.
  \item \textbf{A Lenâpé-English Dictionary} (Brinton, from Moravian MS,
    \url{https://archive.org}) --- 19th-century dictionary based on Zeisberger,
    Heckewelder, and other Moravian sources.
\end{itemize}

\subsection{Central Algonquian (Most Conservative Phonologically)}

Central Algonquian languages are the most conservative phonologically and
provide the best evidence for PA reflex verification.

\subsubsection*{Cree-Innu Group}

\begin{itemize}
  \item \textbf{Eastern James Bay Cree Dictionary}
    (\url{https://dictionary.eastcree.org}) --- Northern and Southern dialects,
    English-Cree-French.
  \item \textbf{itwêwina Plains Cree Dictionary}
    (\url{https://itwewina.altlab.app}) --- Searchable Plains Cree (Y-dialect).
  \item \textbf{Online Cree Dictionary} (\url{https://www.creedictionary.com})
    --- Aggregates 4 dictionaries: Alberta Elders', Earle Waugh, Wolvengrey,
    Maskwacis.
  \item \textbf{A Dictionary of the Cree Language} (Watkins 1865/1938,
    \url{https://www.creeliteracy.org}) --- Historical dictionary for
    diachronic comparison.
\end{itemize}

\subsubsection*{Ojibwe-Potawatomi Group}

\begin{itemize}
  \item \textbf{Ojibwe People's Dictionary}
    (\url{https://ojibwe.lib.umn.edu}) --- 7,000+ entries with audio.
    Southwestern Ojibwe (N-dialect).
  \item \textbf{Nishnaabemwin Online Dictionary}
    (\url{https://dictionary.nishnaabemwin.atlas-ling.ca}) --- 12,000+ words.
    Odawa dialects of Lake Huron.
  \item \textbf{A Dictionary of the Ojibway Language} (Baraga 1853/1992,
    \url{https://archive.org}) --- Historical dictionary by Bishop Frederic
    Baraga.
  \item \textbf{Potawatomi Language Dictionary}
    (\url{https://www.potawatomidictionary.com}) --- Citizen Potawatomi Nation.
\end{itemize}

\subsubsection*{Sauk-Fox-Kickapoo Group}

\begin{itemize}
  \item \textbf{A Concise Dictionary of the Sauk Language}
    (\url{https://www.sacandfoxnation-nsn.gov}) --- Based on fieldwork
    1990s--2000s.
  \item \textbf{Kickapoo Vocabulary} (Voorhis 1988) --- Algonquian and
    Iroquoian Linguistics, Memoir 6. 205 pp.
\end{itemize}

\subsubsection*{Miami-Illinois and Other Central}

\begin{itemize}
  \item \textbf{Myaamia ILDA Dictionary}
    (\url{https://aacimotaatiiyankwi.org}) --- Miami-Illinois online dictionary
    with mobile app. Myaamia Center, Miami University.
  \item \textbf{The Menomini Language} (Bloomfield 1962,
    \url{https://archive.org}) --- 515 pp. The most phonologically conservative
    Algonquian language. Preserves PA vowel length and *h preaspiration most
    faithfully.
  \item \textbf{English-Shawnee Dictionary} (Emerson) --- Practical dictionary.
\end{itemize}

\subsection{Plains Algonquian (Highly Divergent)}

Plains Algonquian languages are the most divergent from PA and provide the
outer bounds for PA reflex verification.

\begin{itemize}
  \item \textbf{Blackfoot Online Dictionary}
    (\url{https://dictionary.blackfoot.atlas-ling.ca}) --- Part of the
    Algonquian Dictionaries Project.
  \item \textbf{English-Cheyenne Dictionary} (Petter 1915,
    \url{https://cheyennelanguage.org}) --- 1,126 pp with handwritten
    annotations. Cheyenne shows extreme PA consonant shifts.
  \item \textbf{Dictionary of the Arapaho Language} (Salzmann 1983, Cowell
    et al.\ 2012, \url{https://www.colorado.edu/project/arapaho}) --- 243 pp
    (4th ed.). Arapaho shows extreme PA vowel reduction.
\end{itemize}

\subsection{Comparative SNEA Dialectology}

\begin{itemize}
  \item \textbf{Costa (2007):} ``The Dialectology of Southern New England
    Algonquian.'' Definitive modern dialect classification of SNEA.
  \item \textbf{Goddard \& Bragdon (1988):} \emph{Native Writings in
    Massachusett.} Comparative phonology/orthography.
  \item \textbf{Bragdon (2009):} Updated linguistic analysis of Massachusett.
  \item \textbf{Cunningham (2024):} Nanticoke phonological analysis ---
    directly parallel methodology.
  \item \textbf{Masthay (2001):} Mahican language word list.
  \item \textbf{Goddard (1978--2010):} Eastern Algonquian surveys, Munsee
    Delaware phonology, Mahican phoneme inventory.
  \item \textbf{Siebert (1975):} Methods for reconstructing from early
    documentary sources (Powhatan).
\end{itemize}

% Open Problems — External Resources
% Draftable now from research cards' open questions sections

\section{Open Problems}
\label{sec:open-questions}

The following problems remain unresolved across the SNEA sources. These are
drawn from the ``Open Questions for Linguist Review'' sections of each
research card.

\subsection{Phonological}

\begin{enumerate}
  \item \textbf{Preaspiration inference rules:} Can a reliable set of
    morphotactic rules predict when preaspiration /h/ was present but
    unwritten? Or does this require per-entry comparative verification?
  \item \textbf{The \textit{gh} grapheme:} Is Edwards' \textit{gh} /x/, /h/,
    or /k/? Cross-linguistic cognate evidence from Munsee Delaware is needed.
  \item \textbf{Nasalization disambiguation:} When does \textit{an}/\textit{aun}
    represent /ã/ vs.\ /an/? Massachusett cognate comparison is the primary
    method, but it is not always available.
\end{enumerate}

\subsection{Methodological}

\begin{enumerate}
  \item \textbf{Pidgin vs.\ genuine dialect:} Are the simplifications in
    Winslow's data evidence of a pre-existing Massachusett Pidgin (Goddard
    1977) or simply a learner's imperfect recording?
  \item \textbf{Dialect stratification:} How much of the variation between
    Natick (Eliot) and Pokanoket (Winslow/Wood) is orthographic vs.\
    phonological vs.\ dialectal?
  \item \textbf{Fielding circularity:} Which Fielding (2013) entries derive
    from Prince-Speck vs.\ independent sources? A source-tracing audit is
    needed.
\end{enumerate}

\subsection{Corpus-Level}

\begin{enumerate}
  \item \textbf{Training data shortage:} Only 3 of 2,134 Narragansett records
    have phonetic fields. Hand-annotating \textasciitilde50 entries is
    recommended before ML approaches are viable.
  \item \textbf{Tiny corpus reliability:} Winslow's 24 records --- can a
    dataset this small yield reliable phonological data?
\end{enumerate}

% Aggregated Linguist Feedback — External Resources
% Template now, content later

\section{Aggregated Linguist Feedback}
\label{sec:feedback}

\textit{This section will be populated as linguist review cycles are completed
for each source workflow. The template below shows the structure.}

\subsection{Feedback Categories}

Each source workflow's feedback section records:

\begin{itemize}
  \item \textbf{Conversion accuracy:} Per-entry verification against linguist
    judgment.
  \item \textbf{Problem identification:} Issues the linguist flags that the
    research phase missed.
  \item \textbf{Methodology critique:} Linguist assessment of the FSM/pipeline
    approach.
  \item \textbf{Revitalization applicability:} Whether the output is usable
    for practical language teaching.
  \item \textbf{Cross-cutting patterns:} Feedback that applies across multiple
    sources.
\end{itemize}

\subsection{Cross-Cutting Patterns}

\texttt{[To be populated after review cycles]}

% Future Work — External Resources
% Draftable now

\section{Future Work}
\label{sec:future}

The combined research across all eight sources suggests several directions
for future work that may feed back into the dependent spec requirements.

\subsection{Immediate Priorities}

\begin{enumerate}
  \item \textbf{Hand-annotated training data:} Create a gold-standard set of
    \textasciitilde50 Narragansett entries with verified phoneme mappings,
    to enable ML classifier approaches.
  \item \textbf{Preaspiration inference model:} Develop morphotactic rules
    for predicting unwritten /h/ across all colonial sources, using Edwards
    and Prince-Speck as training data.
  \item \textbf{Fielding circularity audit:} Trace which Fielding (2013)
    entries derive from Prince-Speck vs.\ independent sources.
\end{enumerate}

\subsection{Longer-Term Directions}

\begin{enumerate}
  \item \textbf{Unified conversion pipeline:} Combine all eight source
    workflows into a single FSM with source-detection frontend.
  \item \textbf{Cross-linguistic calibration database:} Build a database of
    confirmed cognate sets across all SNEA languages for automated
    disambiguation.
  \item \textbf{Revitalization-ready output:} Generate pronunciation guides
    and teaching materials directly from the conversion pipeline.
\end{enumerate}


\end{document}
